\documentclass[8pt]{article}
\linespread{0.1} 
\usepackage{extsizes}
\usepackage{amsmath}
\usepackage{amssymb}
\usepackage{amsthm}
\usepackage{ mathrsfs }
\usepackage{ mathtools}
\usepackage{enumerate}
\usepackage{bbm}
\usepackage{lipsum}
\usepackage{fancyhdr}
\usepackage{tikz-cd} 
\usepackage{adjustbox} 
\usepackage{ stmaryrd }
\numberwithin{equation}{section}
\usetikzlibrary{arrows}
\newcommand{\midarrow}{\tikz \draw[-triangle 90] (0,0) -- +(.1,0);}
\tikzset{commutative diagrams/.cd,
mysymbol/.style={start anchor=center,end anchor=center,draw=none}
}
\newcommand\comsymb[2][\circlearrowleft]{%
  \arrow[mysymbol]{#2}[description]{#1}}

\newtheorem{theorem}{Theorem} [section] 
\newtheorem{proposition}{Proposition}[section] 
\newtheorem{definition}{Definition}[section] 
\newtheorem{lemma}{Lemma}[section] 
\newtheorem{notation}{Notation}[section] 
\newtheorem{remark}{Remark}[section] 
\newtheorem{corollary}{Corollary} [section] 
\newtheorem{terminology}{Terminology}[section] 
\newtheorem{fact}{Fact}[section] 
\newtheorem{example}{Example}[section] 
\newtheorem{claim}{Claim}
\newenvironment{claimproof}[1]{\par\noindent\underline{Proof:}\space#1}{\hfill $\blacksquare$}

\DeclareMathOperator{\tmod}{mod}
\DeclareMathOperator{\triv}{triv}
\DeclareMathOperator{\ind}{ind}
\DeclareMathOperator{\straw}{s.t.}
\DeclareMathOperator{\ev}{ev}
\DeclareMathOperator{\pr}{pr}
\DeclareMathOperator{\Aut}{Aut}
\DeclareMathOperator{\Map}{Map}
\DeclareMathOperator{\Stab}{Stab}
\DeclareMathOperator{\Ind}{Ind}
\DeclareMathOperator{\GL}{GL}
\DeclareMathOperator{\SL}{SL}
\DeclareMathOperator{\SO}{SO}
\DeclareMathOperator{\Sp}{Sp}
\DeclareMathOperator{\esssup}{esssup}
\DeclareMathOperator{\abs}{abs}
\DeclareMathOperator{\rel}{rel}
\DeclareMathOperator{\diam}{diam}
\DeclareMathOperator{\rank}{rank}
\DeclareMathOperator{\Hom}{Hom}
\DeclareMathOperator{\Homk}{Hom_{k-alg}}
\DeclareMathOperator{\Ker}{Ker}
\DeclareMathOperator{\Image}{Im}
\DeclareMathOperator{\Dom}{Dom}
\DeclareMathOperator{\grad}{grad}
\DeclareMathOperator{\divergence}{div}
\DeclareMathOperator{\rk}{rk}
\DeclareMathOperator{\Span}{Span}
\DeclareMathOperator{\mSpec}{mSpec}
\DeclareMathOperator{\MaxSpec}{MaxSpec}
\DeclareMathOperator{\interior}{int}
\DeclareMathOperator{\supp}{supp}
\DeclareMathOperator{\id}{id}
\DeclareMathOperator{\sgn}{sgn}
\DeclareMathOperator{\Mat}{Mat}
\DeclareMathOperator{\PD}{PD}
\DeclareMathOperator{\ext}{ext}
\DeclareMathOperator{\vol}{vol}
\DeclareMathOperator{\inte}{int}
\DeclareMathOperator{\diverg}{div}
\DeclareMathOperator{\curl}{curl}
\DeclareMathOperator{\image}{im}
\DeclareMathOperator{\dR}{dR}
\DeclareMathOperator{\Ob}{Ob}
\DeclareMathOperator{\Mnfd}{Mnfd}
\DeclareMathOperator{\OpEmb}{OpEmb}
\DeclareMathOperator{\Spec}{Spec}
\DeclareMathOperator{\Spa}{Spa}
\DeclareMathOperator{\Lim}{Lim}
\DeclareMathOperator{\Alg}{\mathcal{A}}
\DeclareMathOperator{\glim}{glim}
\DeclareMathOperator{\con}{con}
\DeclareMathOperator{\Spv}{Spv}
\DeclareMathOperator{\clet}{c}
\DeclareMathOperator{\tat}{at}
\DeclareMathOperator{\Sheaf}{Sheaf}
\DeclareMathOperator{\Cont}{Cont}
%\newcommand*{\name}[\num_arguments][default values]{{\color{#1}\Large #2}}
\newcommand{\defeq}{\vcentcolon=}
\newcommand{\Ouv}{\mathscr{O}}
\newcommand{\norm}[1]{\Vert #1 \Vert}
\newcommand{\opNorm}[2]{\norm{#1}_{#2\to#2}}
\newcommand{\normL}[3]{\norm{#1}_{L^{#2}(#3)}}
\newcommand{\N}[0]{\mathbb{N}}
\newcommand{\m}[0]{\mathfrak{m}}
\newcommand{\R}[0]{\mathbb{R}}
\newcommand{\Z}[0]{\mathbb{Z}}
\newcommand{\C}[0]{\mathbb{C}}
\newcommand{\Q}[0]{\mathbb{Q}}
\newcommand{\F}[0]{\mathcal{F}}
\newcommand{\U}[0]{\mathcal{U}}
\newcommand{\B}[0]{\mathcal{B}}
\newcommand{\G}[0]{\mathbb{G}}
\newcommand{\A}[0]{\mathbb{A}}
\newcommand{\T}[0]{\mathbb{T}}
\newcommand{\K}[0]{\mathcal{K}}
\newcommand{\Para}[0]{\mathbb{P}}
\newcommand{\M}[0]{\mathcal{M}}
\newcommand{\maniN}[0]{\mathcal{N}}
\newcommand{\cinf}[0]{\mathcal{C}^\infty}
\newcommand{\torus}[0]{\mathbb{T}}
\newcommand{\cR}[0]{\mathcal{R}}
\newcommand{\sep}[0]{\:|\:}
\newcommand{\pd}[2]{{\frac {\partial #1} {\partial #2}}}
\newcommand{\compinc}[0]{\subset \subset}
\newcommand{\sH}[0]{\mathscr{H}}
\newcommand{\Nfil}[0]{\mathcal{N}}
\newcommand{\g}[0]{\mathfrak{g}}
\newcommand{\p}[0]{\mathfrak{p}}
\newcommand{\QC}[0]{\mathfrak{Q}_C}
\newcommand{\ab}[1]{\vert #1 \vert}
\newcommand{\st}[0]{\:\straw\:}
\newcommand{\cGamma}[0]{c\Gamma}
\newcommand{\sC}[0]{\mathscr{C}}
\newcommand{\sU}[0]{\mathscr{U}}
\newcommand{\sV}[0]{\mathscr{V}}
\newcommand{\sA}[0]{\mathscr{A}}
\newcommand{\sB}[0]{\mathscr{B}}
\newcommand{\sX}[0]{\mathscr{X}}
\newcommand{\sS}[0]{\mathscr{S}}
\newcommand{\cB}[0]{\mathcal{B}}
\newcommand{\cA}[0]{\mathcal{A}}
\newcommand{\cC}[0]{\mathcal{C}}
\newcommand{\nat}[0]{\mathrm{nat}}
\newcommand{\tand}[0]{\text{ and }}
\newcommand{\tor}[0]{\text{ or }}
\newcommand{\dcirc}[0]{{\circ\circ}}
\newcommand{\otherwise}[0]{\text{ otherwise }}
\newcommand{\Rat}[0]{\mathfrak{Rat}}
\newcommand{\polring}[1]{[X_1,\cdots,X_{#1}]}
\newcommand{\compring}[1]{\langle X_1,\cdots,X_{#1}\rangle}
\newcommand{\fib}[1]{%
  \mathbin{\mathop{\times}\limits_{#1}}%
}
\newcommand{\tens}[1]{%
  \mathbin{\mathop{\otimes}\displaylimits_{#1}}%
}

\title{Rigid Analytic Geometry}
\author{Based on Lectures given by\\\\Prof. Dr. Jens Franke}
\date{SoSe 26, University of Bonn}


\begin{document}
\maketitle

%4/13

\section{Introduction - Spectral Spaces and Spaces of Valuation Spectra}
\subsection{Introduction}

\begin{definition}
    Let $\Omega\subseteq\C$ be a domain. We define a sheaf of holomorphic functions as,
\begin{equation*}
    \Omega\supseteq V\mapsto \Ouv(V)=\{f:V\to\C\sep \text{$f$ is holomorphic}\}.
\end{equation*}
\end{definition}

One can equip $\Ouv(V)$ this with the semi-norm such that,
\begin{equation*}
    \norm{f}_{V,A} = \max\{\vert f(z)\vert\sep z\in A\},
\end{equation*}
where $A\subseteq\Omega$ is a closed and bounded in $\C$, namely a compact set $A$. With respect to this semi-norm, we define basis elements of the topology in the following way.
% \par A subset of $\Ouv(V)$ is a neighbourhood of $0$ if and only if there is a compact $A$ as above and a positive number $\varepsilon$ such that it contains 
% \begin{equation*}
%     \{f\in\Ouv(V)\sep \norm{f}_A<\varepsilon\}.
% \end{equation*}

\begin{proposition}
    The semi-norm induces the topology on $\Ouv(V)$ whose basis elements of the neighbourhood of $0$ are of the following form. For $f\in\Ouv(V)$ and $\varepsilon>0$,
    \begin{equation*}
    \{f\in\Ouv(V)\sep\norm{f}_{V,A}<\varepsilon\}.
    \end{equation*}
\end{proposition}
\begin{definition}
    A ring is said to be Bézout if every finitely generated ideal is principal.
\end{definition}
\begin{remark}
    By Weierstrass product theorem, $\Ouv(V)$ is not Noetherian but it is Bézout.
\end{remark}
\begin{proposition}
    The stalks of $\Ouv(V)$ which is a ring of germs of holomorphic functions are Noetherian. This holds for holomorphic functions in several complex variables as well.
\end{proposition}
\par The goal of the course is to generalize these results to $p$-adic cases.
\par We will allow ground field $K$ with an absolute value $\vert\cdot\vert:K\to[0,\infty)$ satisfying,
\begin{equation*}
    \vert0\vert=0,\vert 1\vert = 1, \vert ab\vert = \vert a\vert\vert b\vert, \vert a+b\vert\leq\max\{\vert a\vert,\vert b\vert\}.
\end{equation*}
% If the forth equation holds, we call $\vert\cdot\vert$ an ultra-metric. $K$ must be a complete metric space with the distance $d(x,y) = \vert x-y\vert$ and the topology will non-discrete. It turns out that the $\vert\cdot\vert_K$ has a unique extension to $\vert\cdot\vert_L$ for every algebraic extension $L/K$ that $L$ is complete with respect to $\vert\cdot\vert_L$ where $[L:K]<\infty$ and the completion of the algebraic closure of $K$ (e.g. $\C_p$ where $\Q_p$) is algebraically closed.
\begin{definition}
\label{definition_1_1}
Let $K$ be a field which is complete with respect to the map $\vert\cdot\vert_K:K\to [0,\infty)$ is an absolute value of $K$ if we have the following conditions.
\begin{enumerate}[i).]
    \item $\ab{1}_K=1$,
    \item $\ab{0}_K=0$,
    \item $\forall a,b\in K, \ab{ab}_K = \ab{a}_K\ab{b}_K$.
\end{enumerate}
Furthermore, it is called ultra-metric if we have,
\begin{equation*}
    \forall a,b, \ab{a+b}_K\leq\max\{\ab{a}_K,\ab{b}_K\}.
\end{equation*}
\end{definition}

\begin{proposition}
    Suppose we have an ultra-metric $\abs{\cdot}_K$ on a field $K$. If $K$ is complete with respect to the topology induced by the metric $d(x,y) = \ab{x-y}_K$, then for a finite extension $L/K$, there is an extension $\ab{\cdot}_L:L\to[0,\infty)$ of $\ab{\cdot}_K$ such that $L$ is complete with respect to the canonical metric obtained from $\ab{\cdot}_L$.
\end{proposition}
\begin{proposition}
    \label{proposition_completion_algebraic_closure}
    Let $K$ be a field with an ultrametric $\ab{\cdot}_K$. The completion $\widehat{\overline{K}}$ of the algebraic closure $\overline{K}$ of $K$ with respect to $\ab{\cdot}_K$ is again algebraically closed.
\end{proposition}
\begin{example}
    In Proposition \ref{proposition_completion_algebraic_closure}, we can take for example $\widehat{\overline{K}}$ to be $\C_p$ where $K=\Q_p$. 
\end{example}
We will now give the definition of our basic objects.
\begin{definition}
\label{definition_1_2}
Let $(K,\ab{\cdot}_K)$ be a field which is complete with respect to the absolute value $\ab{\cdot}_K$. We define the Tate algebra over $K$ as
\begin{equation*}
\mathbb{T}_m(K)\defeq\left\{f = \sum_{\alpha\in\N^m}f_\alpha X^\alpha\:\Bigg{|}\:\substack{\text{for all $\varepsilon>0$, there are only}\\ \text{finitely many $\alpha$ such that $\vert f_\alpha\vert>\varepsilon$}}\right\}
\end{equation*}
We define,
\begin{equation*}
    \norm{f|_{\mathbb{T}_m}} \defeq \max\{\vert f_\alpha\vert_K\sep \alpha\in\N^m\}.
\end{equation*}
% where $K$ is algebraically closed, enough to consider
% \begin{equation*}
%     \norm{f|_{\mathbb{T}_m}}=\max\{\vert f(z)\vert\sep z=(z_i)_{i=1,\cdots,m}\in K^m ,\vert z\vert\leq 1\}.
% \end{equation*}
% Note that such maximum is attained.
\end{definition}
\begin{proposition}
Suppose $K$ is algebraically closed in Definition \ref{definition_1_2}. Then we have,
\begin{equation*}
    \norm{f|_{\mathbb{T}_m}}=\max\{\vert f(z)\vert\sep z=(z_i)_{i=1,\cdots,m}\in K^m ,\vert z\vert\leq 1\},
\end{equation*}
and such maximal is actually attained.
\end{proposition}
\begin{remark}
    In $p$-adic settings, $K^m$ is totally disconneted so assertion like the "identity theorem" for holomorphic functions becomes problematic.
\end{remark}
In order to resolve the above obstacle Grothendieck introduced Étale topology.% on the proof of all germs?.
\subsection{Main Theorems of the Course}
\begin{definition}
    An element $a$ is power-bounded with respect to a norm $\norm{\cdot}$ if there is $M>0$ such that for all $n\in\N$, $\norm{a}^n\leq M$.
\end{definition}
\begin{definition}
    Let $A$ be a topological ring. An element $a\in A$ is topologically nilpotent if 
    \begin{equation*}
        \lim_{n\to \infty}a^n = 0.
    \end{equation*}
\end{definition}
\begin{definition}
    We define the following subset of a Tate algebra $\mathbb{T}_m$.
    \begin{align*}
        \mathbb{T}_m^\circ&\defeq\{f\in\mathbb{T}_m\sep \text{$f$ is power bounded}\},\\
        \mathbb{T}_m^{\circ\circ} &\defeq \{f\in\mathbb{T}_m\sep \text{$f$ is topologicaly nilpotent $\lim_{k\to \infty}f^k=0$}\}.
    \end{align*}
\end{definition}
\begin{proposition}
    $\norm{\cdot|_{\mathbb{T}_m}}$ is multiplicative. (i.e. $\norm{fg}=\norm{f}\norm{g}$ and $\norm{1}=1$). In particular,
    \begin{equation*}
        \mathbb{T}_m^\circ = \{f\in\mathbb{T}_m(K)\sep \norm{f|_{\mathbb{T}_m}} \leq 1\}.
    \end{equation*}
   and
    \begin{equation*}
        \mathbb{T}_m^{\circ\circ} = \{f\in\mathbb{T}_m(K)\sep \norm{f|_{\mathbb{T}_m}}< 1\}.
    \end{equation*}
    From the definition of ultrametric, it is clear that $\mathbb{T}^\circ_m$ is a subring and $\mathbb{T}_m^{\circ\circ}$ is an ideal. 
\end{proposition}
\begin{proposition}
    $\mathbb{T}_m(K)$ is Noetherian and all ideals are closed in it. If $\m$ is a maximal ideal, its residue field $k(\m)=\mathbb{T}_m/\m$ is a finite field extension of $K$ (thus $\vert\cdot\vert_K$ extends uniquely to $k(\m)$).
\end{proposition}
\begin{proposition}
    There is a topologically nilpotent unit, thus there are no proper open ideals.
\end{proposition}
\begin{proposition}
$\mathbb{T}_m(K)$ is regular of Krull dimension $m$.
\end{proposition}
\begin{definition}
    A ring $A$ is Jacobson if every prime ideal is an intersection of some maximal ideals.
\end{definition}
\begin{proposition}
    $\mathbb{T}_m(K)$ is a Jacobson ring. (The closed points of $\Spec\mathbb{T}_m$ are dense in every closed subsets)
\end{proposition}
\begin{notation}
For a ring $A$, we denote
\begin{equation*}
    \Sp A \defeq \{\m\in\Spec A\sep\text{$\m$ is a maximal ideal}\}.
\end{equation*}
\end{notation}
\begin{notation}
Let $f\in\T_m(K)$ and $x\in\Sp\T_m(K)$. We denote,
\begin{equation*}
    f(x) \defeq f\tmod{x}\in k(x).
\end{equation*}
\end{notation}
\begin{proposition}
\begin{equation*}
    \norm{f|_{\mathbb{T}_m}}=\max\{\vert f(x)\vert \sep x\in\Sp \mathbb{T}_m\}.
\end{equation*}
\end{proposition}
% We use the set of rational open subsets, i.e. of ideals of the form 
% \begin{equation*}
%     \mathcal{R}_{\Sp \mathbb{T}_m}(f_0|f_1,\cdots,f_m) = \{x\in\Sp\mathbb{T}_m\sep \vert f_0(x)\vert\geq\vert f_i(x)\vert, 1\leq i\leq m\},
% \end{equation*}
% where $\langle f_0,\cdots,f_m\rangle_{\mathbb{T}_m}$ will be $\mathbb{T}_m$. 
We define a basis element of the topology in $\T_m(K)$ as follows.
\begin{definition}[Rational Open Sets]
Suppose $\T_m(K)=\langle f_0,\cdots,f_m\rangle_{\mathbb{T}_m}$. The rational open set associated to the generators $\langle f_0,\cdots,f_m\rangle_{\mathbb{T}_m}$ is 
\begin{equation*}
    \mathcal{R}_{\Sp \mathbb{T}_m}(f_0|f_1,\cdots,f_m) = \{x\in\Sp\mathbb{T}_m\sep \vert f_0(x)\vert\geq\vert f_i(x)\vert, 1\leq i\leq m\}.
\end{equation*}
\end{definition}
\begin{remark}
The Grothendieck topology enforces the quasi-compactness of these subsets.
\end{remark} 
Additional points were :? introduced by van der Put in the late 1970s using a filter machinary equivalent to topos? points. He also ? equivalence with ? valuation ? ? not 
\par Berkovich in the 1980s constructed a multiplicative seminorm $\mathbb{T}_m\to[0,\infty)$ which are similar to $f\mapsto\vert f(x)\vert $ for $x\in\mathbb{T}_m$ or $f\mapsto\norm{f|_{\mathbb{T}_m}}$. One then gets a compact Hausdorff space.
\par Huber constructed more general valuations and one gets a spectral space with a sheaf category equivalent to sheaves on $\Sp\mathbb{T}_m$ (unlike Berkovich).
\par In the Berkovich machinery, one also has a point corresponding to 
\begin{equation*}
    \norm{f}_r=\max\{r^{\vert \alpha\vert}\vert f_\alpha\vert_K\sep \alpha\in\N^m\},
\end{equation*}
where $\vert\alpha\vert = \sum\alpha_i$ and $0\leq r\leq 1$.
\begin{remark}
    \begin{equation*}
        \norm{f}_1=\norm{f|_{\mathbb{T}_m}},\norm{f}_0=0.
    \end{equation*}
\end{remark}
\begin{definition}
    Let $0\leq r\leq 1$. We define, for $f\in\T_m(K)\backslash\{0\}$,
\begin{equation*}
    N_r(f) \defeq (\norm{f}_r,n)
    % N_r(f) \begin{cases}
    %     (\norm{f}_r,n),\quad\text{where }f\not=0, \text{$n$ maximal with $r^n\vert f_n\vert_K=\norm{f}_r$},\\
    %     0\quad f=0.
    % \end{cases}
\end{equation*}
where $n=\max\{i\in\N\sep \norm{f}_r=r^i\ab{f_i}_K\}$.
\end{definition}
\begin{example}
    Consider the case $\T_1(K)$ and $f\in\T_1(K)\backslash\{0\}$ where
    \begin{equation*}
        f = \sum_{n=0}^\infty f_nT^n.
    \end{equation*}
%     \begin{equation*}
%     m=1,f=\sum_{n=0}^\infty f_nT^n \in\mathbb{T}_1.
% \end{equation*}
Then $N_r$ defines a valuation, $\mathbb{T}_1\to\Gamma\cup\{0\}$ where $\Gamma = (0,\infty)_\R\times\Z$ ordered lexicographically with the first factor being the most significant and $0<\gamma$ for all $\gamma\in \Gamma$.
\end{example}
One has 
\begin{equation}
    \label{formula_1_1}
    N(a,b) = N(a)N(b),N(1) = 1,N(0) = 0, N(a+b)\leq\max\{N(a),N(b)\}.
\end{equation}
When $0<r<1$ and $f\in \mathbb{T}_1\backslash\{0\}$, there is $R\in(r,1]_\R$ such that 
\begin{equation}
    \label{formula_1_2}
    f(z) =\norm{f_r}\left({\frac {\vert z\vert} r}\right)^{n'}, r<\vert z\vert\leq R.
\end{equation}
where $N$ is as in Equation \eqref{formula_1_1}, and $\norm{f}_r = \sup\{\vert f(z)\vert\sep \vert z\vert\leq r\}$ where $r=1$, Equation \eqref{formula_1_2} has no analog as no ? $R$ extends as $f(z)$ is undefined where $\vert z\vert>1$. This means that $N_1$ is extended 
\begin{definition}
    \begin{equation*}
    \Spa(\mathbb{T}_m,\mathbb{T}_m^0) = \{\text{continuous valuations $\gamma$ of $\mathbb{T}_m\:\st\:\forall f\in\T_m^\circ,\gamma(f)\leq 1$}\}.
\end{equation*} 
\end{definition}
\begin{remark}
    $N_1(T) = (1,1), N_1(T^{2026}) = (1,2026)$.
\end{remark}
%4/17
\subsection{Spectral Spaces}
\begin{remark}
    $\emptyset$ is not irreducible.
\end{remark}
\begin{proposition}
    Let $X$ be a topological space. For a point $x\in X$, 
    \begin{equation*}
        \overline{\{x\}} = \{ y \in X\sep \text{every neighbourhood of $y$ in $X$ cotains $x$}\}.
    \end{equation*}
\end{proposition}
\begin{definition}
    A point $x\in X$ is a generic point if $\overline{\{x\}} = X$, (i.e. $\{x\}$ is dense in $X$).
\end{definition}
\begin{notation}
    For a topological space $X$, we denote the collection of all open sets to be $\Ouv_X$.
\end{notation}
\begin{definition}
    \label{definition_1}
    Let $X$ be a topological space and $x,y\in X$. We say $x,y$ are topologically distinguishable if there is $U\in\Ouv_X$ such that either $x\in U$ and $y\not\in U$ or $x\not\in U$ and $y\in U$. They are otherwise indistinguishable. 
\end{definition}
\begin{definition}
    A topological space $X$ is $T_0$ or Komogolov if for any pair of points are distinguishable.
\end{definition}
\begin{definition}
    \label{definition_2}
    A topological space $X$ is sober if it satisfies one of the following equivalent conditions.
    \begin{enumerate}[1).]
        \item Every irreducible neighborhood of $X$ has precisely one generic point.
        \item $X$ is $T_0$ and every irreducible subset has at least one generic point.
    \end{enumerate}
\end{definition}
\begin{remark}
    It is easy to see that these are local conditions.
\end{remark}
\begin{definition}
    A topological space is $T_1$ if it satisfies one of the following equivalent conditions.
    \begin{enumerate}[1).]
        \item For any $x,y\in X$, $x\not=y\Rightarrow $ there is $U\in\Ouv_X$ such that $x\in U$ and $y\not\int U$.
        \item Every point of $X$ is closed.
    \end{enumerate}
    It is $T_2$ or Hausdorff if it satisfies one of the following equivalent conditions.
    \begin{enumerate}[1).]
        \item Any pair of points $x,y\in X$ has disjoint neighborhoods $U,V$ such that $x\in U$ and $y\in V$.
        \item The diagonal set $\{(x,x)\sep x\in X\}$ is closed in $X\times X$.
    \end{enumerate}
\end{definition}
\begin{definition}
    $X$ is quasi-compact if every open covering admits a fintie subcovering. A subset $M$ of $X$ is quasi-compact if it is a quasi-compact as a topological space with the subset topology from $X$.
\end{definition}
\begin{proposition}
    If $X$ can be covered by finitely many quasi-compact open neighbourhoods then it is quasi-compact.
\end{proposition}
\begin{proposition}
    Every closed subset of a quasi-compact space is again quasi-compact.
\end{proposition}
\begin{proposition}
    \label{fact_3}
    Let $X\stackrel{f}{\to}Y$ be a continuous with $X$ being quasi-compact and $Y$ being Hausdorff, then $f$ is closed (i.e. every closed set $A$ of $X$ has a closed image $f(A)$ in $Y$).
\end{proposition}
\begin{definition}
    A subset $\mathcal{B}\subseteq\Ouv_X$ is a topology basis for $X$ if for any $U\in\Ouv_X$, we have $U=\bigcup_{V\in\mathcal{B}_U}V$ where $\mathcal{B}_U=\{V\in\mathcal{B}\sep V\subseteq U\}$.
\end{definition}
\begin{remark}
    The subsets of the form 
    \begin{equation*}
        D(f) = \{\mathfrak{p}\in\Spec A\sep f\not\in \mathfrak{p}\}
    \end{equation*}
    form a topology basis for $\Spec A$.
\end{remark}
\begin{proposition}
    \label{proposition_2_5}
    \label{fact_5}
    For a topological space $X$, the following conditions are equivalent.
    \begin{enumerate}[i).]
        \item $\QC(X) \defeq \{\Omega\in\Ouv_X\sep \Omega\text{ is quasi-compact}\}$ is a topology basis closed under finite intersection (including the empty intersection $\bigcap_\emptyset = X$).
        \item There is a topology base $\mathcal{B}\subseteq\QC(X)$ such that $\mathcal{B}$ consists of quasi-compact and it is closed under the finite intersection in $X$.
    \end{enumerate}
\end{proposition}
\begin{definition}
    A topological space $X$ is spectral if it is sober and satisfies one of the conditions in Proposition \ref{proposition_2_5}. 
\end{definition}
\begin{definition}
    A spectral map is a map $X\stackrel{f}{\to}Y$ between spectral spaces such that for any quasi-compact subset $\Omega$ in $Y$, $f^{-1}(\Omega)$ is quasi-compact subset of $X$.
\end{definition}
\begin{remark}
    A spectral map is automatically continuous.
\end{remark}
\begin{proposition}
    \label{facr1.2.6}
    If $\mathcal{B}\subseteq \Ouv_Y$ is a topology base as in the second condition of Proposition \ref{proposition_2_5} for $Y$, then $f$ is spectral if and only if $\forall \Omega\in\mathcal{B},f^{-1}(\Omega)\in\QC(X)$.
\end{proposition}
\begin{example}
    $\Spec A$ is spectral and every ring homomorphism defines a spectral map and $\{D(f)\sep f\in A\}$ are such base as in Proposition \ref{proposition_2_5} for $\Spec A$.
\end{example}
\begin{example}
    The underlying space of qcqs (i.e. quasi-compact and quasi-separated) scheme is spectral.
\end{example}
% \begin{definition}
%     \begin{equation*}
%         \underline{\Ouv_l}(A) = \{\text{$I$ ideals of $A$} \st \sqrt{I}=I\}.
%     \end{equation*}
% \end{definition}
\begin{remark}
    \label{remark_example_1_b}
    Every closed and every quasi-compact open sets of a spectral space is spectral.
\end{remark}
\begin{remark}
    A Noetherian space $X$ (i.e. $\QC(X)=\Ouv_X$) is spectral if and only if it is sober.
\end{remark}
\begin{proposition}
    Every non-empty closed subset of a spectral space $X$ contains a closed point.
\end{proposition}

\begin{proof}
    By Remark \ref{remark_example_1_b}, it suffices to show that every spectral space $X\not=\emptyset$ contains a closed point as non-empty closed subsets are again spectral. Set,
    \begin{equation*}
        \mathcal{M}\defeq\{Z\subseteq X\sep \text{$Z$ is closed}\}.
    \end{equation*}
    Equiq $\mathcal{M}$ with a partial order by $\subseteq$. Since $X$ is quasi-compact, $\mathcal{M}$ has a minimal element with such order. Let $Z$ be a such minimal element in $\mathcal{M}$. If $Z$ has two different element $x,y$, then by $T_0$ condition, we have $U\in\Ouv_X$ such that $x\not\in U$ and $y\in U$. Then $Z'=Z\backslash U$ is in $\mathcal{M}$. As $x\in Z'$, this contradicts to the minimality of $Z$ as $y\in Z\backslash Z'$.
\end{proof}
\begin{definition}
    A filter base of a set $M$ is a non-empty subset $\mathcal{B}$ of $\mathfrak{P}(M)$ such that 
    \begin{enumerate}[i).]
        \item $\emptyset\not\in\mathcal{B}$, 
        \item for $N,O\in\mathcal{B}$ there is $P\in\mathcal{B}$ with $P\subseteq N\cap O$.
    \end{enumerate}
    A filter is a filter base $\mathcal{F}$ such that for $N\in\mathcal{F}$, and $N\subseteq O\subseteq M$ implies $O\in\mathcal{F}$.
    \par A filter generated by a filter base $\mathcal{B}$ is such that 
    \begin{equation*}
        \mathcal{F} \defeq \{N\sep N\subseteq M\text{ and there is $O\in\mathcal{B}$ such that $O\subseteq N$}\}.
    \end{equation*}
    A direct image under a map $M\stackrel{f}{\to}X$ of a filter $\mathcal{F}$ is such that 
    \begin{equation*}
        f_*\mathcal{F}\defeq\{Y\sep Y\subseteq X \text{ and } f^{-1}(Y)\in\mathcal{F}\}.
    \end{equation*}
\end{definition}
\begin{example}
    \label{example_2}
A neighbourhood filter $\mathcal{N}_x$ of $x\in X$ is such that 
\begin{equation*}
    \mathcal{N}_x\defeq\{N\subseteq X\sep \exists U\in\Ouv_X\:\st\: x\in U\subseteq N\}.
\end{equation*}
\end{example}
\begin{example}
    \label{example_3}
    A principal filter $\mathcal{P}_x$ of $x\in X$ is such that 
    \begin{equation*}
        \mathcal{P}_x=\{N\subseteq X\sep x\in N\}.
    \end{equation*}
\end{example}
\begin{remark}
    A neighbourhood base for $x\in X$ is a filter base $\mathcal{B}$ such that it generates the neighbourhood filter of $x$.
\end{remark}
\begin{remark}
    The notion of filters is introduced by Cartan.
\end{remark}
\begin{proposition}
    \label{proposition_2_8}
    \label{fact_8}
    For a filter $\mathcal{F}$ of $X$, the following statements are equivalent.
    \begin{enumerate}[1).]
        \item $\mathcal{F}$ is a $\subseteq$-maximal element of the set of filters on $X$.
        \item If $M,N\subseteq X$ and $M\cup N\in\mathcal{F}$ then $M\in\mathcal{F}$ or $N\in\mathcal{F}$.
        \item For $M\subseteq X$, precisely one of $M$ and $X\backslash M$ is in $\mathcal{F}$.
    \end{enumerate}
\end{proposition}
\begin{definition}
    \label{definition_8}
    Such filter $\mathcal{F}$ satisfying the conditions in Proposition \ref{proposition_2_8} is called an ultrafilter.
\end{definition}
\begin{theorem}[Tarski]
    \label{theorem_tarski}
    Every filter on $X$ is contained in some ultrafilter.
\end{theorem}
\begin{proof}
    Follows from Zorn's lemma.
\end{proof}
\begin{proposition}
    Let $\mathcal{U}$ be an ultrafilter of $X$ then all direct images of $\mathcal{U}$ are ultrafilters. If $\iota:A\to X$ is the inclusion of a subset then this induces,
    \begin{equation*}
        \{\text{ultrafilters on $A$}\}\stackrel{\iota_*}{\to}\{\text{ultrafilters on $X$ containing $A$}\},
    \end{equation*}
    which is a bijectin.
\end{proposition}
\begin{definition}
    \begin{equation*}
        \Lim\mathcal{F} = \{x\in X\sep \mathcal{N}_x\subseteq\mathcal{F}\}.
    \end{equation*}
    is the set of limits of $\mathcal{F}$. If it contains precisely one $x\in X$, we denote it by
    \begin{equation*}
        \Lim\mathcal{F} = x.
    \end{equation*}
\end{definition}
% \begin{definition}
%     \label{definition_9}
%     We say that $x\in X$ is a point of accumulation (denoteds by point of accumulation, sometimes called point of condensation) if for a neighbourhood $V\in\mathcal{N}_x$ and $M\in\mathcal{F}$, we have $V\cap M\not=\emptyset$. 
% \end{definition}
\begin{definition}
    \label{definition_9}
    We say that $x\in X$ is a point of accumulation (denoteds by point of accumulation, sometimes called point of condensation) if for every neighbourhood $V\in\mathcal{N}_x$ and every $M\in\mathcal{F}$, we have $V\cap M\not=\emptyset$. 
\end{definition}
\begin{proposition}
    Let $x\in X$ then
    \begin{equation*}
        \overline{\{x\}} = \Lim\mathcal{P}_x = \text{point of accumulation of the principal filter of $x$} = \Lim\mathcal{N}_x.
    \end{equation*}
\end{proposition}
\begin{proposition}Let $\F,\mathcal{G}$ be filters such that $\F\subseteq\mathcal{G}$ and $\U$ be a ultrafilter.
    \label{fact_11}
    \label{proposition_2_11}
    \begin{enumerate}
        \item Every limit of $\F$ is a point of accumulation of $\F$.
        \item Every limit of $\F$ is a limit of $\mathcal{G}$.
        \item Every point of accumulation of $\mathcal{G}$ is a point of accumulation of $\F$.
        \item A point $x\in X$ is a limit of $\U$ if and only if it is a point of accumulation of $\U$.
        \item The sets of its limits and its point of accumulations are closed for $\F$.
        \item Let $\iota:A\to X$ be a canonical inclusion where $A$ has an induced topology from $X$. Then $\Lim\F = A\cap\Lim\iota_*\F$.
    \end{enumerate}
    % Every limit of a filter $\mathcal{F}$ is a point of accumulation of $\mathcal{F}$. If $\mathcal{F}$ is a subset of some filter $\mathcal{G}$, then every limit of $\mathcal{F}$ is a limit of $\mathcal{G}$ and every point of accumulation of $\mathcal{G}$ is a point of accumulation of $\mathcal{F}$.
    % Furthermore if $\mathcal{U}$ is an ultrafilter, every point of accumulation of $\mathcal{U}$ is a limit of $\mathcal{U}$. For every filter, the set of its limit and of its points of accumulation points are closed. If $A\stackrel{\iota}{\to}X$ is an with induces topology then $\Lim\mathcal{F}=A\cap\Lim\iota_*\mathcal{F}$.
\end{proposition}
\begin{proof}
    For the third assertion, if $a$ is a point of accumulation of $\mathcal{U}$ and $V\in\mathcal{N}_x$, then $X\backslash V\not\in\mathcal{U}$ as $x$ is a point of accumulation. Hence $V\in\mathcal{U}$ by Proposition \ref{proposition_2_8}. For the forth assertion, let $x\not\in\Lim\mathcal{F}$, (respectively $x$ not a point of accumulation of $\mathcal{F}$) then there is $U\in\mathcal{N}_x$ such that $U\not\in\mathcal{F}$ (respectively, such that there is $M\in\mathcal{F}$ such that $U\cap M=\emptyset$). Then no elememt of $U$ is a limit (respectively, point of accumulation) of $\mathcal{F}$. 
\end{proof}
\begin{proposition}
    \label{proposition_1}
    For a subset $M$ of a topological space $X$, the following subsets all coincide.
    \begin{enumerate}[a).]
        \item $\{\text{limits of ultrafilters $\mathcal{U}$ of $X$ such that $M\in\mathcal{U}$}\}$,
        \item $\bigcup_{\substack{\text{$\mathcal{F}$ filter}\\\text{$M\in\mathcal{F}$}}}\Lim\mathcal{F}$. 
        \item $\{x\in X\sep \text{$\exists \mathcal{F}$ filter s.t. $M\in\mathcal{F}$ and $x$ is a point of accumulation of $\mathcal{F}$}\}$.
        \item $\bigcap_{\substack{M\subseteq Z\subseteq X\\\text{$Z$ closed}}}Z$.
    \end{enumerate}
\end{proposition}
\begin{proof}
    Denote those sets by $M_a,M_b,M_c,M_d$. Then,
    \begin{equation*}
        M_a\underbrace{\subseteq}_{\text{trivial}}M_b\underbrace{\subseteq}_{\text{Proposition \ref{proposition_2_11}}} M_c \underbrace{\subseteq}_{\text{definition}} M_d.
    \end{equation*}
    For the reverse direction, note that $M_d = \bigcup_{\substack{U\in\Ouv_X,\\U\cap M=\emptyset}}(X\backslash U)$. If $x\in M_d$, then 
    \begin{equation*}
        \{U\cap M\sep U\in\mathcal{N}_x\}
    \end{equation*}
    is a filter base containing $M$ and in particular containing $\mathcal{N}_x$. By Theorem \ref{theorem_tarski}, it is contained in some ultrafilter $\mathcal{U}$ and by definition $x\in\Lim\mathcal{U}$.
\end{proof}
\begin{definition}
    \label{definition_10}
    The set of above is called the closure of $M$ in $X$.
\end{definition}
%4/20
\begin{remark}
    \label{remark_2}
    The closure of a point s a special case of this.
\end{remark}
\begin{corollary}
    \label{corollary_1}
    For a subset $A$ of a topological space $X$, the following conditions are equivalent.
    \begin{enumerate}[i).]
        \item If $\F$ is a filter on $X$ such that $A\in\F$ and $x$ is a limit of $\F$ then $x\in A$.
        \item If $\F$ is a filter on $X$ such that $A\in\F$ and $x$ is a point of accumulation of $\F$ then $x\in A$.
        \item If $\mathcal{U}$ is a ultrafilter on $X$ such that $A\in\mathcal{U}$, every limit of $\mathcal{U}$ in $X$ is in $A$.
        \item $A$ is closed.
    \end{enumerate}
\end{corollary}
\begin{proof}
    Clear from Proposition \ref{proposition_1}
\end{proof}
\begin{proposition}
    \label{fact_12}
    For a map $f:X\to Y$ of topological spaces. The following conditions are equivalent.
    \begin{enumerate}[i).]
        \item $f$ is continuous.
        \item If $\F$ is a filter on $X$ and $x\in\Lim\mathcal{F}$ then $f(x)\in\Lim f_*\F$.
        \item If $\F$ is a filter on $X$ and $x\in X$ is a point of accumulation of $\F$, then $f(x)$ is a point of accumulation of $f_*\F$.
        \item If $\mathcal{U}$ is a ultrafilter on $X$ and $x\in\Lim\mathcal{U}$, then $f(x)\in\Lim f_*\mathcal{U}$.
    \end{enumerate}
\end{proposition}
\begin{proof}$\:$
    \par i) $\Rightarrow $ii). This follows from that $f_*\mathcal{N}_x=\mathcal{N}_{f(x)}$ and $f_*$ preserves inclusion.
    \par ii) $\Rightarrow$ iii). By Proposition \ref{proposition_2_11}, every limit is a point of accumulation. 
    \par iii)$\Rightarrow$iv). By Proposition \ref{proposition_2_11}, every point of accumulation of a ultrafilter is a limit of the filter and $f_*$ maps ultrafilters to ultrafilters.
    \par i)$\Rightarrow$iv). If $A\subseteq Y$ is closed, we want to show that $f^{-1}(A)$ is closed. Let $\U$ be an ultrafilter in $X$ such that $f^{-1}(A)\in\U$. Then $f_*\U$ is a ultrafilter on $Y$ containing $A$ and by Corollary \ref{corollary_1}, any limit of $f_*\U$ is in $A$. Since limits of $\U$ are mapped to limits of $f_*\U$, $f^{-1}(A)$ contains limits of $\U$. Again by Corollary \ref{corollary_1} as $f^{-1}(A)$ is closed.
\end{proof}
\begin{remark}
    \label{fact_13}
    By Definition \ref{definition_9}, we have the following.
    The set of points of accumulation of a filter $\F$ is the intersection of the closure of the elements of $\F$ or the intersection of all $A\in\F$ which are closed in $X$.
\end{remark}
\begin{proposition}
    \label{proposition_2}
    Let $X$ be a topological space. The following conditions are equivalent.
    \begin{enumerate}[i).]
        \item $X$ is Hausdorff.
        \item Every filter has at most one limit.
        \item Every ultrafilter on $X$ has at most one limit.
    \end{enumerate}
\end{proposition}
\begin{proof}$\:$
    \par i)$\Rightarrow$ii). If $x\not=y$, they have disjoint neighborhoods $x\in U$ and $y\in V$. If $U\in\F$ and $V\in\F$, then $\emptyset=U\cap V\in\F$, which is not possible. Thus precisely one of $x$ and $y$ is in $\Lim\F$. 
    \par ii)$\Rightarrow$iii) is trivial.
    \par iii)$\Rightarrow$i). Assume $x\not=y$ and consider the set 
    \begin{equation*}
    \{U\cap V\sep U\in\Nfil_x,V\in\Nfil_y\}.
    \end{equation*}
    This is a filter unless it contains $\emptyset$. In that case, it is contained in some ultrafilter $\U$. By definition of limits, $x,y$ are both limits of $\U$ thus $x=y$ which is a contradiction.
\end{proof}
\begin{proposition}
    \label{proposition_2_another}
    Let $X$ be a topological space. The following conditions are equivalent.
    \begin{enumerate}[i').]
        \item $X$ is quasi-compact.
        \item Every filter has at least one point of accumulation.
        \item Every ultrafilter has at least one limit.
    \end{enumerate}
\end{proposition}
\begin{proof}$\:$
    \par i')$\Rightarrow$ii') By Remark \ref{fact_13},
    \begin{equation*}
        \{x\sep x\text{ is a point of accumulation of $\F$}\} = \bigcap_{M\in\F}\overline{M}.
    \end{equation*}
    By quasi-compactness $X$ and applying it to the complement, this is a finite intersection of elements in $\F$. By definition, any filter is closed under finite intersection and does not contain the emptyset, this is not empty.
    \par ii')$\Rightarrow$iii'). Every filter is contained in some ultrafilter and limits and points of accumulation are the same in ultrafilter by Proposition \ref{proposition_2_11}.
    \par iii')$\Rightarrow$i'). Let $(A_\lambda)_\Lambda$ be a faimly of closed subsets such that for any finite subset $F\subseteq\Lambda$, we have,
    \begin{equation*}
        \bigcap_{\lambda\in F}A_\lambda\not=\emptyset.
    \end{equation*}
    Then the collection,
    \begin{equation*}
        \left\{\bigcap_{\lambda\in F}A_\lambda\:\bigg{|}\: F\subseteq\Lambda,\vert F\vert<\infty\right\}
    \end{equation*}
    is a filter base, thus it is contained in some ultrafilter $\U$. Let $x\in\Lim\U$. For $\lambda\in \Lambda$, $A_\lambda\in\U$ hence $x\in A_\lambda$ by Corollary \ref{corollary_1}. Hence $x\in\bigcap_{\lambda\in\Lambda}A_\lambda$ and this intersection is not empty.
    This is the definition of quasi-compactness stated in terms of closed sets.
\end{proof}
\begin{proposition}
    \label{proposition_3}
    If $X$ satisfies these equivalent conditions of Proposition \ref{proposition_2_5}, then for every ultrafilter $\U$ on $X$, the closed subset $\Lim\U\subseteq X$ is irreducible. 
\end{proposition}
\begin{proof}
    Let $\Omega,\Theta\in\QC(X)$ such that $\Omega\cap \Lim\U\not=\emptyset$ and $\Theta\cap\Lim\U\not=\emptyset$. By the definition of $\Lim\U$, $\Omega,\Theta\in\U$. Hence $\Omega\cap\Theta\in\U$.
    As $\Omega\cap\Theta$ is quasi-compact $\U\cap\Ouv_{\Omega\cap\Theta}$ has a limit $x\in\Omega\cap\Theta$ by Proposition \ref{proposition_2_another}. By Proposition \ref{proposition_2_11}, $x\in\Lim\U$. 
    \par If $U$ and $V$ are open subsets such that $U\cap\Lim\U$ and $V\cap\Lim\U$ are not emptysets, there are $\Omega\subseteq U$ and $\Theta\subseteq V$ as above (as $\QC(X)$ is a topology base), hence $U\cap V\cap\Lim\U$ is not empty. Finally by Proposition \ref{proposition_2_another}, $\Lim\U\not=\emptyset$ as $X$ is quasi-compact.
\end{proof}

\begin{definition}
    \label{definition_11}
    We say that $x\in X$ is a generic limit of a ultrafilter $\U$ if $\Lim\U=\overline{\{x\}}$.
\end{definition}

\begin{definition}
    \label{definition_12}
    Let $X$ be a topological space in which every ultrafilter $\U$ has a unique generic limit. Then we write,
    \begin{equation*}
        \glim\U = \glim_X\U
    \end{equation*}
    We say that a subset $M\subseteq X$ is closed under generic limit in $X$ if $M\in\U$ for every ultrafilter $\U$ implies $\glim\U\in M$.
\end{definition}

\begin{proposition}$\:$
    \label{fact_14_a}
    \begin{enumerate}[1).]
        \item $X$ is $T_0$ if and only if every ultrafilter on $X$ has at most one generic limit.
        \item If every ultrafilter on $X$ has some generic points then every closed subset $Z\subseteq X$ has some generic points.
    \end{enumerate}
\end{proposition}
\begin{proof}
    For 1). the sufficiency of being $T_0$ is by Definition \ref{definition_2}. If $x$ and $y$ are topologically indistinguishable, then $\overline{\{x\}} = \overline{\{y\}} =Z$ for some $Z$. Furthermore, the principal filter $\U$ of $x$ which has $\Lim\U = Z$ from Example \ref{example_3} has both $x$ and $y$ as generic limits.
    \par 2). Irreducibility of $Z$ implies that $\Ouv_Z\backslash\{\emptyset\}$ is a filter base in $X$. Hence it is contained in some ultrafilter $\U$. By Definition \ref{definition_10} and the closedness of $Z$ implies $Z=\lim\U$. Hence every $\glim\U$ is a generic point of $Z$. 
\end{proof}
\begin{proposition}
    \label{fact_14_b}
    Assume every ultrafilter $\U$ on $X$ has a unique generic limit $\glim\U$. Then the following statements hold.
    \begin{enumerate}[1).]
        \item  The set of subsets of $X$ which are closed under $\glim$ is closed under arbitrary intersection and finite unions.
        \item An open subset $V\subseteq X$ is closed under $\glim$ if and only if $V$ is quasi-comapct.
    \end{enumerate}
\end{proposition}
\begin{proof}
    For 1), the closedness under the intersections is obvious. If $M\cup N\in\U$ then $M\in\U$ or $N\in\U$ by Proposition \ref{fact_8}. Hence $\glim\U\in M\subseteq M\cup N$ or $\glim\U\in N\subseteq M\cup N$.
    \par Let $\mathcal{V}$ be an ultrafilter on $V$ and $\mathcal{U}$ be the image of $\mathcal{V}$ under the embedding $V\hookrightarrow X$. If $V$ is closed under $\glim$ then $\glim\mathcal{V}\in V$ is a limit of $\U$ by Proposition \ref{proposition_2_11}.
    By Proposition \ref{proposition_2}, $V$ is quasi-compact. If $V$ is quasi-compact, $\mathcal{V}$ has a limit $x\in V$ and $x\in\Lim\mathcal{V} = V\cap\Lim\mathcal{U}$ by Proposition \ref{fact_11}. Hence $x\in\overline{\{y\}}$ hence $y\in V$ and $V$ is closed under $\glim$.
\end{proof}

\begin{definition}
    A spectral space $X$ is quasi-separated if for any ultrafilter $\U$ on $X$, if its limit $\Lim\U$ is non-empty then it is an irreducible closed subset.
\end{definition}

\begin{proposition}$\:$
    \label{proposition_4}
    \begin{enumerate}[1).]
        \item A topological space $X$ is spectral if and only if every ultrafilter on $X$ has a unique generic limit and the set of those $\Omega\in\Ouv_X$ which are closed under $\glim$ is a topology base.
        \item An open subset of a spectral space $X$ is quasi-compact if and only if it is closed under $\glim$. (This is proved in Proposition \ref{fact_14_b}).
        \item A map $f:X\to Y$ is spectral if and only if it is continuous and $\glim_Y f_*\U = f(\glim_X\U)$ for every ultrafilter $\U$ on $X$. (This is proved by the above assertion).
    \end{enumerate}
\end{proposition}

\begin{proof}
    If $X$ is spectral, then by Proposition \ref{proposition_3}, $\Lim\U$ is always closed and irreducible subset, hence $\glim\U$ exists uniquely by the sobriety of $X$. 
    \par By Proposition \ref{fact_14_b}, one has $\{U\in\Ouv_X\sep \text{$U$ is closed under $\glim$}\}=\QC(X)$ is a topology base. If any ultrafilter has a unique generic limit, then $X$ is sober by Proposition \ref{fact_14_b} and $\mathcal{\Ouv_C}(X)=\{U\in\Ouv_X\sep \text{$U$ is closed under $\glim$}\}$ which is closed under finite intersection
    in $X$ and by the assumption this is a topology base for $X$.
\end{proof}

\begin{lemma}
    \label{lemma_1}
    Let $\mathcal{B}$ be a topology base of the topological space $X$.
    \begin{enumerate}[a)]
        \item Let $\U$ be an ultrafilter on $X$ and $g\in X$ such that for all $\Omega\in\mathcal{B}$, $g\in\Omega \Leftrightarrow\Omega\in\U$. Then $g$ is a generic limit of $\U$.
        \item If $X$ is $T_0$ and for every ultrafilter $\U$ on $X$, there is $g\in X$ as in a), then $X$ is spectral and the elements of $\B$ are quasi-compact.
    \end{enumerate}
\end{lemma}
\begin{proof}
    Let $U\in\Nfil_g$. Then there is $\Omega\in \B$ such that $g\in\Omega\subseteq U$. By our assumption, $\Omega\in\U$ hence $U\in\U$ thus $g$ is a limit of $\U$. 
    Let $x\not\in\overline{\{g\}}$ then there is $U\in\Nfil_x$ such that $g\not\in U$. As $\B$ is a topology base there is $\Omega\in\B$ such that $x\in\Omega\subseteq U$. Then $g\not\in \Omega$, hence $\Omega\not\in\U$. Since $\Omega$ is a neighbourhood of $x$, this shows $x\not\in\Lim\U$. Thus $g$ is a generic limit. This proves a).
    \par For b), it follows from our assumption and a) that $X$ has unique generic limits of ultrafilters and the elements of $\B$ are closed under such generic limits. The claim now follows from Proposition \ref{proposition_4}.
\end{proof}

\begin{theorem}
    \label{theorem_1}
    A topological space $X$ is spectral if and only if there is a ring $A$ such that $X$ is homeomorphic to $\Spec A$.
\end{theorem}

\begin{proof}
We omit the hard part which is $\Rightarrow$ direction (the existence of such ring $A$ for every spectral space $X$). That $\Spec A$ is spectral can be seen as follows. Let $\U$ be an ultrafilter on $\Spec A$ and set
\begin{equation*}
    \mathfrak{g}_\U \defeq \{a\in A\sep V(a)\in\U\} = \{a\in A\sep \{\mathfrak{p}\in\Spec A\sep a\in\mathfrak{p}\}\in\U\}.
\end{equation*}
We show $\mathfrak{g}_\U\in\Spec A$. Let $a_1,a_2\in\mathfrak{g}_\U$, and $M_i = \{\mathfrak{p}\in\Spec A\sep a_i\in\mathfrak{p}\}\in\U$ hence $M_1\cap M_2\in\U$, hence $\{\mathfrak{p}\sep a_1+a_2\in\mathfrak{p}\}\in\U$ as $\U$ is closed under finite intersection. Hence $a_1+a_2\in\mathfrak{g}_\U$ thus it is an ideal. Let $a_1a_2\in\mathfrak{g}_\U$, but $a_2\not\in\mathfrak{g}_\U$. Then $M_2\not\in\U$ but $M_1\cup M_2\in\U$ since $a_1a_2\in\mathfrak{p}$, one of them must belong to it. Hence 
\begin{equation*}
    M_1\cup M_2=\{\mathfrak{p}\sep a_1a_2\in\mathfrak{p}\}\in\U.
\end{equation*}
$M_2\not\in\U$ implies $M_1\in\U$ hence $a_1\in\mathfrak{g}_\U$. 
\par Since $X=V(f)\cap D(f)$, and that $f\not\in\g_\U$ if and only if $V(f)\not\in \U$ if and only if $D(f)\in\U$. By Lemma \ref{lemma_1}, $\mathfrak{g}_\U$ is $\glim\U$ and all $D(f)$ are $\glim$-closed.
\end{proof}
%4/24
\begin{remark}
    \label{remark_4_1}
    If $\U$ is an ultrafilter on $\Spec A$ and
    \begin{equation*}
        \g_\U = \{a\in A\sep \{\p\in\Spec A\sep a\in\p\}\in\U\}.
    \end{equation*}
    Then the lemma applies with $\B=\{D(a)\sep a\in A\}$.
\end{remark}
\begin{corollary}[Of Lemma \ref{lemma_1} b), and the proof of Theorem \ref{theorem_1}]
    \label{corollary_2}
    The subset $D(a)\subseteq \Spec A, a\in A$ are quasi-compact.
\end{corollary}
\begin{proof}
    From Remark \ref{remark_4_1} and Lemma \ref{lemma_1}.
\end{proof}
\begin{definition}
    Let $X$ be a topological space. A subset $U$ which is open and closed is called clopen.
\end{definition}
\begin{corollary}[Of Proposition \ref{proposition_4}]
    \label{corollary_3}
    For a compact (thus Hausdorff) X, the following conditions are equivalent.
    \begin{enumerate}[a).]
        \item $X$ is spectral.
        \item The clopen subsets of $X$ form a topology base.
    \end{enumerate}
    In this case, the open subsets $U$ is quasi-compact if and only if it is clopen and for a spectral space $Y$, a map $Y\to X$ is spectral if it is continuous.
\end{corollary}
\begin{proof}
    Every ultrafilter $\U$ of $X$ admits a unique limit by Proposition \ref{proposition_2} and Proposition \ref{proposition_2_another}, such limits are trivially generic limits. By Corollary \ref{corollary_1}, the open subsets which are invariant under $\glim$ are closed.% are preimages of the open subsets are also clesed.
    \par To show these two are equivalent. If $U$ is closed, it is quasi-compact as it is a closed subset of a quasi-compact space. If $U$ is quasi-compact then it is closed by Proposition \ref{fact_3}. If $Y\stackrel{f}{\to} X$ is continuous and $\Omega\in\QC(X)$, then $\Omega$ is clopen hence $f^{-1}(\Omega)$ is clopen in $Y$, hence $f^{-1}(\Omega)\in \QC(Y)$.
\end{proof}
\begin{definition}
    A set $\Alg$ of subsets of $X$ is a boolean algebra if the following conditions hold.
    \begin{enumerate}
        \item $\emptyset\in\Alg$.
        \item $M\in\Alg\Rightarrow X\backslash M\in\Alg$.
        \item $\Alg$ is closed under taking finite unions.
    \end{enumerate}
\end{definition}
\begin{remark}
    From the basic axioms of boolean algebras. Clearly we have $X\in \Alg$ and $M_1,M_2\in\Alg\Rightarrow M_1\cap M_2\in\Alg$. 
\end{remark}
\begin{proposition}
    \label{proposition_5}
    Let $X$ be a spectral space. The following boolena algebras of subsets of $X$ coincide.
    \begin{enumerate}
        \item The boolean algebra generated by $\QC(X)$.
        \item $\{M\subseteq X\sep \text{$M$ is closed under generic limits of ultrafilters}\}$.
    \end{enumerate}
    Let $\mathcal{K}_X$ be this boolean algebra of subsets of $X$. The topology on $X$ obtained by taking $\mathcal{K}_X$ as a topology base coincides with the one for which taking the subsets $M\subseteq X$ which are closed under $\glim$ of ultrafilters as closed subsets.
    \par Finally let $X_{\con}$ be $X$ equipped with this topology. Then $X_{\con}$ is a spectral hausdorff space (i.e. $X_{\con}$ as in Corollary \ref{corollary_2}), and $X_{\con}\stackrel{\id_X}{\to}X$ is a spectral map.
    \par For every ultrafilter $\U$ of $X$, $\lim_{X_{\con}}\U=\glim_X\U$. 
\end{proposition}
\begin{proof}[Proof:(Caution : This is not required for the exam)]$\:$
Let $\mathcal{K}_1$ be th Boolean algebra generated by $\QC(X)$ and 
\begin{equation*}
    \mathcal{K}_2 \defeq \{M\sep \text{$M$ and $X\backslash M$ closed under $\glim$}\}.
\end{equation*}
By Proposition \ref{fact_14_a}, $\mathcal{K}_2$ is a Boolean algebra. By Proposition \ref{proposition_4} b), and Corollary \ref{corollary_1}, $\QC(X)\subseteq\mathcal{K}_2$, hence we have an inclusion $\mathcal{K}_1\subseteq\mathcal{K}_2$.
For the time being, $\mathcal{K}_x\defeq\mathcal{K}_1$. Let $X_1$ be the topological space with underlying space $X$ and the topology generated by $\K_x$. Let $X_2$ be the topological space again with the same underlying set such that for a subset $A\subseteq X$ is closed if and only if it is closed under $\glim$ of ultrafilters.
By Proposition \ref{fact_14_a}, $X_2$ is indeed a topological space. Since $\K_x\subseteq \K_2$, $X_2$ is finer or equal to $X_1$.
\par The space $X_1$ is Hausdorff. If $x\not=y$ then without loss of generality exchanging $x$ and $y$ if necessary, there is $U\in\Ouv_X$ such that $x\in U$ and $y\not\in U$. As $\QC(X)$ is a topology base, hence without loss of $U\in\QC(X)$. Then $U\in\K_X$ and $X\backslash U\in\K_X$ and they are distinguished neighbourhoods of $x$ and $y$ in $X_1$.
Moreover, $X_2$ is quasi-compact by Proposition \ref{proposition_2}. Thus it suffices to consider an ultrafilter $\U$ on $X$ and to show that it has a limit in $X_2$. Let $g=\glim_X\U$. We claim that $g\in\Lim_{X_2}\U$. For this, let $U\in\Ouv_{X_2}$ such that $g\in U$, we must show $U\in\U$. Otherwise $X\backslash U\in\U$ But by the definition of $X_2$, $X\backslash U$ is closed undrer $\glim_X$ of ultrafilters. Thus $g\in X\backslash U$ which is a contradiction.
\par By Proposition \ref{fact_3}, the continuous map $X_2\stackrel{\id_X}{\to}X$ is a homeomorphism. Thus first two parts are shown.
\par Let $X_{\con}=X_1=X_2$. Let $\Omega\in\K_2$. By Corollary \ref{corollary_1}, $\Omega$ is clopen in $X_{\con}$ Since $\K_1$ is a topology base. $\Omega$ is a union of elements of $\K_1$ and since $\Omega\in\QC(X_{\con})$, the union can be chosen to be finite. Then $\Omega\in\K_1$ as $\K_1$ is a Boolean algebra of subsets.
\par From now on, we will regard $\K_x=\K_1=\K_2$. Since the elements of $\K_x$ are clopen and $\K_x$ is a topology base for $X_{\con}$, and $X_{\con}$ compact. By Corollary \ref{corollary_3}, $X_{\con}$ is spectral. By Proposition \ref{proposition_4}, $X_{\con}\stackrel{\id_X}{\to}X$ is spectral. 
\end{proof}

\begin{definition}
    \label{definition_13}
    The elements of $\K_X$ are called constructible subsets and the topology of $X_{\con}$ the constructible topology.
\end{definition}
\begin{corollary}
    \label{corollary_4}
    A subset of $X$ is an element of $\K_X$ if and only if it is clopen in $X_{\con}$.
\end{corollary}
\begin{corollary}[Of Proposition \ref{proposition_4}]
    A map $X\stackrel{f}{\to}Y$ between spectral spaces is spectral if and only if it is continuous and $X_{\con}\stackrel{f}{\to}Y_{\con}$ is also continuous.
\end{corollary}
\begin{corollary}
    \label{corollary_15}
    For a spectral space $X$, $\QC(X) = \Ouv_X\cap\K_X$.
\end{corollary}
\begin{definition}
    \label{definition_14}
    Let $X$ be a spectral. We say that $x\in X$ is a specialization of $y\in X$ and conversely $y$ is a generalization of $x$ if $x\in\overline{\{y\}}$. We denote such relation by $x\preceq y$.
\end{definition}
\begin{proposition}$\:$
    \label{proposition_6}
    \begin{enumerate}
    \item A subset $Z\subseteq X$ is closed in $X$ if and only if it is closed in $X_{\con}$ and closed under specialization. That is $x\preceq y,y\in Z\Rightarrow x\in Z$.
    \item A subset $U\subseteq X$ is open in $X$ if and only if it is open in $X_{\con}$ and closed under generalization.
    \item Let $Z$ be closed in $X_{\con}$ and $\overline{Z}$ be its closure in $X$. Then $\overline{Z}$ is the set of specialization of elements of $Z$.
\end{enumerate}
\end{proposition}
\begin{proof}
If $z\in Z$ and $x\preceq z$, then $x\in\overline{\{z\}}\subseteq\overline{Z}$. If $x\in\overline{Z}$ by Proposition \ref{proposition_1}, there is an ultrafilter $\U$ on $X$ such that $Z\in\U$ and $x\in\Lim\U$. Since $Z$ is closed in $X_{\con}$, $g=\glim\U\in\overline{Z}$. But $x\preceq g$, proving the third statement. The first two are equivalent and the first follows from the third.
\end{proof}
\begin{corollary}
    \label{corollary_6}
    For a spectral space $X$, the following statements are equivalent.
    \begin{enumerate}
        \item $X$ is $T_0$.
        \item $X=X_{\con}$ as topological space.
        \item Every quasi-compact subset of $X$ is closed.
        \item $X$ is Hausdorff.
        \item $X$ is $T_4$.
    \end{enumerate}
\end{corollary}
\begin{corollary}
    \label{corollary_7}
    For a map $X\stackrel{f}{\to}Y$ between spectral spaces, the following statements are equivalent.
    \begin{enumerate}
        \item $\glim_Yf_*\U=f(\glim_X\U)$ for all ultrafilters $\U$ on $X$, and $x\preceq_X \xi$ implies $f(x)\preceq_Y f(\xi)$. 
        \item $X_{\con}\stackrel{f}{\to}Y_{\con}$ and $x\preceq_X\xi$ implies $f(x)\preceq_Yf(\xi)$.
        \item $f$ is spectral.
    \end{enumerate}
\end{corollary}
\begin{corollary}
    \label{corollary_8}
    For a subset $Z\subseteq X$, the following statements are equivalent.
    \begin{enumerate}[a).]
        \item $Z$ is spectral with induced topology and $Z\to X$ is spectral.
        \item $X$ is closed in $X_{\con}$.
    \end{enumerate}
    In this case, we dentoe $Z_{\con}$ to be the topology on $Z$ induced from topology $X_{\con}$
\end{corollary}
\subsection{The Valuation Spectrum of a Rings}

\begin{definition}
    \label{definition_1_ordered_abelian_group}
    An ordered abelian group is an abelian group $\Gamma$ (under mulitplication) with a linear ordering $\leq$ such that 
    \begin{equation*}
        \gamma_1\leq \gamma_2, g_1\leq g_2\Rightarrow \gamma_1g_1\leq \gamma_2g_2.
    \end{equation*}
\end{definition}

\begin{example}
    $\{1\}$ and $(0,\infty)_\R^\times$ are ordered abelian groups. 
\end{example}

\begin{remark}
    The inequality,
    \begin{equation*}
        \gamma_1\leq \gamma_2, g_1\leq g_2\Rightarrow \gamma_1g_1\leq \gamma_2g_2.
    \end{equation*}
    is strict if one of the inequalities in the left hand side is strict.
    In particular, $g>1$ implies $g^k>1$ and $g<1$ implies $g^k<1$ for all $k\in\N$. Thus $\Gamma$ has no torsion.
\end{remark}

\begin{definition}
    \label{definition_1_valuation_rings}
    A (general) valuation of a ring $R$ is a map $R\stackrel{\gamma}{\to}\Gamma\cup\{0\}$ where $\Gamma$ is an ordered abelian group.
\end{definition}

\begin{definition}
    \label{definition_5_2}
    A valuation of a ring $R$ is a map $R\stackrel{\gamma}{\to}\Gamma\cup\{0\}$ where $\Gamma$ is an ordered abelian group with multiplication such that the following hold.
    \begin{enumerate}[i).]
        \item $\forall x\in\Gamma\cup\{0\}, 0\cdot x = x\cdot 0 = 0, 0\leq x$.
        \item $\gamma(xy) = \gamma(x)\gamma(y)$
        \item $\gamma(1) = 1$
        \item $\gamma(0) = 0$
        \item $\gamma(x+y)\leq \max\{\gamma(x),\gamma(y)\}$. 
    \end{enumerate}
\end{definition}

\begin{remark}
    \begin{equation*}
        \gamma(-1) = 1,\quad \gamma(y)<\gamma(x)\Rightarrow\gamma(x+y) = \gamma(x).
    \end{equation*}
\end{remark}

\begin{definition}
    \label{definition_5_3}
    Let $\Gamma_\gamma$ be the subgroup of $\Gamma$ generated by $\{\gamma(r)\sep r\in R\}\backslash\{0\}$.
    We say that valuations $\gamma$ and $\omega$ are equivalent if there is an order preserving isomorphism $\Gamma_\gamma\stackrel{i, \sim}{\to}\Gamma_\omega$ such that 
    \begin{equation*}
        \omega(r) = i(\gamma(r)),
    \end{equation*}
    where we put $i(0)=0$.
\end{definition}
\begin{proposition}
    \label{fact_5_1}
    The valuation of $r$ and $w$ are equivalent if and only if 
    \begin{equation*}
        \{(r,s)\in R^2\sep \gamma(r)\leq \gamma(s)\} = \{(r,s)\in R^2\sep w(r)\leq w(s)\}.
    \end{equation*}
\end{proposition}
\begin{definition}
    \begin{equation*}
        \supp\gamma = \{r\in R\sep \gamma(r) = 0\}.
    \end{equation*}
    \label{definition_5_4}
\end{definition}
\begin{remark}
    $\supp\gamma$ is a prime ideal. It stable under multipliation by any element since,
    \begin{equation*}
        \forall s\in R, \gamma(r)=0\Rightarrow\gamma(rs) = \gamma(r)\gamma(s) = 0.
    \end{equation*}
    It is closed under addition since,
    \begin{equation*}
        r,s\in\supp\gamma\rightarrow\gamma(r+s) \leq\max\{\gamma(r),\gamma(s)\} = 0.
    \end{equation*}
    Therefore it is an ideal. Clearly $1\not\in\supp\gamma$. If $\gamma(rs) = 0$, then $\gamma(r)\gamma(s) = 0$. Thus $\gamma(r) = 0$ or $\gamma(s) = 0$. 
\end{remark}
\begin{remark}
    The support only depends on the equivalence classes.
\end{remark}

\begin{example}
    If $\p$ is a prime ideal and 
    \begin{equation*}
        \gamma(r) = \begin{cases}
            1,\quad r\not\in\p,\\
            0,quad r\in\p.
        \end{cases}
    \end{equation*}
    the trivial valuation beloning to $\p$.
    \label{example_5_1}
\end{example}
\begin{definition}
    A valuation is trivial if it is equivalent to this kind of valuation. That is $\Gamma_\gamma=\{1\}$.
\end{definition}
\begin{example}
    \label{example_5_2}
    If $R=K$ is a field, then $\gamma(s) = 1$ where $\gamma^k=1$ for some $k\in\N$ as $\Gamma$ has no torsion. Then for a finite field only admits trivial valuations. 
\end{example}
\begin{example}
    \label{example_5_3}
    The $p$-adic valuation of $\Q$ or $\Q_p$ is a valuation.
\end{example}
\begin{definition}
    \label{definition_5_5}
    Let $\Spv(R)$ be the set of equivalent classes of valuation of $R$. For $x\in\Spv(R)$, we choose a representative $\vert\cdot\vert_x$.
    Let $\Gamma_x$  be the group generated by those $\vert r\vert_x$ which are non-zero. For $m\in\N$ and $(f_i)_{i=0}^m\in R^{m+1}$, let 
    \begin{align*}
        \mathcal{R}(f_0|f_1,\cdots,f_m)& =\mathcal{R}_{\Spv(R)}(f_0|f_1,\cdots,f_m),\\
        &\defeq \{x\in\Spv(R)\sep 0\not=\vert f_0\vert_x\geq\vert f_i\vert_x, 1\leq i\leq m\}.
    \end{align*}
    These are called rational open subsets. We equip $\Spv(R)$ with the topology for which the set of rational open subsets is a topology base. 
    \par If $A\stackrel{\alpha}{\to}B$ is a ring homomorphism then this induces 
    \begin{equation*}
        \Spv(B)\stackrel{\Spv(\alpha)}{\to}\Spv(A)
    \end{equation*}
    sends $x\in\Spv(B)$ to the equivalence class of the valuation,
    \begin{equation*}
        a\mapsto\vert\alpha(a)\vert_x
    \end{equation*}
    of $A$. We set,
    \begin{equation*}
        V(a) = \Spv(A)\backslash \mathcal{R}(a|).
    \end{equation*}
\end{definition}
\begin{remark}
    The topology is therefore the coarsest for which
    \begin{equation*}
        \Spv(A)\stackrel{x\mapsto\supp(x)}{\to}\Spec(A)
    \end{equation*}
    is continuous and all
    \begin{equation*}
        \{x\in\Spv(A)\backslash V(a)\sep \vert\alpha\vert_x\leq\vert a\vert_x\}
    \end{equation*}
    for $(a,\alpha)\in A^2$ are open.
\end{remark}
\begin{remark}
    The set of rational open subset is closed under finite intersections in $\Spv(R)$. Indeed,
    \begin{align*}
        \mathcal{R}(1|) &= \Spv(R),\\
        \mathcal{R}(r_0|r_1,\cdots,r_m)\cap\mathcal{R}(s_0|s_1,\cdots,s_n) &= \mathcal{R}(r_0s_0|r_is_j,0\leq i\leq n,0\leq j\leq n).
    \end{align*}
    If you assume $m\leq 1$ or even $m=1$, you only get a topology subbase.
\end{remark}
One way of constructing generic limits of ultrafilters in $\Spv(R)$ would be to use Proposition \ref{fact_5_1} and 
\begin{equation*}
    \{(r,s)\in R^2\sep \vert r\vert_g\leq\vert s\vert_g\}\defeq \{(r,s)\in R^2\sep \{x\in\Spv(R)\sep \vert r\vert_x\leq \vert s\vert_x\}\in\U\}.
\end{equation*}
This would work and was shown by Huber. This approache will not be taken in this course.
\begin{definition}
    \label{definition_5_6}
    Let $\U$ be an ultrafilter on $\Lambda$ and $(X_\lambda)_{\lambda\in\Lambda}$ be a family of sets indexed by $\Lambda$. We call the elements $(x_\lambda)_{\lambda\in\Lambda}$ and $(\xi_\lambda)_{\Lambda}$ of $\prod_{\lambda\in\Lambda}X_\lambda$, $\U$-equivalent if 
    \begin{equation*}
        \{\lambda\sep x_\lambda=\xi_\lambda\}\in\U.
    \end{equation*}
    The ultraproduct $\prod_{\U}X_\lambda$ is the set of $\U$-equivalence classes of elments of $\prod_{\lambda\in\Lambda}X_\lambda$. If $X_\lambda$-s have group structures, $(x\xi)\in\prod_\U X_\lambda$ is given as the $\U$-equivalece classes of $(x_\lambda\xi_\lambda)_\Lambda$. If $X_\lambda$-s are equipped with linear ordering, put 
    \begin{equation*}
        (x_\lambda)_{\lambda\in\Lambda}/\sim_\U \leq (\xi_\lambda)_{\lambda\in\Lambda}/\sim_\U
    \end{equation*}
    if $\{\lambda\in\Lambda\sep x_\lambda\leq\xi_\lambda\}\in\U$.
\end{definition}
\begin{proposition}
    \label{fact_5_2}
    Suppose $X_\lambda\not=\emptyset$ for all $\lambda\in\Lambda$.
    \begin{enumerate}
        \item If $\U=\{\Theta\subseteq\Lambda\sep\lambda_0\in\Theta\}$ is the principal ultrafilter belonging to $\lambda_0$, then $\prod_\U X_\lambda=X_{\lambda_0}$.
        \item If $F$ is a finite set and all $\lambda\in\Lambda,X_\lambda=F$. Then $\prod_\U F = F$ as for an arbitrary $x\in\prod_\U F$, $\Lambda = \bigcup_{f\in F}\{\lambda\sep x_\lambda =f\}$ (i.e. disjoint finite union) hence there is precisely one $f\in F$ such that $\{\lambda\sep x_\lambda=f\}\in\U$.
        \item If $X_\lambda = X_\lambda'\cup X_\lambda^{''}$ (disjoint unions, $\forall \lambda\in\Lambda, X_\lambda',X^{''}_\lambda\not=\emptyset$) Then $\prod_{\U}(X'_\lambda\cup X^{''}_\lambda) = (\prod_\U X_\lambda')\cup(\prod_\U X^{''}_\lambda)$ which is also a disjoint union as $\Lambda = \{\lambda\in\Lambda\sep x_\lambda\in X_\lambda'\}\cup\{\lambda\in\Lambda\sep x_\lambda\in X_\lambda^{''}\}$.
        \item Then $\prod_\U(\Gamma_\lambda\cup\{0\}) = (\prod_\U\Gamma_\lambda)\cup\{0\}$ and $\prod_\U\Gamma_\lambda$ is an ordered abelian group where all $\Gamma_\lambda$ are ordered abealian groups.
    \end{enumerate}
\end{proposition}
%4/30
\begin{remark}
    When some of the $X_\lambda$ may be empty, it is better to define $\prod_{\U}X_\lambda$ as the set of equivalent classes of pairs $(M,\xi)$ where $M\in\U$ is a subset and $\xi$ a function defined on $M$ such that $\xi(\lambda)\in X_\lambda$ where $\lambda\in M$. Two pairs $(M,\xi),(N,\nu)$ are equivalent if 
    \begin{equation*}
        \{\lambda\in M\cap N\sep \xi(\lambda) = \nu(\lambda)\}\in\U.
    \end{equation*}
    With this, we define,
    \begin{equation*}
        \prod_\U\left(\bigcup_{i=1}^m X_\lambda^{(i)}\right)\cong \bigcup_{i=1}^m\prod_\U X_\lambda^{(i)}.
    \end{equation*}
\end{remark}
\begin{remark}
    $\prod_\U(\text{ordered abelian group})$ is again an ordered abelian group and 
    \begin{equation*}
        \prod_\U (\Gamma_\lambda\cup\{0\}) \cong \left(\prod_\U\Gamma_\lambda\right)\cdot\{0\}.
    \end{equation*}
\end{remark}
\begin{proposition}
    \label{propostion_6_1}
    For every ring $R$, $\Spv(R)$ is a spectral space and a subset of the form $\cR(f_0|f_1,\cdots,f_m)$ with $m\in\N$ and $(f_i)_{i=0}^m\in R^{m+1}$ is quasi-compact.
    \par The generic limit of the ultrafiler $\U$ on $\Spv R$ is given by the equivalence class of the valuation,
    \begin{equation}
        R\stackrel{r\mapsto \left(\substack{\text{the $\U$-equivalence class of}\\\text{ $(\ab{r}_\gamma)_{\gamma\in\Spv R}$}}\right)}{\longrightarrow}\left(\left({\prod_{\gamma\in\Spv R}}\right)_\U\Gamma_\gamma\right)\cup\{0\}=\prod_\U(\Gamma_\gamma\cup\{0\}).\label{proposition_6_1_1}
    \end{equation}
    Moreover, the map,
    \begin{equation*}
        \Spv(R)\stackrel{\gamma\mapsto \supp\ab{\cdot}_\gamma}{\longrightarrow}\Spec R
    \end{equation*}
    is spectral and for every ring homomorphism $R\stackrel{\rho}{\to}S$, the map 
    \begin{equation*}
        \Spv(S)\stackrel{\Spv(\rho)}{\to}\Spv(R)
    \end{equation*}
    is also spectral.
\end{proposition}
\begin{proof}
    It is obvious from the definition and the properties of ultrafilters that \eqref{proposition_6_1_1} is indeed a valuation. To show that it defines a generic limit of $\U$ and that the rational open subsets are quasi-compact, we use Lemma \ref{lemma_1}. Set 
    \begin{equation*}
        \mathcal{B} = \{\cR(f_0| f_1,\cdots,f_m)\sep m\in\N,(f_i)_{i=0}^m \in R^{m+1}\}.
    \end{equation*}
    if $g$ denotes the equivalence class of \eqref{proposition_6_1_1}, 
    \begin{align*}
        [g\in\cR(f_0|f_1,\cdots,f_m)]&\Leftrightarrow\ab{f_0}_g\not=0 \text{ and all $\ab{f_1}_g\leq \ab{f_0}_g$},\\
        &\stackrel{\text{By \eqref{proposition_6_1_1}}}{\Leftrightarrow} (\ab{f_0}_x)_{x\in X}/\sim_\U \not=0\text{ and }(\ab{f_i}_x)_{x\in X}/\sim_\U\leq (\ab{f_0}_x)_{x\in X}/\sim_\U,\\
        &\stackrel{\substack{\text{Definition of $\prod_\U$}\\\text{ and the valuation on it}}}{\Leftrightarrow} \{x\in X\sep \ab{f_0}_x\not=0\}\in\U\text{ and  all }\{x\in X\sep \ab{f_i(x)}\leq f_0(x)\}\in\U,\\
        &\stackrel{\bigcap_{i=0}^m M_i\in\U\Leftrightarrow \forall M_i\in\U}{\Leftrightarrow} \{x\in X\sep \ab{f_0}_x\not=0\}\cap\bigcap_{i=1}^m \{x\in X\sep \ab{f_i(x)}\leq\ab{f_0(x)}\}\in\U,\\
        &\stackrel{\text{Definition of rational open}}{\Leftrightarrow}\cR(f_0\sep f_1,\cdots,f_m)\in\U.
    \end{align*}
    Lemma \ref{lemma_1} can indeed be applied. That the induced map $\Spv(R)\stackrel{\supp}{\to}\Spec(R)$ is spectral follows from 
    \begin{equation*}
        \supp(D(f)) = \cR(f|).
    \end{equation*}
    By Proposition \ref{facr1.2.6}, that $\Spv(\rho)$ is spectral follows immediately from 
    \begin{equation*}
        \Spv(\rho)^{-1}\cR(f_0|f_1,\cdots,f_m) = \cR(\rho(f_0)|\rho(f_1),\cdots,\rho(f_m)).
    \end{equation*}
\end{proof}

\begin{proposition}
    \label{fact_6_3}
    For every ideal $I\subseteq R$, if $R\stackrel{\pi}{\to}R/I$ denotes the canonical projection, then the induced map $\Spv(R/I)\stackrel{\pi}{\to} \Spv(R)$ is a homeomorphism. Explicitly, we have,
    \begin{align}
        \Spv(R/I) = V_{\Spv(R)}(I) &\defeq\{x\in\Spv(R)\sep \supp x\supseteq I\},\\
        & = \Spv(R)\backslash\bigcup_{f\in I}\cR(f).\label{eq_spr_v_formula}
    \end{align}
    If $I$ is finitely generated (i.e. $I = \langle g_1,\cdots,g_m\rangle_R$), then the right hand side of Equation \eqref{eq_spr_v_formula} can be written as,
    \begin{equation*}
        \Spv(R/I) = \bigcup_{i=1}^m \cR(g_i)
    \end{equation*}
    which is the embedded subset. Alternatively, we have,
    \begin{equation*}
        V_{\Spv(R)} = \{x\in\Spv(R)\sep \forall f\in I, \ab{f}_x = 0\}.
    \end{equation*}
    Furthermore, we can identify this with a homeomorphism,
    \begin{equation*}
        \Spv(R/I)\underset{\simeq}{\stackrel{\Spv(\pi)}{\to}} V(I).
    \end{equation*}
    \par Let $S$ be a multiplicative subset of $S$ and $\iota:R\to R_S$ be a canonical inclusion then,
    \begin{equation*}
        \Spv(R_S) = \underset{\Spv(\iota)}{\stackrel{\simeq}{\to}}\bigcap_{s\in S}(\Spv(R)\backslash V(s)) = \bigcap_{s\in S}\cR(s|).
    \end{equation*}
\end{proposition}

\begin{proposition}
    \label{fact_6_4}
    Let $s$ be a specialization of $\eta\in\Spv(R)$, then 
    \begin{equation*}
        \supp\eta\subseteq \supp s.
    \end{equation*}
\end{proposition}
\begin{proof}
    \begin{equation*}
        r\not\in\supp s\iff s\in\cR(r|)\Rightarrow \eta\in\cR(r|)\iff r\not\in\supp\eta.
    \end{equation*}
\end{proof}

\begin{definition}
    \label{definition_6_7}
    Let $\gamma\in\Gamma\cup\{0\}$ and $R\stackrel{v}{\to}\Gamma\cup\{0\}$ be a valuation. We then define,
    \begin{equation*}
    K_{v\leq \gamma}\defeq \{r\in R\sep v(r)\leq \gamma\}.
    \end{equation*}
    For an equivalence class $x\in\Spv(R)$ where $v\in x$, we may denote,
    \begin{equation*}
        K_{x\leq\gamma} = K_{v\leq\gamma}.
    \end{equation*}
    For $\rho\in R$, we denote,
    \begin{align*}
        K_{v\leq\rho} \defeq K_{v\leq v(\rho)}, K_{x\leq \rho} \defeq K_{x\leq\ab{\rho}_x}.
    \end{align*}
    Similarly,
    \begin{align*}
        K_{v<\gamma} &= K_{x<\gamma} \defeq\{r\in R\sep v(r)<\gamma\}\\
        K_{v<\rho} & = K_{x<\rho} \defeq K_{v<v(\rho)}.
    \end{align*}
\end{definition}

\begin{proposition}
    \label{fact_6_5}
    $K_{v\leq 1}$ is an integrally closed subring of $R$, and $K_{v<1}$ is a prime ideal in $K_{v\leq 1}$.
\end{proposition}
\begin{proof}
    These being a subring and an ideal can be checked easily. Let $r\not\in K_{r\leq 1}$, then if $m$ is a positive integer and $(p_i)_{i=1,\cdots,m}\in K_{v\leq 1}^m$, we have,
    \begin{align*}
    v(r^m) &= v(r)^m,\\
    &> v(r)^{m-1}\geq\max_{1\leq i\leq m-1}v(p_i)v(r)^i,\\
    & = \max_{1\leq i\leq m-1}v(p_ir^i),\\
    & \geq v\left(\sum_{i=1}^{m-1}p_ir^i\right).
    \end{align*}
    Hence,
    \begin{equation*}
        r^m \not = \sum_{i=1}^{m-1}p_ir^i.
    \end{equation*}
    Thus $r$ is not integral over $K_{v\leq 1}$, in particular, $K_{v\leq 1}$ is integrally closed.
    \par For $K_{v<1}$ being prime, obviously $1\not\in K_{v<1}$. Let $a,b\in K_{v\leq 1}\backslash K_{v<1}$. Then $v(a) = v(b) = 1$. Therefore $v(ab) = 1$. Showing $K_{v<1}$ is a prime ideal.
\end{proof}

\begin{proposition}
    \label{fact_6_6}
    Let $s$ be a specialization of $\eta\in\Spv(R)$. We have,
    \begin{equation*}
        a\in R\backslash\supp s\Rightarrow K_{\eta<a}\subseteq K_{s<a}.
    \end{equation*}
    Applying this to $a=1$, (recall that $\supp s$ is a prime ideal), we get,
    \begin{equation*}
        K_{\eta<1}\subseteq K_{s<1}.
    \end{equation*}
\end{proposition}

\begin{proof}
    By the assumption and Proposition \ref{fact_6_4}, we know that $\ab{a}_\eta\not=0$. Therefore,
    \begin{equation*}
        b\not\in K_{\eta<a}\iff\eta\not\in\cR(b|a)\Rightarrow s\not\in \cR(b|a)\iff b\not\in K_{s<a}.
    \end{equation*}
\end{proof}
\begin{proposition}
    \label{fact_6_7}
    Let $s$ be a specilization of $\eta\in\Spv(R)$. 
    \begin{equation*}
        a\in R\backslash\supp(s)\Rightarrow K_{s\leq a}\subseteq K_{\eta\leq a}.
    \end{equation*}
    In particular for $a=1$, we get,
    \begin{equation*}
        K_{s\leq 1}\subseteq K_{\eta\leq 1}.
    \end{equation*}
\end{proposition}

\begin{proof}
    By Proposition \ref{fact_6_4}, $\ab{a}_\eta\not=0$. Hence,
    \begin{equation*}
        b\in K_{s\leq a}\iff s\in \cR(a|b)\Rightarrow \eta\in\cR(a|b)\iff b\in K_{\eta\leq a}.
    \end{equation*}
\end{proof}
\begin{definition}
    \label{definition_6_8}
    Let $X$ be a linearly ordered set. For $a,b\in X$ with $a\leq b$, we define the line segment between $a$ and $b$ to be 
    \begin{equation*}
        [a,b]_X \defeq\{x\in X\sep a\leq x\leq b\}.
    \end{equation*}
    A subset $C\subseteq X$ is convex if for any two elements $c,d\in C,c\leq d$, we have $[ c,d]_X\subseteq C$.
\end{definition}

\begin{definition}
    A convex subgroup $\Theta$ is a subgroup of a linearly ordered group $\Gamma$ which is convex.
\end{definition}

\begin{proposition}
    \label{fact_6_8.5}
    A subgroup $\Theta$ of a linearly ordered group $\Gamma$ is convex if and only if for any $\theta\in \Theta$ with $\theta<1$, $[\theta,\theta^{-1}]_\Gamma\subseteq\Theta$.
\end{proposition}
Recall that for a valuation $v:R\to \Gamma\cup\{0\}$, $\Gamma_v$ is the subgroup of $\Gamma$ generated by $v(R)\backslash\{0\}$.
\begin{definition}
    \label{definition_6_9}
    $\cGamma_v$ is the convex subgroup of $\Gamma_v$ generated by $\{v(r)\sep r\in R, v(r)\geq1\}$. 
\end{definition}
\begin{remark}
    Equivalently, we can define $\cGamma_v$ to be a subgroup generated by $\{v(r)\sep r\in R, v(r)>1\}$.
\end{remark}

\begin{proposition}
    \label{fact_6_10}
    \begin{equation*}
        \cGamma_v = \{\gamma\in\Gamma_v\sep \exists r\in R \:\st\: \min\{\gamma,\gamma^{-1}\} v(r)\geq1\}
    \end{equation*}
\end{proposition}
\begin{proof}
    Let $\cGamma_v' =  \{\gamma\in\Gamma_v\sep \exists r\in R \:\st\: \min\{\gamma,\gamma^{-1}\} v(r)\geq1\}$. We will show that $\cGamma_v=\cGamma_v'$.
    If $r\in R$ is as in the construction of $\cGamma_v'$, then $v(r)\geq1$ thus $v(r)\in\cGamma_v$. If $\gamma\in\cGamma_v'$, then 
    \begin{equation*}
        v(r)^{-1}\leq \gamma\leq v(r).
    \end{equation*}
    Hence $\gamma\in\cGamma_v$. In particular $\cGamma_v'\subseteq\cGamma_v$. On the other hand $\cGamma_v'$ is a convex subgroup of $\Gamma_v$ containing all $v(r)$ such that $v(r)\geq1$. Thus by the minimality of $\cGamma_v$, we have shown the statement.
\end{proof}
\begin{remark}
    Even if $H$ is a convex subgroup of $\Gamma_v$, it does not imply that $\cGamma_v\subseteq H$.
\end{remark}

%5/11

\begin{remark}
    If $R$ is a field then $\cGamma_v=\Gamma_v$.
\end{remark}

\begin{definition}
    \label{definition_7_10}
    Let $A\stackrel{v}{\to}\Gamma$ be a valuation of $A$ and $H\subseteq\Gamma_v$ a convex subgroup containing $\cGamma_v$. Let 
    \begin{equation*}
        v|_{H}(r)\defeq\begin{cases}
            v(r)\quad(v(r)\in H)\\
            0\quad(\text{otherwise}).
        \end{cases}
    \end{equation*}
\end{definition}

\begin{proposition}
    \label{fact_7_8}
    Under the assumption of Definition \ref{definition_6_8}, $v|_H$ is indeed a valuation of $A$ and $v|_H$ is a specialization of $v$ in $\Spv(A)$.
\end{proposition}
\begin{proof}
    Let $M$ be the convex hull of $\{0\}\cup H$ in $\Gamma_v\cup\{0\}$. For $m\in M$, set,
    \begin{equation*}
        \varphi(m) =\begin{cases}
            m,\quad(m\in H),\\
            0,\quad (\text{otherwise}).
        \end{cases}
    \end{equation*}
    Since $H\supseteq\cGamma_v$, we have $v(a)\in M$ for all $a\in A$. (note that $[0,1]_{\Gamma_v\cup\{0\}}\supseteq M$ and $v(a)\in\cGamma_v\subseteq H\subseteq M$ where $v(a)\geq1$ )
    Set $w\defeq v|_H = \varphi v$. Since $v$ is monotonic, $\varphi(0) = 0$, and $\varphi(1) = 1$, the axioms
    \begin{equation*}
        w(0) = 0,w(1) = 1,w(a+b)\leq\max\{w(a),w(b)\}
    \end{equation*}
    follows from their counterparts for $v$. Also $M$ is multiplicativly closed and $\varphi(m\mu) = \varphi(m)\varphi(\mu)$. (This is obvious when $\varphi(m)\varphi(\mu)\not=0$. Otherwise, without loss of generality $\varphi(m)=0$, hence $m\not\in H$. In particular, $m<1$. Also $hm<1$ for all $h\in H$ by the convexity of $H$. Thus, $nm<1$ for all $n\in M$, in particular $m\mu<1$. Now $m\mu\in H$ would imply $\mu\in H$ as $m\mu\leq \mu$ and $\mu\in M$ (hence $\mu$ is bounded from above by some $h\in H$). But then $m\in H$ a contradiction). 
    \par To see $w$ is a specialization of $r$, observe that 
    \begin{align*}
        w\in\cR(f_0|f_1,\cdots,f_m) &\iff [w(f_0)\not=0\text{ and } w(f_0)\geq w(f_i),i=1,\cdots,m],\\
        &\iff[v(f_0)\in H\text{ and }v(f_0)\geq w(f_i),i=1,\cdots,m],\\
        &\iff [v(f_0)\in H\tand v(f_0)\geq v(f_i),i=1,\cdots,m],\\
        &\iff [v(f_0)\not=0\tand v(f_0)\geq v(f_i),i=1,\cdots,m]\\
        &\iff (v/\sim)\in\cR(f_0|f_1,\cdots,f_m).
    \end{align*}
    Note that $v(f_i)>v(f_0)\in H$ would imply $v(f_i)\in H$ by the convexity of $H$, as $v(f_i)\in H$, hence $v(f_i)$ bounded from above by some elements of $H$. Hnce $w(f_i)=v(f_i)>v(f_0) = w(f_0)$ if $v(f_i)>v(f_0)$.
\end{proof}

\begin{remark}
    Motivated by the above notion, we say $s\in\Spv(R)$ is a horizontal specialization of $\eta\in\Spv(R)$ if there is $v\in\eta$ and a convex subgroup $\cGamma_v\subseteq H\subseteq\Gamma_v$ such that $s$ is equivalent to $v|_H$.
\end{remark}

\begin{corollary}
    \label{fact_7_9}
    If $x$ is a closed point of $\Spv(A)$, then $\cGamma_x=\Gamma_x$. 
\end{corollary}
\begin{proof}
    If this is not the case, $x|_{\cGamma_x}$ is a proper specialization of $x$.
\end{proof}

\begin{proposition}
    \label{fact_7_10}
    If $s\in\Spv(A)$ is a horizontal specialization of $\eta\in\Spv(A)$ and $\Omega\in\cR(f_0|f_1,\cdots,f_m)$ then $s\in\Omega\iff(\eta\in\Omega\tand\ab{f_0}_s\not=0)$.
\end{proposition}

\begin{proof}
    Indeed, the construction of the end of the proof of Proposition \ref{fact_7_8} imply 
    \begin{equation*}
        (w/\sim)\in\Omega\iff[(v/\sim)\in\Omega\tand w(f_0)\not=0].
    \end{equation*}
\end{proof}

\begin{definition}
    Let $A\stackrel{v}{\to}\Gamma\cup\{0\}$ be a valuation and a convex subgroup $H\subseteq \Gamma$ or $H\subseteq\Gamma_v$. The vertical specialization of $v$ with respect to $H$ is
    \begin{equation*}
        (v/H)(a)=\begin{cases}
            v(a)H, \quad(v(a)\not=0),\\
            0,\quad(\otherwise),
        \end{cases}
    \end{equation*}
    where $v(a)H$ denotes the coset represented by $v(a)$.
\end{definition}

\begin{remark}
    $s$ is a vertical specilization of $\eta\in\Spv(A)$ if and only if $s$ is a specialization of $\eta$ and $\supp s = \supp \eta$. 
\end{remark}
\begin{remark}
    In Huber's literature, he refers horizontal specialization as primary specialization and the vertical specializations as the secondary specializations.
\end{remark}

\begin{remark}
    Note that $v|_H$ is a specialization of $v$ while $v$ is a specialization of $v/H$.
\end{remark}

\subsection{Valuation Rings}

\begin{definition}
    Let $K$ be a field. A subring $R\subseteq K$ is a valuation ring of $K$ if for all $x\in K^\times, x\in R\tor x^{-1}\in R$.
\end{definition}

\begin{example}
    Every discrete valuation rings is a valuation ring of its field of fraction. More concretely, $\Z_p$ is a valuation ring of $\Q_p$. 
\end{example}
\begin{example}
    If $S$ is a valuation ring of $L$ and $K\subseteq L$ is a subfield then $S\cap K$ is a valuation ring of $K$.
\end{example}

\begin{example}
    If $\mathcal{L}$ is a set of valuation rings which is linearly ordered by $\subseteq$. Then $\bigcap_{R\in\mathcal{L}}R$ is a valuation ring of $K$ (if $\mathcal{L}=\emptyset$ then the intersection would be $K$ itself which is trivially a valuation ring) and if $\mathcal{L}\not=\emptyset$ then $\bigcup_{R\in\mathcal{L}}R$ is also a valuation ring of $K$.
\end{example}

\begin{example}
    If $R$ is a valuation ring of $K$ and $R\subseteq S\subseteq K$ then $S$ is a valuation ring of $K$.
\end{example}

\begin{proposition}Let $R$ be a valuation ring of $K$. We have the following.
    \label{fact_7_1}
    \begin{enumerate}[a).]
        \item The field of fractions of $R$ is $K$. 
        \item For $a,b\in R$, either $a\mid b$ or $b\mid a$. 
        \item $R$ is a Bézout domain. That is every finitely generated ideals are principal. 
        \item $R$ is integrally closed in $K$. 
        \item The following conditions are equivalent for $R$
        \begin{enumerate}[i).]
            \item $R$ is Noetherian.
            \item $R$ is a discrete valuation ring.
            \item The unique maximal ideal $\m_R$ is principal.
            \item $R$ is a principal ideal domain.
        \end{enumerate}
        \item The set of ideals of $R$ is linearly ordered by the inclusion.
        \item $\Spv(R)$ is linearly orded by specizalization.
        \item An ideal $I\subseteq R$ is prime if and only if $I=\sqrt{I}$ and $1\not\in I$.
        \item We have an order reversing bijection such that 
        \begin{center}
        \begin{tikzcd}
\Spec(R) \arrow[r, "\p\mapsto R_\p", shift left] & \{S\sep R\subseteq S\subseteq K\tand S\text{ is a subring of $K$}\} \arrow[l, "R\cap\m_S\mapsfrom S", shift left]
\end{tikzcd}
        \end{center}
        \item The set of $S$ such that $R\subseteq S\subseteq K$ and $S$ is a subring of $K$ is linearly ordered by inclusion.
    \end{enumerate}
\end{proposition}

\begin{proposition}
    \label{fact_7_2}
    If $R$ and $S$ are valuation ring of $K$ then $R\subseteq S$ is equivalent to $\m_S\subseteq\m_R$.
\end{proposition}
\begin{proof}
    \begin{equation*}
        \m_R=\{0\}\cup\{x\in K^\times\sep x\not\in R\}.
    \end{equation*}
\end{proof}
\begin{proposition}
    \label{proposition_7_1}
    We have following bijections,
    \begin{center}
        \begin{tikzcd}
\Spv(K) \arrow[r, "x\mapsto \{k\in K\sep \ab{k}_x\leq1\}", shift left] & \{\text{valuation rings $R$ of $K$}\} \arrow[l, "v_R/\sim\mapsfrom R", shift left]
\end{tikzcd}
    \end{center}
    where 
    \begin{equation*}
        v_R(x) = \begin{cases}
            x\cdot R^\times ,\quad(x\not=0),\\
            0,\quad(x=0).
        \end{cases},\Gamma_R = K^\times/R^\times.
    \end{equation*}
    The ordering on $\Gamma_R$ is given as follows,
    \begin{equation*}
        xR^\times\leq yR^\times \iff x\in yR.
    \end{equation*}
\end{proposition}

%5/15 

\begin{remark}
    \label{fact_8_2}
    Under the bijection in Proposition \ref{proposition_7_1}, if the embedding of $K$ into a field extension $L$ is denoted $i$, $\Spv(i)(x_R) = x_{K\cap R}$ where $R$ is a valuation ring of $L$, $x_R\in\Spv(L)$. Moreover, $x_R\in\cR(f_0|f_1,\cdots,f_m)$ if and only if $f_0=0$ and $f_i\in f_0\cR$. 
\end{remark}

\begin{proposition}
    \label{fact_8_3}
    $x_R$ be a specialization of $x_S$ if and only if $R\subseteq S$. 
\end{proposition}

\begin{proof}$\:$
    \par If $R\subseteq S$, then by Remark \ref{fact_8_2} and
    \begin{align*}
        x_R\in\cR(f_0|f_1,\cdots,f_m) &\iff [f_0\not=0\tand \forall 1\leq i\leq m, f_i\in f_0R]\\
        &\Rightarrow [f_0\not=0\tand \forall f_i\in f_0S],\\
        &\iff [x_S\in\cR(f_0|f_1,\cdots,f_m)].
    \end{align*}
    If $x_R$ is a specialization of $x_S$, then $R = K_{x_R\leq 1}\subseteq K_{x_S\leq 1}=S$ by Proposition \ref{fact_6_7}.
\end{proof}


\begin{lemma}
    \label{lemma_8_aux}
    Let $B$ be an $A$-algebra and $M\subseteq B$ be a finitely generated $A$-submodule of $B$ such that $bM=0$ for any $b\in B$. Then it implies $aM=0$. If $b\in B$ is such that $bM\subseteq M$ then $b$ is integral over $A$. 
\end{lemma}
\begin{proof}
    Let $(m_i)_{i=1}^l$ be generators of $M$ as an $A$-module. 
    \begin{equation*}
        bm_i = \sum_{i=1}^l \alpha_{ij}m_j.
    \end{equation*}
    Set $p\in A[T]$ be the characteristic polynomial of $(\alpha_{ij})$. Then $p(b) = 0 $ by Cauchy-Hamilton.
\end{proof}

\begin{proposition}
    \label{proposition_8_2}
    Let $A$ be a subring of the filed $K$. Then the integral closure $A$ in $K$ is equal to the intersection of the valuation rings of $K$ containing $A$.
\end{proposition}

\begin{proof}
    Let $A'$ be the intersection of all valuation rings containing $A$. $A'\supseteq A$ is clear as all valuation rings are integrally closed. Thus it is sufficient to consider an element $x\in K$ which is not integral over $A$ and show there is a valuation ring $R$ of $K$ which doesn't contain $x$.
    To this end, set 
    \begin{equation*}
        \mathcal{M}\defeq \{A\subseteq B\subseteq K\sep\text{$B$ is an intermediate ring with $x$ is not integral over $B$}\}.
    \end{equation*}
    Since $A\in\mathcal{M}$, $\mathcal{M}\not=\emptyset$. Let $\mathcal{L}\subseteq\mathcal{M}$ be a non-empty set which is linearly ordered with respect to inclusion. Then $C=\bigcup_{B\in\mathcal{L}}B$ is in $\mathcal{M}$. That is $\mathcal{L}$ is bounded in $\mathcal{M}$.
    By Zorn's lemma, there is a maximal element $R$. It remains to show that $R$ is a valuation ring.
    \par Let $k\in K\backslash R$ such that $k^{-1}\not\in R$. Then $x$ is integral over $R[x^{\pm}]$, hence there are polynomials $(p_j^{\pm})_{j=0}^{d_{\pm}-1}$ such that 
    \begin{equation*}
        x^{d_\pm} = \sum_{j=0}^{d_\pm-1}x^jp_j^{(\pm)}(k^{\pm1}).
    \end{equation*}
    Let $d = \max\{d_+,d_-\}$ and $e$ the maximum of all $\deg p_j^{(\pm)}$ with $0\leq j<d^{(\pm)}$. Let $N\subseteq K$ be the $R$-submodule generated by the $x^ik^j$ where $0\leq i<d, 0\leq j< 2e$. Then $N$ is a finitely generated $R$-module
    and $xN\subseteq N$. Indeed, consider a generator $y=x^ik^j$ of $N$. We can assume that $j<e$. If $i<d_+-1$ then $xy$ is in the set generators of $N$. Otherwise, 
    \begin{equation*}
        xy = x^{i+1-d_+}\sum_{l=0}^{d_+-1}x^lp_l^{(+)}(k)k^j\in N.
    \end{equation*}
    Now assume $j\geq e$. If $i<d-1$, then $xy$ is in the set of generators of $N$. Otherwise, 
    \begin{equation*}
        xy = x^{i+1-d_-}\sum_{l=0}^{d_--1}x^lk^j p_l^{(-)}(k^{-1})\in N.
    \end{equation*}
    By Lemma \ref{lemma_8_aux}, $x$ is integral over $R$. This contradicts that $R\in\mathcal{M}$.
\end{proof}

\begin{remark}
    The book of Kniebuch and Schrecker also has a proof. 
\end{remark}

\section{$G_+$-topological span and their sheaf cohomology}
\subsection{Grothendieck topologies and $G_+$-topological span}
Let $\sC$ be a category. In the application we are interested in, it will be given by a poset $(P,\preceq)$. Objects of $\sC$ are the all of $P$ and a morphism $x\stackrel{\mu}{\to}x$ exists if and only if $x\preceq y$, in which case $\mu$ is unique. In fact $P=\Ouv_X$ (all open subsets of $X$) or $P$ is some topology base of $X$, and $\preceq$ is containment $\subseteq$. However, the case where therer may be more than one morphism between objects of $\sC$ are very important in other applications like étale or crystalline cohomology.
\begin{definition}
    \label{definition_8_1}
    A sieve on (or over) an object $X$ of $\sC$ or an $X$-sieve is a class $\mathscr{S}$ of morphisms $U\to X$ such that $[U\stackrel{\nu}{\to}X]\in\mathscr{S}$ then $\nu\omega\in\mathscr{S}$ where $V\stackrel{\omega}{\to}U$ is any morphism in $\sC$.
\end{definition}
\begin{definition}
    \label{definition_8_1_1}
    If $\sS$ is a $X$-sieve and $Y\stackrel{\xi}{\to}X$ then $\xi^*\sS$ is the $Y$-sieve of all morphism $U\stackrel{\nu}{\to}Y$ such that $\xi\nu\in\sS$.
\end{definition}
\begin{definition}
    \label{definition_8_2}
    A Grothendieck topology on $\sC$ is a collection $(\mathcal{T}_X)_{X\in\Ob(\sC)}$ of $X$-sieves (where $\mathcal{T}_X$ being called the class of covering sieves of $X$) satisfying the following conditions.
    \begin{enumerate}[I).]
        \item (GT-Triv) The all-sieve $\sS$ on $X$ (i.e. the class of all morphisms with its target to $X$) is in $\mathcal{T}_X$.
        \item (GT-Trans) If $\sS\in\mathcal{T}_X$, then for any morphism $Y\stackrel{\nu}{\to}X$ in $\sC$, we have $ \nu^*\sS\in\mathcal{T}_Y$.
        \item (GT-Loc) Let $\mathscr{T}$ be a $X$-sieve. If there is $\sS\in\mathcal{T}_X$ such that $\nu^*\mathscr{T}\in\mathcal{T}_U$ for any morphism $(U\stackrel{\nu}{\to}X)\in\sS$, then $\mathscr{T}\in\mathcal{T}_X$.
    \end{enumerate}
\end{definition}
\begin{remark}
    Recall that a presheaf on $\sC$ is a contravariant functor $\sC\stackrel{\mathcal{G}}{\to}(\text{Sets})$. The idea is that $\mathcal{G}$ is a sheaf if and only if $\mathcal{G}(X)\to\lim_{\sS}\mathcal{G}$ where morphisms in $\sS$ are morphisms of $X$-objects, is an isomorphism for all $\sS\in\mathcal{T}_X$.
\end{remark}
\begin{remark}
    If $\sC$ is given by a poset, then $\sS$ is also a poset (of $\preceq$-predecessor of $X$) and 
    \begin{equation*}
        \varprojlim_{\sS}\mathcal{G}=\left\{(g_Y)_{Y\in\sS}\in\prod_{Y\in\sS}\mathcal{G}(Y)\:\bigg{|}\: g_Y|_Z = g_Z\text{ where }Z\preceq Y\right\}.
    \end{equation*}
    Here $\mathcal{G}(\text{the unique morphisms $Z\to Y$})$ are denoted by $\cdot|_Z$.
\end{remark}
% \begin{remark}
%     In general (fpqc-topology or proétale topology), set theoretic patterns arise which we can ? because they were $\sC$ is gien by a poset $P$ (all all sieves the being subsets of $P$).
% \end{remark}
\begin{proposition}$\:$
    \label{fact_8_1} 
    \begin{enumerate}
        \item If $\sS\subseteq\mathscr{T}$ and $\sS\in\mathcal{T}_X$, then $\mathscr{T}\in\mathcal{T}_X$. 
        \item If $\sS_1,\sS_2\in\mathcal{T}_X$, then $\sS_1\cap \sS_2\in\mathcal{T}_X$.
    \end{enumerate}
\end{proposition}
\begin{proof}
    If $[U\stackrel{\nu}{\to}X]\in\sS$,
    \begin{equation*}
    \nu^*\sS = \{\sigma:V\to U\sep \nu\sigma: V\to X\in\sS\}.
    \end{equation*}
    Thus for any $\sigma:V\to U\in\sC$, $\sigma\in\nu^*\sS$ by using the definition of sieves. Thus $\nu^*\sS$ is the all-sieve on $U$. Observing that $\nu^*\mathscr{T}\supseteq \nu^*\sS$, hence $\nu^*\mathscr{T}=(\text{all sieve on $U$})\in\mathcal{T}_U$ by GT-Triv.
    By GT-Loc, $\mathscr{T}\in\mathcal{T}_X$. 
    \par Let $\mathscr{T}=\sS_1\cap \sS_2$. If $[U\stackrel{\nu}{\to}X]\in\sS_1$, as in the previous argument $\nu^*\mathscr{S}_1$ is the all-sieve, thus $\nu^*\mathscr{T}=\nu^*\sS_2\in\mathcal{T}_U$. By GT-Trans and $\nu^*\sS_2=\mathscr{T}\in\mathcal{T}_X$. By GT-Loc $\mathscr{T}\in\mathcal{T}_X$.
\end{proof}
From now on, $\sC$ will be given by some topology base $\mathcal{B}$ in a topology space $X$ turned into a poset by the inclusion. We then give a more specific definition for sieves in $\sC$.
\begin{definition}
    Let $\sC$ be a category whose objects are open subsets of $X$ and morphisms are inclusions.
    \begin{enumerate}[1).]
        \item A sieve $\sS$ on $U\in\sC$ is a set of open subsets of $U$ such that for any $V\in\sS$, every open subsets $W$ of $V$ is also in $\sS$.
        \item Let $V\subseteq U$ and $\sS$ be a sieve on $U$, then $\sS|_V = \sS\cap\{\text{open subsets of $V$}\}$.
    \end{enumerate}
\end{definition}
\begin{definition}
    \label{definition_8_3}
    Let $X$ be a topological space and $\cB_X$ be the topology base of $X$. Let $(U_i)_{i\in I}$ be a collection of elements of $\cB_X$. Then,
    \begin{align*}
        [U_i|i\in I]&\defeq \bigcap_{\substack{\text{$X$-sieve $\mathscr{S}$}\\ \st\:\forall i\in I,U_i\in\sS}}\sS,\\
        & =\{U\in\mathcal{B}_X\sep \exists i\in I\:\st\: U\subseteq U_i\}.
    \end{align*}
    % where all $U_i\in\mathcal{O}_W=\{\Omega\in\mathcal{O}\sep \Omega\subseteq W\}$.
\end{definition}
\begin{remark}
    In order to see the two definitions given above note that the latter is a sieve and contains all $U_i$ and indeed a sieve on $X$. Note that this is a sieve in the case $I$ is finite by Proposition \ref{fact_8_1}.
\end{remark}
\begin{definition}
    A $W$-sieve is finitely generated if it has the form $[\Omega_1,\cdots,\Omega_n]$ where all $\Omega\in\mathcal{B}_W$. If several topology bases are under consideration, we will use $[\Omega_1,\cdots,\Omega_n]_\mathcal{B}$ for disambiguation.
    \par Let $U\in\Ouv_X$. $U$-sieve $\sS$ is admissible (denoted by $\sS\vDash U$) by the Grothendieck group $(\mathcal{T}_U)_{U\in\Ouv_X}$ if $\sS\in\mathcal{T}_U$. Also for an arbitrary sieve $\sS$, we write $\sS//U$ if $\sS|_U\vDash U$.
\end{definition}
%5/18
% \begin{definition}
%     \label{definition_9_3}
%     We denote,
%     \begin{equation*}
%         [U_1,\cdots,U_n] = (\text{$\Ouv_U$-sieve generated by $U_i$}) = \bigcup_{i=1}^n\Ouv_{U_i}.
%     \end{equation*}
%     If we specify the topology base $\mathcal{B}$, we denote,
%     \begin{equation*}
%         [\Omega_1,\cdots,\Omega_n]_{\mathcal{B}}=(\text{$\mathcal{B}$-sieves generated by $\Omega_i$}) = \bigcup_{i=1}^n\mathcal{B}_{\Omega_i}.
%     \end{equation*}
%     A sieve is finitely generated if it can be generated in this way by finitely many of its elements. 
%     \begin{equation*}
%         \sS\vDash  U \iff \sS \text{ is a $U$-sieve and $\sS\in\mathcal{T}_U$}.
%     \end{equation*}
%     Also,
%     \begin{equation*}
%         \sS//U\iff \sS|_U\vDash U.
%     \end{equation*}
% \end{definition}


\begin{proposition}
    \label{proposition_9_1}
    Let $X$ be a topological space and $\cB$ be a topology base for it. Then we have a bijection between the following two sets.
    \begin{enumerate}[1).]
        \item Grothendieck topologies $\mathcal{T}$ on $\cB$.
        \item Grothendieck topologies $\mathcal{T}$ on $\Ouv_X$ with the additional property that $[\cB_U]\in\mathcal{T}_U$ holds for all $U\in\Ouv_X$.
    \end{enumerate}
    The bijection is explicitly given by the following way.
    \begin{enumerate}[a).]
        \item If a Grothendieck topology $\mathcal{T}^{(\cB)}$ on $\cB$ is given, $\sS\in\mathcal{T}_U$ if and only if $\sS\cap\cB_\Omega\in\mathcal{T}_\Omega^{(\cB)}$ holds for all $\Omega\in\cB_U$.
        \item If a Grothendieck topology $\mathcal{T}$ on $\Ouv_X$ is given, $\sS\in\mathcal{T}_\Omega^{(\cB)}$ if and only if $[\sS]\in\mathcal{T}_\Omega$, (where $[\sS]=[\Theta\sep \Theta\in\sS]$ is the sieve on $\Ouv_\Omega$ generated by $\sS$).
    \end{enumerate}
\end{proposition}
\begin{proof}
    Ommited. Hint : it becomes easier if it is assumed that $\Omega,\Theta\in\cB\Rightarrow \Omega\cap\Theta\in\cB$.
\end{proof}
\begin{definition}
    \label{definition_9_4}
    A $G_+$-topological space is a $T_0$-space $X$ with a Grothendieck topology on $\Ouv_X$ such that
    \begin{enumerate}
        \item $\emptyset\vDash \emptyset$ (implies $\mathcal{G}(\emptyset)=\{0\}$ for sheaves $\mathcal{G}$ of abelian groups).
        \item If $\sS\vDash U$ then $U = \bigcup_{V\in\sS}V$.
    \end{enumerate}
\end{definition}
\begin{definition}
    \label{definition_9_5}
    If $X$ is a $G_+$-space then a $G_+$-base for $X$ is a subset $\cB\subseteq\Ouv_X$ such that $\cB$ is a topology base in the ordinary sense and $[\cB_U]\vDash U$ for all $U\in\Ouv_U$. We say that $\cB$ is $G_{++}$ if in addition,
    \begin{enumerate}[a).]
        \item $\Omega,\Theta\in\cB \Rightarrow \Omega\cap \Theta \in\cB$.
        \item Membership in $\cB$ is local in the sense that $\Omega\in\cB,U\in\Ouv_\Omega$ such that $U\subseteq \Omega$ and $[\Theta\in\cB_\Omega\sep U\cap \Theta\in\cB]\vDash \Omega$ then $U\in\cB$.
    \end{enumerate}
\end{definition}
\begin{definition}
    \label{definition_9_6}
    If $X$ is a $G_+$-space then $V\in\Ouv_X$ is called $G_+$-quasicompact (abbreviated as $(G_{+qc})$) if for all $\sS\in\mathcal{T}_U$, there are finitely many $(V_i)_{i=1,\cdots,n}$ such that 
    \begin{equation*}
        [V_1,\cdots,V_n]\vDash V.
    \end{equation*}
    (i.e. every covering sieve of $U$ has a finitely generated covering subsieve).
\end{definition}
\begin{example}
    If $X$ is a $T_0$-space it gets the structure of a $G_+$-space by
    \begin{equation*}
    \mathcal{T}_U=\left\{\text{sieves $\sS$ in $U$ such that $U = \bigcup_{V\in \sS}V$}\right\}.
    \end{equation*}
\end{example}
\begin{proposition}
    \label{proposition_9_2}
    Let $X$ be a $T_0$-space, $\cB$ a topology base for it such that $\Omega,\Theta\in\cB$ implies $\Omega\cap\Theta\in\cB$. Put 
    \begin{equation*}
        \mathcal{T}_U = \left\{\text{sieves $\sS$ on $U$}\:\bigg{|}\:\forall \Omega\in\cB_U,\exists (\Theta_i)_{i=1}^n\in\sS^n\text{ such that $\Omega=\bigcup_{i=1}^n\Theta_i$ }\right\}
    \end{equation*}
    Then $X$ becomes a $G_+$-topological space and $\cB$ is a $G_+$-base and the elements of $\cB$ are $G_{+qc}$.
\end{proposition}
\begin{proof}
    %The axiom for a Grothendieck topology are easily verified (see sheet 5). 
    GT-Triv follows since for any open set $U$, $U$ itself is a subset of $U$. Let $\sS$ be an admissible $U$-sieve and $\nu:V\hookrightarrow U$. Then as $\mathcal{B}_V\subseteq\mathcal{B}_U$, we have GT-Trans. Let $\mathscr{S},\mathscr{T}$ be $U$-sieves such that $\sS$ is admissible. Then for $\nu:V\hookrightarrow U\in\mathscr{S}$, $\nu^*\mathscr{T}\in\mathcal{T}_V$ means that 
    $\{\Theta_1,\cdots,\Theta_n\}$ as in the construction exists for each element of $\mathcal{B}_V$ in $\mathscr{T}$. Thus $\mathscr{T}\in\mathcal{T}_U$.
    \par To show that it is a $G_+$-topological space. First, $\emptyset\vDash \emptyset$ is clear from the definition. The other parts of $G_+$ for $X$ is also clear. Let $U\in\Ouv_X$ and $\sS=[\cB_U]$. If $\Omega\in\cB_U$ then $\Omega\in\sS$ hence $n=1,\Theta_1=\Omega$ can be taken.
    \par To show that the membership of $\mathcal{B}$ is local, note that for $\Omega\in\cB$ and $U\in \Ouv_X$, 
    \begin{equation*}
        [\Theta\in\cB_\Omega|U\cap\Theta\in\cB]\vDash \Omega
    \end{equation*}
    is generated in $\cB$, thus it implies that 
    \begin{equation*}
        \Omega = \bigcup_{i=1}^n \Theta_i\cap U = \Omega\cap  U = U. 
    \end{equation*}
    Thus $U\in\cB$.
    \par Finally let $\sS\vDash \Omega$ for $\Omega\in\mathcal{B}$. If $\Theta_1,\cdots,\Theta_n$ are as in the formulation of the Proposition. Then $\mathscr{T} = [\Theta_1,\cdots,\Theta_n]$ then $\mathscr{T}\vDash \Omega$ as when $\Xi\in\cB_\Omega$ and $\Xi_i = \Xi\cap \Theta_i\in\cB\cap \mathscr{T}$ by the definition of topology bases and $\Xi=\bigcup_{i=1}^n\Xi_i$. %Therefore the seond part of the definition of $G_+$-space is shown.
\end{proof}
\begin{remark}
    The assertion of the Proposition \ref{proposition_9_2} will still hold when we relax the condition as follows.
    \begin{equation*}
        \Omega,\Theta\in\cB\Rightarrow \exists n\in\N, \exists (\Xi_i)_{i=1,\cdots,n}\in\cB^n\:\st \: \Omega\cap\Theta = \bigcup_{i=1}^n \Xi_i.
    \end{equation*}
\end{remark}
\begin{definition}
    \label{definition_9_7}
    This is called the $G_+$-topology on $X$ enforcing the quasicompactness of the elements of $\cB$.
\end{definition}
\begin{definition}
    \label{definition_9_8}
    A sieve in $\Ouv_X$ is called a prime sieve if $n\in\N$ and $\bigcap_{i=1}^nU_i\in\sS$, (intersection in $X$ and all $U_i\in\Ouv_U$) implies that there is $1\leq i\leq n$ such that $U_i\in\sS$.
\end{definition}
\begin{remark}
    It is sufficient to require this when $n=0,(X\not\in\sS)$ and $n=2,(U\cap V\in\sS\Rightarrow U\in\sS$ or $V\in\sS$).
\end{remark}
\begin{proposition}
    \label{proposition_9_3}
    Let $X$ be a $G_+$-space. For a subset $\xi\subseteq\Ouv_X$ the following conditions are equivalent.
    \begin{enumerate}[a).]
        \item $\xi$ satisfies 
            \begin{enumerate}[i).]
                \item $U\in\xi,V\in\Ouv_X,U\subseteq V\Rightarrow V\in\xi$.
                \item $U,V\in\xi\Rightarrow U\cap V\in\xi$ and moreover $X\in\xi$.
                \item If $U\in\xi$ and $\sS\in\mathcal{T}_U$ then $\sS\cap\xi\not=\emptyset$.
            \end{enumerate}
        \item $\mathscr{T}=\Ouv_X\backslash\xi$ is a prime sieve containing every open subset covered by it ($U\in\Ouv_X$ and $\mathscr{T}//U$ implies $U\in\mathscr{T}$).
    \end{enumerate}
\end{proposition}
\begin{proof}
    The primeness follows from i) and ii). 
\end{proof}
\begin{definition}
    \label{definition_9_9}
    We call $\xi$ a van der Put point.
\end{definition}
% \begin{remark}
%     The set of topos points (refering to SGA 4.1). If the category of sheaves of sets on $X$ (as defined later on) is in canonical bijection between van der Put points.
% \end{remark}
\begin{definition}
    \label{definition_9_10}
    The set of vander Put points of $X$ is denoted by $X^*$.
\end{definition}
\begin{proposition}
    \label{fact_9_2}
    Let $X$ be a $G_+$-space and $U\in\Ouv_U$ be an open subset. Equip $U$ with the $G_+$-structure $\mathcal{T}^{(U)}=\mathcal{T}|_{\Ouv_U}$, we then have a bijection
    \begin{center}
        \begin{tikzcd}
\{\text{van der Put poitns $\xi_U$ of $U$}\} \arrow[r, "\xi_U\mapsto \xi=\{V\in\Ouv_X\sep V\cap U\in\xi_U\}", shift left=2] & U^*=\{\xi\in X^*\sep U\in\xi\} \arrow[l, "\xi_U=\xi\cap\Ouv_X\mapsfrom \xi", shift left=2]
\end{tikzcd}
    \end{center}
    We have $(U\cap V)^*=U^*\cap V^*$.
\end{proposition}
\begin{proof}
    $\xi_U\mapsto\xi$ is justified since $\cdot\cap U$ preserves inclusion and $\cap$ is associative.
\end{proof}
%current Current
\begin{definition}
    \label{definition_9_11}
    We equip $X^*$ with the topology for which $\{U^*\sep U\in\Ouv_X\}$ is a topology base.
\end{definition}
\begin{proposition}$\:$
    \label{fact_9_3}
    \begin{enumerate}[a).]
        \item If $\sS\vDash U$ then $U^*=\bigcup_{V\in\sS}V^*$. (This follows from Proposition \ref{proposition_9_3} a) iii)).
        \item If $\cB$ is a $G_+$ base then $\cB^*=\{\Omega^*\sep \Omega\in\cB\}$ is a topology base for $X^*$. (Follows from Definition \ref{definition_9_5}).
        \item $\xi\in X^*$ is a specialization (i.e. $\xi\in\overline{\{u\}}$ of $u\in X^*$ ) if and only if $\xi\subseteq u$.
        \item If $\sS$ is a maximal element of the set of sieves on $X$ which are not in $\mathcal{T}_X$. Then $\xi=\Ouv_U\backslash\sS$ is a closed point of $X^*$.
    \end{enumerate}
\end{proposition}

\begin{proof}
    We show the proof of d). Obviously, $X\not\in\sS$ by GT-Triv. Let $U_1,U_2\not\in\sS$ and set $U=U_1\cap U_2 \in\sS$. By maximality of $\sS$, $\sS_i=\sS\cup\Ouv_{U_i}\in\mathcal{T}_X$. But $\sS|_V=\sS_2|_V$ for all $V\in\sS_1$.
    Then $\sS\in\mathcal{T}_X$ by GT-Loc, a contradiction. Thus $\sS$ is a prime sieve. if $\sS//V$ then $\sS\cup\Ouv_X\not\in\mathcal{T}_X$ (otherwise $\sS\vDash X$ by GT-Loc). Then $V\in\sS$ by maximality of $\sS$. By Proposition \ref{proposition_9_3}, $\xi=\Ouv_X\backslash\sS$ is a van der Put point. If $\xi$ has a proper specialization $u$ then $\mathscr{T}=\Ouv_X\backslash u\not\in\mathcal{T}_X$ and $\mathscr{T}\supsetneq\sS$ a contradition..
\end{proof}

% 5/22

\begin{definition}
    \label{definition_10_10}
    Let $X$ be a $G_+$-space. We say $X$ has sufficiently many van der Put points if for all $U\in\Ouv_X$, and every $U$-sieve $\sS$, $U^*=\bigcup_{V\in\sS}V^*$ implies $\sS\vDash U$. 
\end{definition}

\begin{remark}
    Definition \ref{definition_10_10} can be restated by saying that the converse of a) in Proposition \ref{fact_9_3} also holds.
\end{remark}

\begin{lemma}
    If $X$ has sufficently many van der Put points, then it is $G_+$-quasicompact and if $\sS$ is an $\Ouv_X$-sieve such that $X^*=\bigcup_{U\in\sS}U^*$ then $\sS = X$.
\end{lemma}

\begin{proof}
    Assume that $\sS$ is not covering $X$. Set,
    \begin{equation*}
    \mathcal{M} \defeq\{\mathscr{T}\sep \sS\subseteq\mathscr{T},\mathscr{T}\not\in\mathcal{T}_X\}.
    \end{equation*}
    We apply Zorn's lemma to $(\mathcal{M},\subseteq)$. This is possible as $\mathcal{M}\not=\emptyset$ because $\sS\in\mathcal{M}$ and if $\mathcal{L}\subseteq\mathcal{M}$ is non-empty subset linearly ordered by $\subseteq$, then 
    \begin{equation*}
    \widetilde{\sS}=\bigcup_{\mathscr{T}\in\mathcal{L}}\mathscr{T}
    \end{equation*}
    is a sieve containinig $\sS$. If $\widetilde{\sS}$ $\in\mathcal{T}_X$, then by the quasi-compactness of $X$, there are $U_1,\cdots,U_n\in\widetilde{\sS}$ such that $[U_1,\cdots,U_n]\vDash X$. Then there is $\mathscr{T}\in\mathcal{L}$ such that all $i=1,\cdots,n,U_i\in\mathscr{T}$ therefore 
    \begin{equation*}
    [\Theta_1,\cdots,\Theta_n]\vDash x,[\Theta_1,\cdots,\Theta_n]\subseteq\mathscr{T}\stackrel{\text{Proposition \ref{fact_8_1}}}{\Rightarrow}\mathscr{T}\vDash X.
    \end{equation*}
    This is a contradiction to $\mathscr{T}\in\mathcal{M}$.
    \par It follows that $\mathcal{M}$ has a $\subseteq$-maximal element. By Proposition \ref{fact_9_3}, $\xi=(\Ouv_X\backslash\sS_{\max})\in X^*$. If $U\in\sS\subseteq\sS_{\max}$, then $\xi\not\in U^*$ (otherwise $U\in\xi$, but by Proposition \ref{fact_9_2}, $\xi\cap\sS_{\max}\not=\emptyset$). A contradiction to our assumption that $X^*=\bigcup_{U\in\sS}U^*$.
\end{proof}

\begin{proposition}
    If $X$ is a $G_+$-space which has a $G_+$-base $\cB$ where elemnts are $G_+$-quasicompact. Then $X$ has sufficiently many van der Put points. 
    If $\cB$ is closed under finite intersection in $X$ (including intersection over an empty set thus $X\in\cB$), then $X^*$ is a spectral space.
\end{proposition}

\begin{proof}
    Let $U\in\Ouv_X$ and $\sS$ any $U$-sieve such that $U^* = \bigcup_{V\in\sS}V^*$. Let $\Omega\in\cB_U$ and $\sS_\Omega = \sS|_\Omega$. Then $\Omega^*\subseteq U^*$ hence $\Omega^*=\bigcup_{V\in\sS}\Omega^*\cap V^*=\bigcup_{V\in\sS}(\Omega\cap V)^*$. (Note that $(V\cap W)^* = V^*\cap W^*$ by definition).
    Since $\Omega\cap V\in\sS_\Omega$, it follows that $\Omega^* = \bigcup_{V\in\sS_\Omega}V^*$. By Lemma \ref{lemma_1}, $\sS\vDash \Omega$. Let $\mathcal{T}$ be the sieve generated by $\cB_U$. From GT-Trans and what we have seen $\sS//W$ where $W\in\mathcal{T}$. By GT-Loc, $\sS\vDash U$ hence $X$ has sufficiently many van der Put points.
    \par For the spectrality of $X^*$, it is easy to see that $\Omega^*$ is quasicompact if $\Omega$ is $\Omega$ is $G_+$-quasicompact. Then $\{\Omega^*\sep \Omega\in\cB\}$ satisfy the requirement for the quasicompactness part of the definition of spectrality. If $Z\subseteq X^*$ is a van der Put point, then $\xi=\{U\in\Ouv_X\sep U^*\cap Z\not=\emptyset\}$ is a generic point of $Z$. $T_0$ is also easy to show.
\end{proof}

\begin{remark}
    \begin{equation*}
        \glim_{X^*}\U = \{U\in\Ouv_X\sep \text{$U$ contains some $\Omega\in\cB$ such that $\Omega^*\in\mathcal{F}$}\}
    \end{equation*}
    All $\Omega^*$ with $\Omega\in\cB$ are closed under this. 
\end{remark}

\subsection{Sheaves}
\begin{definition}
    \label{definition_10_1}
    A presheaf on $X$ with values in (sets, abelian groups, rings) is a contravariant functor $\mathcal{G}$ from $\Ouv_X$ to one of these categories. Then all $\mathcal{G}(U)$ are objects in these categories and we have
    \begin{equation*}
    V\subseteq U\Rightarrow \mathcal{G}(U)\stackrel{g\mapsto g|_V}{\to}\mathcal{G}(V),
    \end{equation*}
    which is a morphism in the target category such that $g\in\mathcal{G}$ then $g|_U = g,(g|_V)_W=g|_W$ if $W\subseteq V\subseteq U$.
    If this is only given on elements of a certain topology base $\cB$, we speak of a presheaf on $\cB$. In all cases, if a Grothendieck topology in $\cB$ is given, then a sheaf is a presheaf $\mathcal{G}$ satisfying the additional assumption that 
    \begin{equation*}
        \mathcal{G}(\Omega)\stackrel{\sim}{\to}\varprojlim_{\Theta\in\sS}\mathcal{G} = \left\{(g_\Theta)_{\Theta\in\sS}\in\prod_{\Theta\in\sS}\mathcal{G}(\Theta)\sep g_\Theta|_{\Xi} = g_\Xi\text{ where $\Theta\in\sS, \Xi\in\cB_\Theta$}\right\}.
    \end{equation*}
    If the map is only injective, then 
\end{definition}

\begin{definition}$\:$
    \label{definition_10_2}
    \begin{enumerate}
        \item A stalk of $\mathcal{G}$ at $\xi\in X^*$ is such that
        \begin{equation*}
            \mathcal{G}_\xi = \varinjlim_{\Omega\in\xi}\mathcal{G}(\Omega) = \{(\Omega,g)\sep \Omega\in\xi\cap\cB,g\in\mathcal{G}(\Omega)\}/\sim
        \end{equation*}
        where the equivalent relation is given by $(g,\Omega)\sim (r,\Theta)$ if there is $\Xi\subseteq\Omega\cap \Theta$ such that $\Xi\in\cB$ and $g|_{\Xi} = g|_{\Xi}$.
        \item A $G$-skyscraper sheaf of $\xi\in X^*$ is \begin{equation*}
            ([G]_\xi)(U) = \begin{cases}
                G,\quad U\in\xi,\\
                \{0\}
            \end{cases}
        \end{equation*}
    %     \item When $X$ is a $G_+$-space with sufficiently many van der Put points %need revision
    %     \begin{equation*}
    %         (\mathrm{Gerb}(\mathcal{G}))(U) = \left\{r\in\prod_{\xi\in U^*}\mathcal{G}_\xi\:\bigg{|}\:\sS_r \vDash U\right\}
    %     \end{equation*}
    %     where $\sS_\gamma$ is the sieve of or in $\Ouv_X$ generated by as in Proposition \ref{proposition_9_1} those $\Omega\in\cB_u$ such that there is $g\in\mathcal{G}(\Omega)$ such that $\gamma_\xi$ the component of $\xi$ of $\gamma$) is the image of $g$ under $\mathcal{G}(\Omega)\to\mathcal{G}_\xi$, for all $\xi\in\Omega^*$.
    \end{enumerate}
\end{definition}

\begin{definition}
    Let $X$ be a $G_+$-space with sufficiently many van der Put points and $\mathcal{G}$ be a sheaf on $X$. For an open subset $U\subseteq X$ and
    \begin{equation*}
        r\in \prod_{\xi\in U^*}\mathcal{G}_\xi.
    \end{equation*}
    We denote the sieve $\mathscr{S}_r$ on $U$ generated in $\Ouv_X$ to be such that 
    \begin{equation*}
        \mathscr{S}_r \defeq \left[\Omega\in\cB_U\:\bigg{\vert}\:\exists g\in\mathcal{G}(\Omega) \st\forall \xi\in \Omega^*, r_\xi = g_\xi\right].
    \end{equation*}
    Then we define,
    \begin{equation*}
        (\mathrm{Gerb}(\mathcal{G}))(U) = \left\{r\in\prod_{\xi\in U^*}\mathcal{G}_\xi\:\bigg{\vert}\:\sS_r\vDash U\right\}.
    \end{equation*}
    For the well-definedness, see Proposition \ref{proposition_9_1}.
\end{definition}


\begin{proposition}
    \label{proposition_10_1}
    Let $X$ be a $G_+$-space with sufficiently many van der Put points.
    \begin{enumerate}[a)]
        \item The stalk is exact on presheaves.
        \item $[G]_\xi$ is a sheaf and $\mathcal{G}\mapsto\mathcal{G}_\xi$ is left adjoint to $G\mapsto [G]_\xi$ on (pre)sheaves.
        \item If $\mathcal{G}$ is a presheaf, the morphism $\mathcal{G}\to\Sheaf(\mathcal{G})$ sending $g\in\mathcal{G}(U)$ to $($image of $g$ in $\mathcal{G}_\xi)_{\xi\in\U^*}\in\prod_{\xi\in U^*}\mathcal{G}_\xi$ defines an isomorphisms on stalks. 
        \item $(\mathcal{G}_\cB)_\xi\stackrel{\simeq}{\to}\mathcal{G}_\xi$ for all $x\in\in X$. 
        \item A morphism of presheaves defining isomorphism on stalks defines an isomorphism on sheafification. 
        \item If $\mathcal{G}$ is a sheaf, $\mathcal{G}\to\Sheaf(\mathcal{G})$ is an isomorphism. If $\mathcal{G}$ is only defined on $\cB$, $\mathcal{G}\stackrel{\simeq}{\to}\Sheaf(\mathcal{G})|_\cB$.
        \item A morphism of sheaves defining isomorphism on all stalks is an isomorphism.
        \item Sheafification $(\text{Presheaves on $\cB$})\to(\text{sheaves on $\Ouv_X$})$ is left adjoint to the functor of restriction to $\cB$.
        \item One has an equivalence of category between sheaves on $X$ (for the Grothendieck topology on $X$) and (ordinary) sheaves on $X^*$, sending a sheaf $\mathcal{G}$ on $X$ to the sheafification of the (pre)sheaf $U^*\to \mathcal{G}(U)$ on $\{U^*\sep U\in\Ouv_X\}$ ( a topology base for $X^*$) and a sheaf $\mathcal{H}$ on $X^*$ to $\mathcal{G}(U)\to\mathcal{H}(U^*)$. 
    \end{enumerate}
\end{proposition}

\begin{proof}
    a)-h) are showwn in some ways as in derived sheaf theory. i) follows from these 
\end{proof}

%6/1

\section{General non-archimedian topological ring}
\subsection{General topological rings and modules}
\begin{definition}
    \label{definition_11_1_1}
    A topological group is a group $G$ with a topology on the underlying set such that multiplication and inversion are continuous. That is 
    \begin{equation*}
        \mu(g,h)=gh,\iota(g) = g^{-1}
    \end{equation*}
    are continuous.
\end{definition}
\begin{definition}
    \label{definition_11_1_2}
    A topological ring is a ring $R$ with a topology on the underlying set such that it is a topological group under addition and 
    \begin{equation*}
        R\times R\ni(a,b)\mapsto ab\in R
    \end{equation*}
    is continuous. (We do not assume the continuity of $a\mapsto a^{-1}$ for $a\in R^\times$).
\end{definition}
\begin{definition}
    \label{definition_11_1_3}
    Let $R$ be a topological ring. A topological $R$-module $M$ is an $R$-module with topology on its underlying set such that $(M,+)$ is a topological group and 
    \begin{equation*}
        R\times M\ni (a,m)\mapsto am\in M
    \end{equation*}
    is continuous.
\end{definition}
\begin{proposition}
    \label{fact_11_1}
    For a topological group $G$, the following conditions are equivalent.
    \begin{enumerate}[a).]
        \item $G$ is $T_0$ (Komogolov).
        \item $G$ is $T_1$ (Fréchet).
        \item $G$ is $T_2$ (Hausdorff).
        \item $\{0\}$ is closed in $G$.
    \end{enumerate}
\end{proposition}
\begin{proof}$\:$
    \par If $x\in\overline{\{0\}}$ then $x^{-1}\in\overline{\{0\}}$ by the continuity of inversion. Hence $1\in\overline{\{x\}}$ by continuity of $\mu(x,\cdot)$ hence $T_0$ is satisfied. For a)$\Rightarrow $d), let 
    \begin{equation*}
        \Delta = (\text{preimage of $\{u\}$ under $G\times G\stackrel{(g,h)\mapsto g^{-1}h}{\to} G$})
    \end{equation*}
    is closed thus $T_2$.
\end{proof}
\begin{proposition}
    \label{fact_11_2}
    Every open subgroup $H$ of a topological group $G$ is also closed, and every group $H\subseteq H'\subseteq G$ is also clopen.
\end{proposition}
\begin{proof}
We have,
\begin{equation*}
    G\backslash H = \bigcup_{g\in G/H\backslash\{0\}}gH
\end{equation*}
is open therefore $H$ is closed. Every $H'$ subset is a union of right or left $H$-cosets in $G$ which is also open.
\end{proof}
\begin{definition}
    \label{definition_11_2}
    Let $R$ be a topological ring and $M$ be a topological $R$-module. A subset $X\subseteq M$ is bounded or $R$-bounded if for every neighborhood $U$ of $0$ in $M$, there is a neighborhood $V$ of $0$ in $R$ such that 
    \begin{equation*}
        V\cdot X = \{vx \sep v\in V,x\in X\}\subseteq U.
    \end{equation*}
\end{definition}
\begin{example}
    When $R$ has the discrete topology, every subset is bounded (take $V=\{0\}$).
\end{example}
\begin{example}
    This defines a boundary on $X$. If $X_1,X_2$ are bounded then so is $X_1\cup X_2$. Every $\{m\}$ is bounded, $M\subseteq N$ and $N$ bounded then $M$ is also bounded.
\end{example}
\begin{example}
    If $R=K$, equipped with an valuation $\ab{\cdot}$ and $M$ is an (valuative) $K$-vector space and $K$ is not discrete, then a subset $X$ of $M$ is bounded if and only if there is $B\in[0,\infty)_\R$ such that 
    \begin{equation*}
        \forall x\in X, \norm{x}\leq B.
    \end{equation*}
    For example, take $U = \{x\in X\sep \norm{x}\leq 1\}, B = \ab{x}^{-1}$ where $x\in V\backslash\{0\}$ and $V$ is as in the definition.
\end{example}
\begin{definition}
    \label{definition_11_3}
    For subsets $X_1,X_2\subseteq R$, we define,
    \begin{equation*}
        X_1X_2 = \{x_1x_2|x_1\in X_1,x_2\in X_2\}.
    \end{equation*}
    Let $X^0 = \{1\}$. We define recursively $X^n$ such that 
    \begin{equation*}
        X^{n+1} = X\cdot X^n.
    \end{equation*}
    We say that. a subset $X\subseteq R$ is power bounded if 
    \begin{equation*}
        \bigcup_{n\in\N}X^n
    \end{equation*}
    is bounded. It is topologically nilpotent if for every neighborhood $U$ of $0$ in $R$, there is $n\in\N$ such that 
    \begin{equation*}
        \bigcup_{m=n}^\infty X^m \subseteq U.
    \end{equation*}
    We call an element $r\in R$ topologically nilpotent / power bounded if $\{r\}$ satisfies these properties.
\end{definition}
\begin{proposition}
    \label{fact_11_3}
    Let $R$ be a topological ring and $M$ be a topological $R$-module. We have the following statements.
    \begin{enumerate}[a).]
        \item A finite union of bounded subsets of $M$ is bounded. 
        \item Trivially, every finite subset of $M$ is bounded.
        \item If $X_1,X_2\subseteq M$ are bounded then so is $X_1+X_2 = \{x_1+x_2\sep x_1\in X_1,x_2\in X_2\}$. 
        \item If $X\subseteq R$ and $Y\subseteq M$ are bounded then so is $XY\subseteq M$. 
    \end{enumerate}
\end{proposition}
\begin{proof} For c), let $U$ be a neighborhood of $0$ in $M$, there are a neighborhood $\tilde{U}$ of $0$ in $M$ such that $\tilde{U}+\tilde{U}\subseteq U$ and another neighborhoods $V_1,V_2$ of $0$ in $R$ such that $V_1X_1,V_2X_2\subseteq \tilde{U}$. Take $V = V_1\cap V_2$ in Definition \ref{definition_11_2}. Proofs for the other statements are also easy.
\end{proof}
\begin{proposition}
    \label{fact_3_another}
    If $X\subseteq M$ is $R$-bounded and $M\stackrel{f}{\to}N$ a continuous morphism of $R$-modules then $f(X)$ is bounded in $N$.
\end{proposition}
\begin{proof}
    If $U\subseteq N$ is a neighborhood of $0$, let $\tilde{U} = f^{-1}(U)$ and $V$ a neighborhood of $0$ in $R$ such that $VX\subseteq\tilde{U}$. Then $Vf(X)\subseteq U$.
\end{proof}
It is not always true that the image of a bounded subset of a $R$ under. a morphism $R\to S$ of topological ring is bounded. (The image is $R$-bounded but not necessarily $S$-bounded) for example take $R=\Q_p,S=\Q_q$. However, the following conditions on $R$ asserts that this does not happen.
\begin{align*}
    &\text{$R$ has a topologically nilpotent unit}\tag{T}\label{T}\\
    &\text{$R$ has a sequence of units converging to $0$}\tag{H}\label{H}\\
    &\text{Every neighborhood of $0$ intersects $R^\times$}\tag{E}\label{E}
\end{align*}
\begin{proposition}
    \label{fact_11_4}
    Let $R\stackrel{f}{\to}S$ be a continuous morphism of topological ring. If $R$ has \eqref{T},\eqref{H},\eqref{E}, then so does $S$ and the image $f(X)$ of a bounded subset of $R$ is a bounded subset of $S$.
\end{proposition}
\begin{proof}
    If $U$ is a neighbourhood o $0$ in $S$ and $R$ has \eqref{E}, put $\tilde{U}=f^{-1}(U)$ and $r\in R^\times\cap \tilde{U}$. $f(r)\in S^\times \cap U$. If 
    \begin{equation*}
        \lim_{j\to\infty} r_j = 0
    \end{equation*}
    in $R$ and all $r_j\in R^\times$,%(where $r_j = r^j$ and $r\in R^\times$)
    \begin{equation*}
        \lim_{j\to\infty}f(r_j) = 0
    \end{equation*}
    in $S$ and all $f(r_j)\in S^\times$. 
    %(respectively $f(r_j) = s^j$ with $s = f(r)\in S^\times$) 
    showing that \eqref{H}. For \eqref{T}, take a topologically nilpotent unit $r\in R$, then by setting $r_j = r^j$ in the previous argument, we see \eqref{T} is inheritted.
    \par If $R$ has \eqref{E} and $X\subseteq R$ bounded and $U$ is a neighborhood of $0$ in $S$, there is $\tilde{U}$ a neighborhood $0$ in $S$ such that $\tilde{U}\tilde{U}\subseteq U$. Let $U' = f^{-1}(\tilde{U})\subseteq R$ a neighborhood of $0$. There is a neighbourhood $V'$ of $0$ in $R$ such that $V'X\subseteq U'$. As $R$ has \eqref{E}, we can take $\rho\in R^\times\cap f^{-1}(\tilde{U})$. Take $V = \rho V'$ then $Vf(X)\subseteq U$.  
\end{proof}
\begin{proposition}
    \label{fact_11_5}
    Let $R$ be a topological ring. 
    \begin{enumerate}[a).]
        \item Every power bounded subset is bounded.
        \item A finite union of power bounded subsets is power bounded.
        \item A subset $X\subseteq R$ which is both bounded and topologically nilpotent is power bounded.
        \item A finite union of subsets which are both bounded and topologically nilpotent is topologically nilpotent.
        \item If $X$ is power bounded or topologically nilpotent and $Y$ is topologically nilpotent, $XY$ is then topologically nilpotent.
        \item The image of a topologically nilpotent $X\subseteq R$ under a morphism $R\to S$ of topological rings is topologically nilpotent.
        \item If $R$ has \eqref{E}, $R\stackrel{f}{\to}S$ is a morphism of topological rings and $X\subseteq R$ is bounded in $R$. Then $f(X)$ is bounded in $S$.
    \end{enumerate}
\end{proposition}
\begin{proof}
    Mostly straightforward.
\end{proof}
\begin{definition}
    \label{definition_11_4}
    Let $R$ be a topological ring, we define the following subsets.
    \begin{equation*}
        R^\circ = \{r\in R\sep \text{$r$ is powewr bounded}\},R^{\circ\circ}=\{r\in R\sep \text{$r$ is topologically nilpotent}\}.
    \end{equation*}
\end{definition}
\begin{remark}
    Without further conditions, these may fail to be additionally closed. By Proposition \ref{fact_11_3} b), and Proposition \ref{fact_11_5} a), $R^\dcirc\subseteq R^\circ$ always holds. %Needs revision
\end{remark}
\begin{definition}
    \label{definition_11_5_1}
    Let $(m_n)_{n\in\N}$ be a sequence in $M$. It is Cauchy if for every neighbourhood $U$ of $0$ in $M$, there is $k$ such that $m_i-m_j\in U$ where $\min\{i,j\}\geq k$. $M$ is sequentially complete if every Cauchy sequence has a limit in $M$.
\end{definition}
\begin{definition}
    \label{definition_11_5_2}
    A filter $\mathcal{F}$ of $M$ is Cauchy if every neighborhood $U$ of $0$ in $M$, there is $X\in\mathcal{F}$ such that $x-y\in U$ where $x,y\in X$. $M$ is complete if every Cauchy filter has a limit.
\end{definition}
\begin{definition}
    \label{definition_11_5_3}
    The sequential completion of $M$ is a set of equivalence classes of Cauchy sequences in $M$ where two Cauchy sequences $(m_i),(n_i)$ are equivalent if for every neighborhood $U$ of $M$, there is $k\geq0$ such that $i\geq k$ then $m_i-n_i\in U$.
\end{definition}
\begin{remark}
    If $G$ is a non-commutative topological group, then it typically requires that both $g_ig_j^{-1}\in U$ and $g_i^{-1}g_j\in U$ where $\min\{i,j\}\geq k$. Respectively, $x^{-1}y\in U$ and $xy^{-1}\in U$ for $x,y\in X$ for filters.
\end{remark}
\begin{remark}
    One can complete by taking equivalence classes of Cauchy filters. But how to continus multiplication in that completion of $R^2$. (Sequential completion does not have this problem).
\end{remark}
\begin{remark}
    It is known that topological vector space over $\C$ or $\R$ is the quotient of a complete topological vector space $V$ over that field such that every bounded subset of $V$ is contained in a finite dimensional subspace.
\end{remark}
\begin{proposition}
    \label{fact_11_7}
    Every quotient of a complete topological group with a countable neighborhood base of $1$ by a closed normal divsor is complete.
\end{proposition}
\subsection{Continuous Valuations}
For an ordered group $\Gamma$ equip $\Gamma\cup\{0\}$ with the topology where $U$ is open if and only if $0\not\in U$ or there is $\gamma\in\Gamma$ such that $(0,\gamma)_\Gamma\subseteq U$.
% 6/5
\par Consider a valuation $A\stackrel{v}{\to}\Gamma$ where $A$ is a topological ring. Recall that $\Gamma_v$ is the convex subgroup of $\Gamma$ such that
\begin{equation*}
    \Gamma_v = \langle\{v(a)\sep a\in A\cdot\supp v\}\rangle_\Gamma.
\end{equation*}
Also recall that for $a\in A$ and $\gamma\in \Gamma$,
\begin{equation*}
    K_{v<a} = \{\alpha\in A\sep v(\alpha)<v(a)\},K_{v<\gamma} = \{\alpha\in A\sep v(\alpha)<\gamma\}.
\end{equation*}
Similarly,
\begin{equation*}
    K_{v\leq a} = \{\alpha\in A\sep v(\alpha)\leq v(a)\},K_{v\leq\gamma} = \{\alpha\in A\sep v(\alpha)\leq\gamma\}.
\end{equation*}
They come with an exception such that 
\begin{equation*}
    K_{v<a}= \emptyset,
\end{equation*}
when $a\in\supp v$. They are subgroups of $(A,+)$. Note that a subgroup $H\subseteq G$ of a topological group $G$ is open if and only if it is a neighbourhood of $0$. Also recall that subgroups containing open subgroups are also open.
%current Current 
\begin{proposition}
    \label{fact_12_1}
    The following conditions are equivalent.
    \begin{enumerate}[a).]
        \item $A\stackrel{v}{\to}\Gamma_v\cup\{0\}$ is continuous.
        \item For all element $\gamma$ of the convex hull of $\Gamma_v$ in $\Gamma$, $K_{v<\gamma}$ is open.
        \item For all $\gamma\in\Gamma_v$, $K_{r<\gamma}$ is open.
        \item For all $a\in A\backslash\supp(v)$, $K_{r<a}$ is open.
        \item For all $a\in A\backslash\supp(v)$, $K_{r<a}$ is a neighbourhood of $0$.
    \end{enumerate}
\end{proposition}
\begin{proof}
    By the previous argument d) and e) are equivalent. d)$\Rightarrow$c) is trivial. For c)$\Rightarrow$a), consider $\gamma\in\Gamma_v\cup\{0\}$ and a neighbourhood $U$ of $\gamma$ in $\Gamma_v\cup\{0\}$. We must show that $v^{-1}(U)$ is a neighbourhood of $a$ in $A$. If $v(a)\not=0$, it is sufficient to consider 
    \begin{equation*}
        U = \{v(a)\}.
    \end{equation*}
    Then,
    \begin{equation*}
        v(b) = v(a)
    \end{equation*}
    where $b\in K_{r<a}+a$, and the set of such $b$ is a neighborhood of $a$ in $A$. If $v(a)=0$, we may assume $U=[0,\gamma)_{\Gamma_v\cup\{0\}}$ with $\gamma\in\Gamma_v$. Then $v^{-1}(U)=K_{v<\gamma}$ is open by the assumption c).
    \par a)$\Rightarrow$d) is trivial by the definition of the topology on $\Gamma_v\cup\{0\}$, $K_{v<a}$ being $v^{-1}([0,v(a))_{\Gamma_v\cup\{0\}})$ where 
    \begin{equation*}
        [0,v(a))_{\Gamma_v\cup\{0\}}
    \end{equation*}
    is a neighbourhood of $0$ in $\Gamma_v\cup\{0\}$. 
    \par d)$\Rightarrow$b), every $\gamma\in\Gamma_v$ has the form,
    \begin{equation*}
        \gamma  = {\frac {v(a)} {v(b)}},
    \end{equation*}
    with $a,b\in A\backslash\supp v$. Then $K_{v<\gamma}$ is the inverse image under the endomorphism $A\stackrel{\cdot b}{\to}A$ of the topological group $(A,+)$ of the open subgroup $K_{r<a}$. Hence it is also an open subgroup.
\end{proof}
\begin{definition}
    \label{definition_12_1}
    A valuation of $A$ is called continuous if the equivalent conditions in Proposition \ref{fact_12_1} hold.
\end{definition}
\begin{remark}$\:$
    \begin{enumerate}
        \item  This only depends on the equivalent class of $v$.
        \item $\Gamma_v\cup\{0\}$ carries the topology induced from $\Gamma\cup\{0\}$ only if $\Gamma$ is the convex hull of $\Gamma_v$ in $\Gamma$. Hence $\Gamma$ can be replaced by $\Gamma_v$ in c) only if $\Gamma$ is equal to that convex hull (in which case the c) is b)).
    \end{enumerate}
\end{remark}
\begin{example}$\:$
    \label{example_12_1}
    \begin{enumerate}
        \item $\Q_p\stackrel{\ab{\cdot}_p}{\to}[0,\infty)_\R$ is a continuous valuation.
        \item The trvial valuation for $\p\in\Spec A$ is continuous if and only if $\p$ is open.
        \item If $\p$ is an open prime ideal, then any valuation with support $\p$ is open.
        \item For $r\in(0,1]_\R$, the valuation $N_r$ from 1.1(1) (need to look for the reference) is continuous. %later
    \end{enumerate}
\end{example}
\begin{proposition}
    \label{fact_12_2}
    If $v$ is a continuous valuation. Then $\supp v$ is closed.
\end{proposition}
\begin{proof}
    Indeed $\supp v = v^{-1}(\{0\})$ and $\{0\}$ is closed in $\Gamma_v\cup\{0\}$. 
\end{proof}
\begin{definition}
    \label{definition_12_2}
    We denote the subset equivalence classes of continuous valuations of $A$ to be $\Cont(A)\subseteq\Spv(A)$.
    Let $A^+\subseteq A$ be a subset we then denote $\Spa(A,A^+)$ to be the set of equivalence classes of continuous valuations $v$ of $A$ such that $a\in A^+\Rightarrow v(a)\leq 1$.
\end{definition}
\begin{remark}$\:$
    \begin{enumerate}
        \item At this point, we do not yet introduce a topology on these sets.
        \item If $A=\T_1$ as in 1.1.(1)(look at the first part and take the reference) and $A^+=A^\circ$. $N_r\in\Spa(A,A^+)$ if and only if $r\in (0,1)_\R$.
    \end{enumerate}
\end{remark}
\begin{definition}
    \label{definition_12_3}
    A Oka system is a ring $A$ and a set $\mathfrak{M}$ of ideals in $A$ such that for any ideal $I$ in $A$, $a\in A$ and the ideal $I+aA$ and 
    \begin{equation*}
        (I:a) \defeq \{\alpha\in A\sep a\alpha\in I\}
    \end{equation*}
    both belong to $\mathfrak{M}$ then $I\in\mathfrak{M}$.
\end{definition}
\begin{example}
    \label{example_12_2}
    The set of finitely generated ideals in $A$ form a Oka system. This observation gives Eakin-Nagata or Cohen's theorem which state that a ring is Noetherian if and only if all of its prime ideals are finitely generated.%in Oka topology in Proposition \ref{proposition_1} as this Oka system gives Eitoin-Nagata?. If all prime ideals of $A$ are finitely generated, $A$ is Noetherian.
\end{example}
\begin{proposition}
    \label{proposition_12_1}
    Let $\mathfrak{M}$ be an Oka system and $\p\not=A$ be a $\subseteq$-maximal element of the set of ideals of $A$ not belonging to $\mathfrak{M}$. Then $\p\in\Spec A$. 
\end{proposition}
\begin{proof}
    Assume $ab\in\p$. Let $a\not\in \p,b\not\in\p$. Then $\p+aA$ and $(\p:a)$
    contain $\p$ and not equal to $\p$ as $b\in(\p:a)$. By the maximality, both belong to $\mathfrak{M}$ thus $\p\in\mathfrak{M}$ which is a contradiction.
\end{proof}
\begin{proposition}
    \label{fact_12_3}
    Let $\mathfrak{M}$ be a set of ideals such that $I\subseteq J$ and $I\in\mathfrak{M}$ implies $J\in\mathfrak{M}$. Suppose further that $\mathfrak{M}$ is closed under products, that is $I,J\in\mathfrak{M}$ then $I\cdot J\in\mathfrak{M}$. Then $\mathfrak{M}$ is an Oka system.
\end{proposition}
\begin{proof}
    Indeed $I\supseteq (I+aA)(I:a)$ holds. Therefore by definition, it is Oka.
\end{proof}
\begin{remark}
    The Oka system from Example \ref{example_12_2} satisfies the assumptions if and only if $A$ is Noetherian.
\end{remark}
\begin{example}
    \label{example_12_3}
    The set $\mathfrak{M}_S$ of ideals not disjoint form a multiplicative subset $S\subseteq A$ is an Oka system. Then every $\subseteq$-maximal non-element of $\mathfrak{M}_S$ is a prime ideal. (This could also be seen by assuming that all maximal ideals of $A_S$ are prime). By Zorn's lemma, every ideal not belonging to $\mathfrak{M}_S$ is contained in an ideal which is maximal with respect to this non-containment. %is equal to the set of all ideals (i.e. $0\not\in S$)%there is a prime ideal $\p\in\mathfrak{M}_S\cap\Spec A$.
    By Proposition \ref{proposition_12_1}, such ideal is prime.
\end{example}
\begin{proposition}
    \label{fact_12_4}
    Let $v$ be a valuation of a topological ring $A$ satisfying the following conditions.
    \begin{enumerate}
    \item $K_{v\leq 1}$ is a neighbourhood of $0$ in $A$.
    \item For each $a\in A\backslash\supp v$, there is $b\in A$ with $v(ab)\geq1$.
    \end{enumerate}
    % $K_{r\leq 1}$ is a neighbourhood of $0$ in $A$ and such that for each $a\in A\backslash\supp v$, there is $b\in A$ such that $v(ab)>1$. Then $v$ is continuous.
    Then $v$ is continuous.
\end{proposition}
\begin{proof} 
    If $a\in A\backslash\supp v$, and $b$ is as above, then $K_{v<a}$ cotains the inverse image under the continuous endomorphism, $A\stackrel{\cdot b}{\to} A$ of $K_{r\leq 1}$ which is a neighbourhood of $0$. Then Proposition \ref{fact_12_1} e) is satisfied.
\end{proof}
% The following conditions [Hub 93a, Lemma 3.26],[Hub 93b, Lemma 3.3.(i)] or Prop III.1.4(i)/III.4.4.(i) in Morel's text.
\begin{proposition}
    \label{proposition_12_2}
    Let $A$ be a topological ring, $A^+$ be an open subring of $A$ which is integrally closed in $A$ and $\alpha\in A\backslash A^+$. Then there is $v\in\Spa(A,A^+)$ such that $v(\alpha)>1$.
\end{proposition}
\begin{proof}%later review the references.
    By assumption, $\alpha$ is not integral over $A^+$ ,hence $S=\{p(\alpha)\sep p\in A^+[T]\text{ monic}\}$ is a multiplicative subset of $A$ not contaiing $0$. By Example \ref{example_12_3}, there is $\p\in\Spec A$ maximal among all ideals of $A$ disjoint from $S$. It follows from the maximality of $\p$, that for every $a\in A\backslash\p$, there is $b\in A$ and a monic polynomial $P\in A^+[T]$ such that $P(\alpha) -ab\in\p$. Since $\p$ is disjoint from $S$, the image $\alpha\mod\p$ of $\alpha$ in $k(\p)$ is not integral over the image $A^\#$ of $A^+$ in $k(\p)$. By Proposition 1.4.1, this is a valuation ring $R$ of $k(\p)$ such that $A^\#\subseteq R$ and $\alpha\not\in R$. Consider $v$, the compsition,
    \begin{equation*}
        A\to k(\p)\stackrel{\text{valuation defined by R Prop 1.4.1}}{\to}\Gamma.
    \end{equation*}
    Since $\alpha\mod\p\not\in R, v(\alpha)>1$. By the definition, $K_{v\leq1}$ contains $A^+$ as $A^+$ and $\p\subseteq A^\#\subseteq R$, hence $K_{v\leq1}$ is open. If $a\in A$ and $v(a)\not=0$< then $a\not\in \p$ and we have $b\in A$ and a monic $P\in A^+[T]$ such that $P(\alpha)-ab\in\p$. Since $P$ is monic and $v(A^+)\subseteq[0,1]_{\Gamma\cup\{0\}}$, $P(\alpha)>1$. Since $\supp v=\p$, $v(ab)=v(P(\alpha))>1$. By Proposition \ref{fact_12_4}, $v$ is continuous. 
\end{proof}
\begin{corollary}
    \label{corollary_12_1}
    For a subset $A^+$ of $A$, the following conditions are equivalent.
    \begin{enumerate}[a).]
        \item $A^+ = \bigcap_{x\in\Spa(A,A^+)} K_{x\leq 1} = \{a\in A\sep \ab{a}_x\leq 1$ for all $x\in Spa(A,A^+)\}$
        \item $A^+$ is the intersection of a set of open and integrally closed subrings of $A$.
    \end{enumerate}
\end{corollary}
\begin{proof}
    For b)$\Rightarrow$a), note that $\subseteq $ in a) is trivial and $\supseteq$ is by Proposition \ref{proposition_12_2}. For a)$\Rightarrow$b), note that $K_{x\leq1}$ for $x\in\Spa(A,A^+)$ are integrally closed (by Fact 1.3.5) and open (by Proposition \ref{fact_12_1} a)) subrings of $A$.
\end{proof}
%6/8
\begin{proposition}
    \label{fact_13_5}
    Let $\eta\in\Cont(A)$ and $s\in\Spv(A)$ a specialization of $\eta$. Then $s\in\Cont(A)$.
\end{proposition}
\begin{proof}
    Let $a\not\in\supp s$. Then by Proposition \ref{fact_6_4}, $\supp\eta\subseteq\supp s$. Hence $\ab{a}_{\eta}\not=0$. Thus $K_{\eta<a}$ is a neighbourhood of $0$ in $A$ by Proposition \ref{fact_12_1} e). By Fact 1.3.6, $K_{\eta<a}\subseteq K_{s<a}$ hence $K_{s<a}$ is a neighbourhood of $0$ in $A$ and again by Proposition \ref{fact_12_1} e), hence for $s$ as well.
\end{proof}
WIthout further assumptions, $\Spa$ and $\Cont$ (to be equipped with a topology below) fail to be spectral spaces and there are ultrafilters $\U$ on $\Spa$ such that $\Cont(A)\in\U$ but $\glim_{\Spa A}\U$ fails to have any continuous specializations.
\par The following condition on $A$ are often useful.
\begin{align*}
    &\text{If $I\subseteq A$ is an open ideal, $I^2$ is also open} \tag{oisq}\\
    &\text{Every open ideal $I\subseteq A$ contains a finitely generated open ideal of $A$.}\tag{oifin}
\end{align*}
\begin{proposition}
    \label{fact_13_6}
    If $I$ and $J$ are open ideals of a topological ring $A$ satisfying (oisq), then $IJ$ is open. In particular, by induction on $n$, all $I^n$ are open.
\end{proposition}
\begin{proof}
    As $I\cap J$ is open $(I\cap J)^2 \subseteq I\cdot J$ is also open by (oisq).
\end{proof}
\begin{proposition}
    \label{fact_13_6_dash}
    If $a\in A^\dcirc$ and $x\in\Cont(A)$, then $\ab{a}_x<1$.
\end{proposition}
\begin{proof}
    By assumption, we have,
    \begin{equation*}
        \lim_{n\to\infty}\ab{a}_x^n = 0
    \end{equation*}
    in $\Gamma_x\cup\{0\}$. By the continuity of $x$, the above equality does not hold if $\ab{a}_x\geq1$. 
\end{proof}
We further introduce following conditions.
\begin{align*}
    &\text{The ideal generated by $A^{\dcirc}$ of $A$ is open}\tag{tn;0}\\
    &\text{$A^{\dcirc}$ is a neighbourhood of $0$ in $A$.}\tag{tn;1}
\end{align*}
Obviously (tn;1)$\Rightarrow$(tn;0). %All four conditions will be easily verified for $f$-adic map = Huber map.
\begin{proposition}
    \label{fact_13_7}
    If $A$ contains an open subring $B$ satisfying one of (oifin),(oisq), (tn;0),(tn;1), then $A$ also has that property.
\end{proposition}
\begin{definition}
    \label{definition_13_4}
    We say that $x\in\Spa(A)$ can be made continuous if $K_{x<a}$ is a neighbourhood of $0$ in $A$.
\end{definition}
\begin{remark}
    The condition in Definition \ref{definition_13_4} is already satisfied if $x$ is continuous by Propositoin \ref{fact_12_1} e).
\end{remark}
\begin{proposition}$\:$
    \label{fact_13_8}
    \begin{enumerate}[a).]
        \item If $x\in\Spv(A)$ can be made continuous then $K_{x<\gamma}$ is a neighbourhood of $0$ where $\gamma\in\cGamma$.
        \item If $K_{x\leq 1}$ is a neighbourhood of $0$ then $K_{x<\gamma}$ is a neighbourhood of $)$ where $\gamma\in\cGamma\backslash\{0\}$.
        \item If $K_{x\leq 1}$ is a neighbourhood of $0$, then $\cGamma_x=\{1\}$ or $x$ can be made continuous.
    \end{enumerate}
\end{proposition}
\begin{proof}
    For b), If $\gamma>1$ then $K_{x\leq1}\subseteq K_{x<\gamma}$ and the claim follows. Otherwise $\gamma<1$ and there is $a\in A$ such that $\gamma\ab{a}_x\geq1$. In that case, $\gamma\ab{a}_x^2>1$ and $K_{x<\gamma}$ contains the preimage under the continuous map $A\stackrel{a^2}{\to} A$ of $K_{x\leq1}$, a neighbourhood of $0$.
    \par a). For $\gamma=1$, this follows from Definition \ref{definition_13_4} and for $\gamma\not=1$, we can apply b). 
    \par c). Assume $\cGamma_x\not=\{1\}$. There is $\gamma\in\cGamma_x\backslash\{1\}$. Without loss of generality, assume $\gamma<1$. Then $K_{x<1}$ is a neighbourhood of $0$ as it contains $K_{x<\gamma}$, a neighbourhood of $0$ by b).
\end{proof}
%current Current
\begin{proposition}
    \label{proposition_13_3}
    Let $A$ be a topological ring with (oisq) and (tn;0), and let $\eta\in\Spv(A)$. Assume $\ab{x}_\eta<1$ for all $x\in A^{\dcirc}$. In the following the word "specialization" always means "specialization in $\Spv$".
    \begin{enumerate}[a).]
        \item There is a continuous specialization of $\eta$ if and only if $\eta$ can be made continuous in the sense of Definition \ref{definition_13_4}.
        \item If this is the case, then,
            \begin{enumerate}[I).]
                \item The following subset of $\Gamma_x$ is a convex subgroup containing $\cGamma_x$.
                \begin{equation*}
                    s\Gamma_\eta \defeq \{\gamma\in\Gamma_x\sep K_{\eta<r}\text{ is a neighbourhood of $0$ in $A$}\}.
                \end{equation*}
                \item The horizontal specialization $s=\eta|_{s\Gamma_\eta}$ (which by I) is well-defined) is continuous.
                \item The sets of specializations of $s$ and of continuous specializations of $\eta$ coincide.
                \item Let $\Omega=\mathcal{R}(f_0|f_1,\cdots,f_m)$ with $(f_i)_{i=0}^{m+1}\in A^{m+1}$ be such that the ideal of $A$ generated by $K_{\eta\leq f_0}$ is open, then $s\in\Omega$ if and only if $\eta\in\Omega$.
            \end{enumerate}
    \end{enumerate}
\end{proposition}
\begin{proof}
    Assume that $\eta$ has the continuous specialization $s$. As $K_{s\leq 1}\subseteq K_{\eta\leq 1}$ by Fact 1.3.7 and $K_{s<1}$ is a neighbourhood of $0$ by Proposition \ref{fact_12_1} e)., $K_{\eta\leq1}$ is a neighbourhood of $0$. By Proposition \ref{fact_13_8}, $\cGamma_\eta=\{1\}$ or $\eta$ can be made continuous.
    If $\cGamma_x=\{1\}$, then $A=K_{\eta\leq1}$ and $K_{\eta<1}$ is an ideal in $A$. By assumption about $\eta$, $A^\dcirc\subseteq K_{\eta<1}$ and by (tn0), $K_{\eta<1}$ is an open ideal of $A$. Then $\eta$ can be made continuous.
    \par Assume that $\eta$ can be made continuous. To finish a), and to show b), it suffices to show I)-IV) under this assumption. I) follows from Proposition \ref{fact_13_8} a). That is $\cGamma_\eta\subseteq s\Gamma_\eta$. The convexity of $s\Gamma_\eta$ as a neighbourhood of $\Gamma_\eta$ is obvious. One easily checks that the subgroup property follows (and I) is fully shown) where it is shown that $\gamma^2\in s\Gamma_\eta$ when $\gamma\in s\Gamma_\eta$. Note that
    \begin{equation*}
        \min\{\gamma_1,\gamma_2\}^2\leq \gamma_1\gamma_2\leq \max\{\gamma_1,\gamma_2\}^2.
    \end{equation*}
    It is sufficient to do this where $\gamma\leq1$. When $\gamma\in \cGamma_\eta$, then $\gamma^2\in\cGamma_\eta$ and the claim follows from a). Otherwise, $\gamma\ab{a}_\eta<1$ for all $a\in A$. Consider the (open as $\gamma\in s\Gamma$) ideal of $A$ generated by $K_{\eta<\gamma}$. By (oisq) and Proposition \ref{fact_13_6}, $I^3$ is also open. The elements of $I^3$ have the form $\alpha=\sum_{i=1}^n a_ib_ic_id_i$ with $a_i\in A$ and $b_i,c_i,d_i\in K_{\eta<\gamma}$. By b), $\ab{a}_\eta,\ab{b}_\eta<\ab{a_i}_\eta\gamma<1$. Hence 
    \begin{equation*}
        \ab{\alpha}_\eta\leq\max_{1\leq i\leq n}\{\ab{a_i}_\eta,\ab{b_i}_{\eta},\ab{c_i}_\eta,\ab{d_i}_\eta\}<\max_{1\leq i\leq n}\{\ab{c_i}_\eta,\ab{d_i}_\eta\}<\max_{1\leq i\leq n}\gamma\cdot\gamma = \gamma^2.
    \end{equation*} 
    Thus, $K_{\eta<\gamma^2}$ contains the open ideal $I^3$, and the claim follows.
    \par For II), note that $\ab{a}_s\not=0$ implies $\ab{a}_\eta\in s\Gamma_\eta$. hence by the definition,
    \begin{equation*}
        K_{s<a} = K_{\eta<a}.
    \end{equation*}
    This is a neighbourhoodof $0$ by definition of $s$.
    \par For IV), by Fact 1.3.10, it suffices to show that $\ab{f_0}_s\not=0$. That is $\ab{f_0}_\eta\in s\Gamma_\eta$. if $\ab{f_0}_\eta\in\cGamma_\eta$, then this follows from a). Otherwise $\ab{a}_\eta\ab{f_0}_\eta<1$ for all $a\in A$.
    Moreover, the ideal $I$ of $A$ generated by $K_{\eta<f_0}$ is open by our assumption about $f_0$. By (0;aq), $I^2$ is also open. Its elements have the form, $\alpha=\sum_i^n a_ib_ic_i$ with all $a_i\in A, b_i,c_i\in K_{\eta\leq f_0}$. Therefore,
    \begin{equation*}
        \ab{\alpha}_\eta\leq\max_{1\leq i\leq n}\{\ab{a_i}_\eta,\ab{b_i}_\eta,\ab{c_i}_\eta\}<\max_{1\leq i\leq n}\ab{c_i}_\eta\leq\ab{f_0}_\eta.
    \end{equation*}
    Then $K_{\eta<f_0}$ contains the open ideal $I^2$, hence is a neighbourhoodof $0$. Thus $\ab{f_0}_\eta\in s\Gamma_\eta$ as claimed.
    %Proof unfinished and will be continued on friday.
    \par For III, every specialization of the continuous. (II) s is continuous by \ref{fact_13_5}. Let $t\in \Spv(A)$, be any continuous specialization of $\eta$. For every rational open subset $\Omega = \cR(f_0|f_1,\cdots,f_n)$, we want to show that $t\in\Omega$ implies $s\in\Omega$. But $t\in\Omega$ implies $\ab{f_0}_t\not=0$ by definition, hence $K_{t< f_0}$ is a neighbourhood of $0$ by \ref{fact_13_1} d) (and the containing of $0$), hence $K_{t\leq f_0}\supseteq K_{t<f_0}$ is a neighbourhood of $0$ By Fact 1.3.7 $K_{\eta\leq f_0}\supseteq K_{t\leq f_0}$ hence this is a neighbourhood of $0$ as well. Hence IV is applied and $\eta\in\Omega$ follows since $t\in\Omega$ and $t$ is a specialization of $\eta$. 
    Thus, $\eta\in\Omega$ by IV and II is shown.
\end{proof}
%6/12
\begin{definition}
    \label{definition_14_5}
    The strong topology on $\Cont(A)$ or on $\Spa(A,A^+)$ is the one induced from $\Spv$. That is for which the set of all $\cR(f_0|f_1,\cdots,f_m)$ for arbitrary $n\in\N$ and $(f_i)_{i=0}^n\in A^{n+1}$ form a topology base. 
    \par If $A$ satisfies (oisq), the weak topology on both spaces is the one for which the set $\Rat_A = \Rat_A^{(t)}$ of $\cR(f_0|f_1,\cdots,f_n)$ for which $\langle f_0,\cdots,f_n\rangle_A$ is an open ideal of $A$ is a topology base.
\end{definition}

\begin{remark}
    Note that $\cR(f_0|f_1,\cdots,f_m)\cap\cR(g_0|g_1,\cdots,g_m) = \cR(f_0g_0|f_ig_j)$ and 
    \begin{equation*}
        \langle f_ig_j\rangle_A = \langle f_0,\cdots,f_n \rangle_A \langle g_0,\cdots,g_m\rangle_A
    \end{equation*} 
    is open by (oisq). Hence $\Rat_A^{(t)}$ is closed under finite intersection. In most case, $A$ will be equipped with a topology and may the notation $(t)$ be dropped.
    When the ring is equipped with a topology only the elements of $\Rat_A^{(t)}$ will be called rational open subsets of $\Cont(A)$ or $\Spa(A,A^+)$.
\end{remark}
\begin{proposition}
    \label{fact_14_9}
    If $A$ has (oisq) and (oifin), then these topologies coincide.
\end{proposition}
\begin{proof}
    Let $U = \cR(f_0|f_1,\cdots,f_n)\subseteq X$ where $X=\Cont(A)$ or $X=\Spa(A,A^+)$. For $x\in U$, we want to show that there is $\Omega\in\Rat_A^{(t)}$ such that $x\in\Omega\subseteq U$. As $x$ is a valuation, $K_{x<f_0}$ is open and by (oifin), there are finitely many $\varepsilon_1,\cdots,\varepsilon_l\in K_{x<f_0}$ such that $\langle \varepsilon_1,\cdots,\varepsilon_l\rangle_A$ is open. Then $x\in\Omega=\cR(f_0|f_1,\cdots,f_n,\varepsilon_1,\cdots,\varepsilon_l)\subseteq U$ and $\Omega\in\Rat_A^{(t)}$. 
\end{proof}
\begin{proposition}
    \label{fact_14_10}
    $\Cont(A)$ and $\Spa(A,A^+)$ are $T^0$ for the strong topology. This also holds for the weak topology when it coincides with the strong one. 
\end{proposition}

\begin{proposition}
    \label{proposition_14_4}
    Let $A$ be a topological ring with (oisq), (oifin) and (tn0). For an ultrafilter $\U$ on $\Cont(A)$, let $\sV$ be its image on $\Spv A$. The following conditions are equivalent.
    \begin{enumerate}[a).]
        \item $\U$ has a generic limit in $\Cont(A)$.
        \item $\U$ has some limit in $\Cont(A)$. 
        \item The generic limit $g=\glim_{\Spv(A)}\sV$ of $\sV$ in $\Spv A$ can be made continuous in the sense of Definition \ref{definition_13_4}. 
        \item The set 
        \begin{equation}
            \{a\in A\sep \{v\in\Cont(A)\sep \ab{a}_x<1\}\in\U\}
            \label{equation_14_4}
        \end{equation}
        is a neighbourhood of $0$ in $A$.
    \end{enumerate}
    In this case, a generic limit $\gamma$ of $\U$ in $\Cont(A)$ is given by the continuation $g|_{s\Gamma_g}$ of $g$. If $\Spa(A,A^+)\in\U$, then $\gamma\in\Spa(A,A^+)$ and $\gamma$ is a generic limit of the intersection of $\U$ iwth the set of subsets of $\Spa(A,A^+)$.
\end{proposition}
\begin{proof}
    a)$\Rightarrow$b) is obvious. For b)$\Rightarrow$c), if $\lambda$ is any limit of $\U$, it is a limit of $\sV$ in $\Spv(A)$ by the continuity of  (recall that we are dealig with the strong topology Propositio n\ref{fact_14_9}). Hence $\lambda$ is a specialization of $g$ in $\Spv(A)$. By Proposition \ref{proposition_13_3} a), $g$ can be made continuous.
    \par For c)$\Rightarrow $ a). If $g$ can be made continuous and $\gamma$ is its continuation. This is in the closure of $g=\glim_{\Spv(A)}\sV$, hence is a limit of $\sV$. As $\Cont(A)$ carries the induced topology, it is a limit of $\U$ in $\Cont(A)$. Let $\lambda$ be any limit of $\U$ in $\Cont(A)$. Then $\lambda$ is a limit of $\sV$ in $\Spv(A)$. Hence Proposition \ref{proposition_13_3}, a specialization of $\gamma$ in $\Spv(A)$. Using once again the fact that $\Cont(A)$ carries the induced topology, $\lambda$ is a specialization of $\gamma$ in $\Cont(A)$, showing that $\gamma$ is a $\glim_{\Cont(A)}\U$. 
    \par We have shown the equivalence of first three statement and also the 
    \par Since 
    \begin{equation*}
        \{x\in\Cont(A)\sep \ab{a}_x<1\} = \{x\in\Spv(A)\sep \ab{a}_x\}\cap\Cont(A)
    \end{equation*}
    as $a\in A$ belongs to the subset considered in \eqref{equation_14_4} if and only if 
    \begin{equation*}
        M = \{x\in\Spv(A)\sep \ab{a}_x<1\}\in\sV.
    \end{equation*}By the description of $g$ in Proposition 1.3.1, 
    \begin{equation*}
        \{a\sep M_a\in\sV\}
    \end{equation*}
    coincides with $K_{g<a}$. By Definition \ref{definition_13_4}, thisis a neighbourhood of $0$ in $A$ if and only if $g$ can be made continuous, and c)$\Rightarrow $d) follows.
    \par Finally let $\Spa(A,A^+)\in\U\subseteq\sV$. We have 
    \begin{equation}
        \Spa(A,A^+)=\left[\bigcap_{\alpha\in A^+}\cR_{\Spv(A)}(1|\alpha)\right]\cap\Cont(A)
        \label{equation_14_4_2}
    \end{equation}
    and since the open sets $\cR_{\Spv(A)}(1|\alpha)$ are all closed under intersection. $g$ is contained in all of them, hence $\ab{\alpha}_g\leq1$ where $\alpha\in A^+$. Since $\gamma$ is a horizontal specialization of $g$, this implies $\ab{\alpha}_\gamma\leq1$ and $\gamma$ belongs to the filter generated in the outher $\bigcap$ in \eqref{equation_14_4_2}.
    Since $\gamma$ is also continuous, $\gamma\in\Spa(A,A^+)$ by \eqref{equation_14_4_2}. Since $\Spa(A,A^+)$ carries the topology induced from $\Cont(A)$, $\gamma$ is a limit of 
    \begin{equation*}
        \U' = \U\cap(\text{subsets of $\Spa(A,A^+)$}).
    \end{equation*}
    If $\lambda$ is any limit of $\U'$ in $\Spa(A,A^+)$, it is a limit of the image $\U$ of $\U'$ under the continuous map $\Spa(A,A^+)\to \Cont(A)$. Hence it is a specialization of $\gamma$ in $\Cont(A)$. Hence also in $\Spa(A,A^+)$ which has the induced topology.
\end{proof}
\begin{corollary}
    \label{corollary_14_2}
    If $A$ is a topological ring satisfying (oisq),(oifin), and (tn;1). Then $X=\Cont(A)$ and all $X=\Spa(A,A^+)$ are spectral spaces and the elements of $\Rat_X^{(t)}$ are quasicompact open subsets of $X$.
\end{corollary}
\begin{proof}
    If $a\in A^\dcirc$, then 
    \begin{equation*}
        \{x\in\Cont(A)\sep \ab{a}_x<1\}=\Cont(A)\in\U 
    \end{equation*}
    (in the situation of Proposition \ref{proposition_14_4}), hene the subset of $A$ contained in \eqref{equation_14_4} contains $A^\dcirc$ which is a neighbourhood of $0$ if (tn1) holds. It follows that all ultrafilters on $X$ has generic limits. To show that the set of open subsets of $X$ which are closed under generic limit of ultrafilters in $X$ contains $\Rat_X^{(t)}$ (which is a topology base). Let 
    \begin{equation*}
        \Omega = \cR_X(f_0|f_1,\cdots,f_m)=\Omega'\cap X,
    \end{equation*}
    where 
    \begin{equation*}
        \Omega' = \cR_{\Spv(A)}(f_0|f_1,\cdots,f_m),
    \end{equation*}
    with $I=\langle f_0,\cdots,f_m\rangle_A$ being open. Let $\U'$ be an ultrafilter on $X$, $\U$ its image in $\Cont(A)$ and the remaining situation as in Proposition \ref{proposition_14_4}. Then $g\in\Omega'$ as $\Omega'\in\sV$ (note $\Omega'\cap X = \Omega\in\U'$) and $\Omega'$ is closed under $\glim_{\Spv(A)}$ by Proposition 1.3.1. We claim that $K_{g\leq f_0}$ is a neighbourhood of $0$ in $A$. Then $g\in\Omega'$ implies $\glim_X\U'=\gamma\in\Omega'\cap X=\Omega$. Assume ? that there is $a\in A$ such that $\ab{af_0}\geq1$. Then $K_{g\leq f_0}$ contains the preimage under $A\stackrel{f_0}{\to}A$ of $K_{g\leq 1}$, which is in a neighbourhood of $0$ as it contains $K_{g\leq 1}$ and $g$ ca nbe made continuous. Hence $K{g\leq f_0}$ is a neighbourhood of $0$ in this case. 
    \par If no such a exists , $K_{g\leq f_0}$ contains $I^2$, as elements of $I^2$ have the form 
    \begin{equation*}
        \alpha = \sum_{i,j=1}^m a_{ij}f_if_j.
    \end{equation*}
    Hence,
    \begin{equation*}
        \ab{\alpha}_g\leq\max_{\substack{0\leq i\leq m\\0\leq j\leq m}}\ab{a_{ij}}_g\ab{f_i}_g\ab{f_j}_g \leq \ab{a_{00}}_g\ab{f_0}_g^2<\ab{f_0}_g 
    \end{equation*}
    as $g\in\Omega$.
\end{proof}

%6/15

\begin{remark}
    Even when $\Spa(A,A^+)\in\U$ is the situation of Proposition \ref{proposition_14_4}, $\glim_{\Cont(A)}=\glim_{\Spa(A,A^+)}(\U\cap\Ouv_{\Spa(A,A^+)})$ may have specialization in $\Cont(A)$ which do not belong to $\Spa(A,A^+)$. (However, horizontal specialization of $g$ are the elements of $\Spa(A,A^+)$). For example $A=\T_1,A^+ = A^\circ$ and $g$ is given by $\norm{\cdot\vert \T_1}$, $\U$ is the principal filter for $g$, then $g = \glim_{\Cont(\T_1)}$ and the valuation (1.1.1) (denoted by $N_1$) is nto in $\Spa(\T_1,\T_1^\circ)$.
\end{remark}

\begin{proposition}
    \label{fact_15_11}
    If $X = \Cont(A)$ or $\Spa(A,A^+)$, and $x\in X$ is closed then $\cGamma_x = \Gamma_x$.
\end{proposition}
\begin{proof}
    Problem 5 from Sheet 8
\end{proof}

\begin{proposition}
    \label{fact_15_12}
    If $X = \Cont(A)$ or $\Spa(A,A^+)$, and $s\in X$ is a specialization of $\eta\in X$ in $X$. Let $x\in\Spv(A)$ is a specialization in $\Spa(A)$ of $\eta$ such that $s$ is a specialization in $\Spv(A)$ of $x$, then $x\in X$ and the two specializations in $\Spv(A)$ involving $x$ also hold in $X$ 
\end{proposition}

\begin{proof}
    Sheet 8 Problem 7. There are ? Huber (1.5. of the script) for $\Spv$ also hold in $X$. Definition \ref{definition_14_6} and Fact 13 about this are omitted
\end{proof}

\begin{proposition}
\label{proposition_15_5}
Let $A$ be a topological ring and $\mathcal{M}$ be the set of open ideals in $A$.
\begin{enumerate}[a).]
    \item If $A$ has (oifin) then every ideal $I$ not in $\mathcal{M}$ is contained in some ideal of $A$ which is $\subseteq$-maximal in the set of ideals of $A$ which are not in $\mathcal{M}$.
    \item If $A$ has (oisq), then $\mathcal{M}$ is an Oka-system. If $\p$ is any ideal of $A$ which is $\subseteq$-maximal with the property of begin $\not\in\mathcal{M}$. Then $\p$ is a prime ideal which is closed but not open and such that every ideal strictly contains it is open, or $\p$ is a dense maximal ideal.
\end{enumerate}
\end{proposition}
\begin{proof}
    Let $\mathcal{N}$ be the set of ideals in $A$ which are not open. For a), let $\mathcal{L}\subseteq\mathcal{N}$ be linearly $\subseteq$-ordered non-empty subset. Let $I = \bigcup_{J\in\mathcal{L}}J$. Then $I$ is an ideal in $A$. If $I$ is open then by (oifin), there are $f_1,\cdots,f_m\in I$ such that 
    \begin{equation*}
        K = \langle f_1,\cdots,f_m\rangle_A
    \end{equation*}
    is open. As $\mathcal{L}$ is linearly orderd, there is $J\in\mathcal{L}$ containing all $f_i$. Hence $K\subseteq J$ and $J$ is open. This contradicts to that $J\in\mathcal{N}$. Applying Zorn's lemma, we derived the statement.
    \par For b),that $\mathcal{M}$ is an Oka system follows from Proposition Fact 3 and Proposition Fact 6. That $\p$ must be prime follows from \ref{proposition_12_1}. If $\p$ is not closed let $a\in\overline{\p}\backslash\p$. As $\overline{\p}$ is an ideal and $\p$ is prime, $a^2\in\overline{p}\backslash\p$. Now, $\p+a^2A$ is open by the maximality of $\p$, hence closed hence $\overline{p} = \p+a^2A$. Sinc $a\in\overline{p} = \p+a^2A$, there is $b\in A$ such that $a = a^2b\in\p$. As $a\in\p, 1 -ab\in \p$. Thus $1\in\overline{\p}$ as $a\in\overline{\p}$ and $\overline{\p}$ is an ideal.
    Now the previous consideration can be applied to all $a\in A$, and show that $A/\p$ is a field. Thus $\p$ is a dense maximal ideal.
\end{proof}

\begin{definition}
    \label{definition_15_7}
    A continuous varuation $v$ of $A$ and its equivalence class is called analytic if the closed prime ideal $\supp v$ fails to be open. Let $\Cont(A)_a$ and $\Spa(A,A^+)_a$ be the sets of $x\in\Cont(A)$ or $x\in\Spa(A,A^+)$ which are analytic.
\end{definition}

\begin{proposition}
    \label{fact_15_14} For a topological ring $A$ with (tn;0), the following conditions are equivalent.
    \begin{enumerate}[a).]
        \item The ideal of $A$ generated by $A^\dcirc$ is $A$.
        \item A has no open maximal ideal.
        \item $A$ has no open prime ideal.
        \item $A$ has no open proper ideal.
        \item For every subset $A^+\subseteq A$, all elements of $\Spa(A,A^+)$ are analytic.
    \end{enumerate}
\end{proposition}
\begin{proof}
    a)$\Rightarrow$d), as every open ideals is contained in $A^\dcirc$. d)$\Rightarrow$c)$\Rightarrow$b) and c)$\Rightarrow$d) are all trivial. d)$\Rightarrow$a) is by (tn;0). b)$\Rightarrow$d), every proper ideal $I$ is contained in some maximal ideal which is open when $I$ is. 
    \par e)$\Rightarrow $c) If $\P$ is an open ideal, its finite valuation given as elements of $\Spa(A,A^+)$. 
\end{proof}
\begin{definition}
    A ring is called analytic if the ideal of $A$ generated by $A^\dcirc$ is $A$.
\end{definition}
\begin{example}
    Eery ring with (E) from 3.1 ($A^\times$ intersects every neighbourhood of $0$) is analytic. In particular, every ring with (T) is analytic.
\end{example}

\begin{proposition}
    \label{fact_15_15}
    When $A$ is a topological ring with (oifin) and (tn;0) and $X=\Cont(A)$ or $X=\Spa(A,A^+)$. Then $X_0$ is a finite union of elements of $\Rat_X$. In particular, it is open.
\end{proposition}
\begin{proof}
    By our assumption, there are finitely man yelements $f_1,\cdots,f_m\in A^\dcirc$ generating an open ideal in $A$. Then $\p\in\Spec A$ is open if and only if $f_1,\cdots,f_m\in\p$. Then 
    \begin{equation*}
        X_a = \bigcup_{i=1}^m\cR_x(f_i|f_1,\cdots,f_m).
    \end{equation*}
\end{proof}

\begin{remark}
    If only (tn;0) holds then,
    \begin{equation*}
        X_a = \bigcup_{f\in A^\dcirc}\cR_X(f|)
    \end{equation*}
    is still open.
\end{remark}

\begin{proposition}
    \label{fact_15_16}
    Let $A\stackrel{p}{\to} B$ be a morphism of topological rings which is open. If $M\subseteq A$ is bounded then $p(M)$ is bounded in $B$.
\end{proposition}

\begin{proof}
    Exercise. Sheet 8 problem 1.
\end{proof}

\begin{proposition}
    \label{proposition_15_6}
    Let $A$ be a topological ring with (oisq), (oifin), and (tn0) and without dense maximal ideals. 
    Let $A^+$ be an open subring of $A$ which is a subset of $A^\circ$. For $m\in\N$ and $(a_i)_{i=1}^m\in A^m$, the following conditions are equivalent.
    \begin{enumerate}[a).]
        \item The ideal $I=\langle a_1,\cdots,a_n\rangle_A$ is open.
        \item For all $x\in\Spa(A,A^+)$, there is $1\leq i\leq m$ such that $\ab{a_i(x)}\not=0$.
    \end{enumerate}
\end{proposition}

\begin{proof}
    By the definition, $\max_{1\leq i\leq m}\ab{a_i}_x\not=0$ if and only if $I\subseteq\supp x$. If $I$ is open, no such $x$ can be analytic. Thus a)$\Rightarrow$b).
    \par Assume that $I$ is not open. By Proposition \ref{proposition_15_5}, there is a closed prime ideal $\p$ such that $I\subseteq \p$ and such that $\p$ is not open but every larger ideal is open. By (tn;0), there is $x\in A^\dcirc\backslash\p$. As $A^+$ is a neighbourhood of $0$, without losing generality, ($s\mapsfrom s^j$ for some positive $j\in\N$), $s\in A^+$.
    We denote the image of $A^+$ in $A/\p\subseteq k(\p)$ by $B$. By Proposition \ref{fact_15_16}, 
    \begin{equation*}
        B\subseteq(A/\p)^\circ\tag{1}.
    \end{equation*}
    Now $(s\mod\p)^{-1}\in k(\p)$ be not integral over $B$. Otherwise, there will be $(p_i)^d_{i=1}\in B$ such that 
    \begin{equation*}
        1 = \sum_{i=1}^d p_i(s\mod\p)^i = (s\mod\p)\sum_{i=0}^{d-1}p_{i+1}(s\mod\p)^i \in (s\mod\p)B.
    \end{equation*}
By a) and $(s\mod\p)\in (A/\p)^\dcirc$, this implies that $1$ is topologically nilpotent in $A/\p$. This contradicts the fact that $\p$ is closed in $A$.
By Proposition 1.4.2, there is a valuation ring $R$ of $k(\p)$ containing $B$ but not $(s\mod\p)^{-1}$. Let $\gamma$ denote the composition,
\begin{equation*}
    A\to A/\p \to k(\p)\stackrel{\text{valuation defined by $R$}}{\to}(k(\p)^\times/R^\times)\cup\{0\}.
\end{equation*}
There is a valuation and its support $\p$ is not open and $\gamma(a)\leq1$ for $a\in A^+$, as $a\mod\p\in B\subseteq R$, for such a. Thus the equivalence class of $\gamma$ is in $\Spa(A,A^+)_a$ and all $\gamma(f_i)=0$, contradicting b), provided that we can show that $\gamma$ is continuous. 
\par We first show that $K_{\gamma\leq s^m}$ is a neighbourhood of $0$, for all $m\in\N$> If $m=0$, this follows from $A^+\subseteq K_{\gamma\leq 1}$ and the openness of $A^+$. If there is $a\in A$ such that $\gamma(as^m)\geq1$, then $K_{\gamma\leq s^m}$ contains the preimage under $A\stackrel{a}{\to}A$ of $K_{\gamma\leq1}$, hence it is a neighbourhood of $0$.
Let $\alpha\in A\backslash\p$. Then $\p+\alpha A$ is an open ideal. Since $s\in A^\dcirc$, $s^n\in\p+\alpha A$ for some positive integer. Thus there is $b\in A$ such that $s^n-b\alpha\in\p$ or 
\begin{equation*}
    \gamma(b)\gamma(\alpha) = \gamma(s^r)
\end{equation*}
hence $\gamma(b)\gamma(\alpha)>\gamma(s^{n+1})$. (Note that $(s\mod\p)\in\m_R$ as $(s\mod\p)^{-1}\in R$, hence $\gamma(s)<1$).
Thus $b^{-1}K_{\gamma\leq s^{n+1}}$ which is a neighbourhood of $0$ which is contained in $K_{\gamma<\alpha}$. By Fact 1 d), $\gamma$ is continuous.
\end{proof}

%6/19
\begin{corollary}
    \label{corollary_16_3}
    Under the assumption of Proposition \ref{proposition_15_6} on $A$ and $A^+$, $\langle a_1,\cdots,a_m\rangle_A = A$ if and only if 
    \begin{equation*}
        \max_{1\leq i\leq m}\vert a_i\vert_x >1
    \end{equation*}
    for all $x\in\Spa(A,A^+)$.
\end{corollary}
\begin{proof}
    The necessity is obvious. If the condition holds, $I=\langle a_1,\cdots,a_m\rangle_A$ is open by Proposition \ref{proposition_15_6}. If it is different from $A$, then there is an open maximal ideal $\m$ such that $I\subseteq \m$. If $x\in\Spa(A,A^+)$ is given by the trivial valuation for $\m$, then $\vert a_i\vert_x=0$, contradicting to our assumption.
\end{proof}
\begin{remark}
    In particular, $a\in A^\times$ if and only if $\vert a\vert_x\not=0$ for all $x\in\Spa(A,A^+)$ under the same assumption as in Proposition \ref{proposition_15_6}.
\end{remark}

\begin{lemma}
    If the map $gf$ is open then $g$ is open provided that $f$ and $g$ are continuous homomorphism of topological groups.
    \label{lemma_16_17_1}
\end{lemma}
\begin{proof}
    Such $g$ are open if and only if $g(U)$ is a neighbourhood of $1$ if $U$ is a neighbourhoodof $1$. (Then as $g(U) = g(Ux^{-1})g(x)$ is a neighbourhood of $g(x)$ under $U$ is a neighbourhood of $x$). But $g(U)$ contains $(gf(f^{-1}(U)))$ and this is a neighbourhood of $1$ as $f^{-1}(U)$ is a neighbourhood of $1$ and $gf$ is open.. 
\end{proof}

\begin{proposition}
    \label{fact_16_17}
    Let $A$ be a topological ring and $M$ be a topological $A$-module. Let $m,n\in\N$ and $\vec{a} = (a_i)_{i=1}^m\in A^m$, $I= \langle a_1,\cdots,a_m\rangle,\vec{b} = (b_i)_{i=1}^n\in I^n$. If 
    \begin{equation*}
        M^n\stackrel{\vec{b}}{\to}M, (m_i)_{i=1}^n\mapsto \sum_{i=1}^n b_im_i
    \end{equation*}
    is open, then $M^m\stackrel{\vec{a}}{\to}M$ is also open.
\end{proposition}
\begin{proof}
    Set 
    \begin{equation*}
        b_i =\sum_{j=1}^m c_{ij}a_j.
    \end{equation*}
    Then we have the commutative diagram,
    \begin{center}
        \begin{tikzcd}
M^n \arrow[d, "(c_{ij})_{ij}"'] \arrow[rd, "\vec{b}"] &   \\
M^m \arrow[r, "\vec{a}"']                             & M
\end{tikzcd}
    \end{center}
    Now apply Lemma \ref{lemma_16_17_1}.
\end{proof}

\begin{remark}
    In particular, the openess of the map $\vec{a}:M^m\to M$ depends only on $I = \langle a_1,\cdots,a_m\rangle$.
\end{remark}

\begin{proposition}
    \label{fact_16_18}
    Let $A$ be a topological ring and $M$ be a topological $A$-module.Then the following conditions are equivalent.
    \begin{enumerate}[a).]
        \item For every neighbourhood $U$ of $0$, there are $m\in\N$ and $\vec{a}\in U^m$ such that $M^m\stackrel{\vec{a}}{\to}M$ is open.
        \item When $I\subseteq A$ is an open ideal, then a) holds with $U=I$. (A direct corollary of Proposition \ref{fact_16_17}). 
    \end{enumerate}
    If $A$ has (oifin), this is equivalent to 
    \begin{enumerate}[c).]
        \item If $m\in\N$ and $\vec{a}\in A^m$ generates an open ideal, then $M^m\stackrel{\vec{a}}{\to}M$ is open. (Also trivial as $I$ in b) can be assumed to be finite generated under (oifin)).
    \end{enumerate}
\end{proposition}
\begin{definition}
    An $A$-module $M$ is called adic if and only if the equivalent conditions a) and b) hold in Proposition \ref{fact_16_18}. An $A$-algebra is called adic if and only if it is adic as an $A$-module.
    \par $A$ is called $\id$-adic if and only if $\id_A$ is an adic morphism.
\end{definition}

\begin{proposition}
    \label{fact_16_19}
    Every $\id$-adic ring satisfies (oifin) and (oisq).
\end{proposition}
\begin{proof}
    If $I$ is an open ideal. Then there are $(a_j)^m_{j=1}\in I^m$ such that $A^m\stackrel{\vec{a}}{\to}A$ is open. Then $\langle a_1,\cdots,a_m\rangle_A$ is open and it is the image of $\vec{a}$. Thus (oifin). Moreover, the image of $I^m$ under $\vec{a}$ is also open and contained in the ideal square $gI$, which is therefore open.
\end{proof}

\begin{proposition}
    \label{fact_16_20}
    If $A$ is analytic or satisfies (3.1.E), then very $A$-module (hence every $A$-algebra) is adic. in particular, $A$ is $\id$-adic. 
\end{proposition}
\begin{proof}
     If $A$ is analytic, there are $(a_i)_{i=1}^m\in (A^\dcirc)^m$ such that $\langle a_1,\cdots,a_m\rangle=A$. Then also $\langle a_1^j,\cdots,a_m^j\rangle =A$ for all $j\in\N$. Thus $I=A$ for all open ideal $I$, and one can take $n=1, a_1=1$ in Proposition \ref{fact_16_18} b). Condition (3.7 E0)likewise implies that $A$ is the only open ideal.
\end{proof}

\begin{proposition}
    \label{fact_16_21}
    Let $A\stackrel{f}{\to} B$ be a morphism of topological rings and $M$ be a topological $B$-module which is adic as an $A$-module .The n$M$ is adic as a $B$-module. In particular, if $B$ is an adic $A$-algebra, then $B$ is $\id$-adic.
\end{proposition}
\begin{proof}
Let $I\subseteq B$ be an open ideal. There are $\vec{a}\in(f^{-1}I)^m$ such that $M^m\stackrel{\vec{a}}{\to} M$ is open (induced by multiplicatio n$A\times M\to M$). Then $M^m\stackrel{f(\vec{a})}{\to} M$ (induced from the multiplication $B\times M\to M$) is also open.
\end{proof}

% \begin{proposition}
%     \label{fact_16_22}
%     If $A\stackrel{f}{\to} B$ is an adic morphismof topologica rings and $X\subseteq A$ bounded, then $f(X)$ is bounded in $B$.
% \end{proposition}

% \begin{proof}
%     Let $U\subseteq B$ be a neighbourhood of $0$, we must show the existence of a neighbourhood $V$ of $0$ such that $V f(X)\subseteq U$. Let $U'\subseteq U$ be a neighbourhood of $0$ such that $U'^2\subseteq U$, and for $W=f^{-1}(U)\subseteq A$. As $X$ is bounded, there is a neighbourhood $\Omega$ of $0$ in $A$ such that $\Omega X\subseteq W$. As $B$ is an adic $A$-algebra, there are $(\omega_i)_{i=1}^m$ such that $B^m\stackrel{\vec{\omega}}{\to}B$ is open. Then the image $V$ of ${U'}^m$ under $\vec{\omega}$ will do.
%     \par If all $U_i$
% \end{proof}

\begin{proposition}
    \label{fact_22}
    Let $B\subseteq A$ be an open subring, $M$ be a topoological $A$-module, $N$ an open sub-$B$-moudle of $M$. Then $M$ is adic as an $A$-module if and only if $N$ is adic over $B$. 
\end{proposition}
\begin{proof}
    
\end{proof}
\subsection{Nat-rings, Nat-modules and Huber Ring}

\begin{definition}
    A commutative (additively written) topological group is called nat (for non-Archimedian topological) if and only if it has a neighbourhood base of $0$ where each element is an open subgroup. A nat-ring is a topological ring $A$ such that $(A,+)$ is a nat. For a nat-ring $A$, a nat-A-module is a topological $A$-module $M$ such that $(M,+)$ is a nat. A morphism of nat-rings (of nat-$A$-modules) is a morphism of topological rings ($A$-modules) where source and target are nat.
\end{definition}

\begin{example}
    $\Q_p$ with topology given by $\ab{\cdot}_p$ or any other field with ultrametric absolute value.  Any ring with its discrete topology. 
\end{example}
\begin{remark}
    For any topological ring $A$, $A^{\mathrm{nat}}$ is a topological ring $A$ equipped with the topology where the open subgroups of $(A,+)$ are a neighbourhood base of $0$. Then $A^{\nat}$ is a nat-ring. Every morphism $A\to B$ of topological ring with $B$ being nat factors over $A^{\nat}$, and $\Cont(A^\nat)\to \Cont(A)$ is bijective.
\end{remark}

\begin{proposition}
    \label{fact_16_1}
    Let $R$ be a topological ring $M$ be a $R$-module which is nat, $X\subseteq M$ a subset of $M$, $X'\subseteq M$ be the subgroup generated by $X$. Then $X$ is bounded if and only if $X'$ is.
\end{proposition}
\begin{proof}
    $X'X = X$, thus $\Leftarrow$ is clear. For the other, let $U\subseteq M$ be a neighbourhood of $0$. Without loss of generality, $U$ is a subgroup of $(M,+)$. If $X$ is bounded, there is a neighbourhood $V$ of $0$ in $R$ such that $VX\subseteq U$. Since $U$ is a subgroup, $VX'\subseteq U$. Thus $X'$ is bounded.
\end{proof}

\begin{proposition}
    \label{fact_16_1_b}
    Let $A$ be a nat, $X\subseteq A$ be a subset, $X'$ be a subgroup generated by $X$. Then $X$ is power-bounded if and only if $X'$ is. (Respectively, it is topologically nilpotent and power bounded if $X'$ is).
\end{proposition}

\begin{proof}
    Let $Y = \bigcup_{n=0}^\infty X^n$ and $Y'\subseteq A$ be the subgroup generated by $Y$. Then $Y''$ also be the subgropu generated by $\bigcup_{n=0}^\infty {X'}^m$. $Y'\subseteq Y''$ is clear. Moreover, $X'Y'\subseteq Y'$, $(XY'\subseteq Y')$ is clear and 
    \begin{equation*}
        \{\xi\in A\sep \xi Y\subseteq Y\}
    \end{equation*}
    is a subgroup. Thus ${X'}^m\subseteq Y'$ by induction on $m$. Thus , $Y''\subseteq Y'$. By Proposition \ref{fact_16_1}, $X$ is powerbounded if and only if $Y$ is bounded if andonly if $Y'= Y''$ bounded if and only if $X'$ is power bounded.The latter case is similar.
\end{proof}

\begin{proposition}
    \label{fact_16_2}
    Let $A$ be a nat-ring, $X\subseteq A$ be a subset and $X'\subseteq A$ be the subring generated by $X$. Then the following condition are equivalent.
    \begin{enumerate}[a).]
        \item $X$ is powerbounded.
        \item $X'$ is bounded.
        \item $X'$ is powerbounded.
    \end{enumerate} 
\end{proposition}
\begin{proposition}
    \label{fact_16_3}
    Let $A$ be a nat ring, $(X_i)_{i=1}^m$ are power bounded subsets. Then the subring generated by $X_i$ is powerbounded. (Apply Proposition \ref{fact_16_2} to $\bigcup_{i=1}^m X_i$ which is bounded, see 3.1).
\end{proposition}
\begin{proposition}
    \label{fact_16_4}
    For bounded subgroups $B_1,B_2$ of a nat ring $A$ are contained in a bounded subring $B\subseteq A$.
\end{proposition}

\begin{proposition}
    \label{fact_16_5}
    For a nat ring $A$, $A^\circ$ is a subring of $A$ which is integrally closed in $A$, $A^\dcirc$ is an ideal in $A^\circ$ and $\sqrt{A^\dcirc}=A^\dcirc$. If $A$ is (tn;1), then $A^\circ$ is clopen in $A$.
\end{proposition}

\begin{proof}
    That $A^\dcirc$ is a subgroup of $(A^\circ,+)$ follows from Proposition \ref{fact_16_1} and $A^\circ A^\dcirc\subseteq A^\dcirc$ follows from Fact 3.1.5. That $\sqrt{A^\dcirc}=A^\dcirc$ follows from the definition. If $r^n\in A^\dcirc$ with $n>0$, then 
    \begin{equation*}
        \lim_{j\to\infty}r^{i+jn}=0
    \end{equation*}
    for all $0\leq i <n$. Hence $\lim_{j\to\infty}r^j=0$. If $x\in A$ is integral over $A^\circ$, then there are $n\in\N$ and $(p_j)_{j=0}^{n-1}\in A^\circ$ such that 
    \begin{equation*}
        x^n = \sum_{j=0}^{n-1}p_jx^j.
    \end{equation*}
    By Proposition \ref{fact_16_2}, there is a bounded subring $S\subseteq A$ containning all $p_j$. Then $X^N$ is contained in the subgroup generated by $S$ and $\{(x_i)_{i=0}^{n-1}\}$ which is bounded by Proposition \ref{fact_16_1} and Fact 3.1.3. If (tn;1) holds then $A^\dcirc$ is a neighbourhood of $0$ hence as in $A^\circ$ and $A^\circ$ is open, hence closed subgroup of $(A,+)$.
\end{proof}

\begin{remark}
    \label{remark_16_2}
    Note that $A^\circ$ can be unbounded (e.g. $A=\Q_p[T]/T^2\Q_p[T]$, $A^\circ = \Z_p+T\cdot\Q_p$) and $A^\dcirc$ will often happen to be topologically nilpotent. For example, $A=\C_p$, then 
    \begin{equation*}
        A^\dcirc = \{x\in\C_p\sep \vert x\vert<1\}
    \end{equation*}
    such that $(A^\dcirc)^2 = A^\dcirc$.
\end{remark}

\begin{proposition}
    \label{proposition_16_1}
    Let $A$ be a complete nat with (tn;1) (e.g. $A^\dcirc$ is open subgroup). If $A$ is sequentially compact, then $A^\times$ is an open subset of $A$, $A^\times\stackrel{a\mapsto a^{-1}}{\to} A$ is continuous, and $A$ has no dense proper ideals (hence no dense maximal ideals).
\end{proposition}

\begin{proof}
    $A^\times$ contains $A^\dcirc+1$ ( a neighbourhood of $1$), hence  by multiplicative invariance, it is a neighbourhood of each of its elements, by $(1+n)^{-1}=\sum_{j=0}^\infty n^j$. If $I\subseteq A$ is a dense ideal, $1\in\overline{I}$, hence $\overline{I}\cap A^\times \not=\emptyset $, hence $I\cap A^\times\not=\emptyset$, hence $I=A$.
\end{proof}

\begin{proposition}
    \label{fact_17_6}
    Let $A$ be a nat ring, $B$ be a nat $A$-algebra which is adic over $A$, $X\subseteq AA$ bounded. Then the image of $X$ in $B$ is bounded.
\end{proposition}

\begin{proof}
    Let $U$ be a neighbourhood of $0$ in $B$, we need to find a neighbourhood $V$ of $0$ in $B$ such that $Vf(X)\subseteq U$, where $A\stackrel{f}{\to}B$ is the homomorphism making $B$ into an $A$-algebra. There is a neighbourhood $U$ of $0$ in $B$ such that $U'U'\subseteq U$ and without loss of generality, $U$ and $U'$ can be assumed to be subgroups of $(A,+)$. Let $W=f^{-1}(U')\subseteq A$, there is $\Omega\subseteq A$ such that $\Omega X\subseteq W$ and by the assumption, there are $(\omega_i)_{i=1}^n\defeq\vec{\omega}$ such that $B^n\stackrel{\vec{\omega}}{\to}B$ is an open map. Then $V$ which is the image of ${U'}^n$ under this open map will do. For $r = f(\omega_i)u$ with $u\in U'$, then 
    $rf(X) \subseteq f(\Omega X)U'\subseteq f(W)U'\subseteq U'U'\subseteq U$, and write? $U$ as a subgroup and the general elements of $V$ is a sum of such $r$, we indeed have $Vf(X)\subseteq U$. 
\end{proof}

\begin{remark}
    \label{remark_17_2}
    In general, $A^\circ$ may well be unbounded in $A$ for example,
    \begin{equation*}
        A=\Q_p[T]/T^2,A^\circ = \Z_p+\Q_pT.
    \end{equation*}
    But every finite subseet of it is contained in some bounded subring.
\end{remark}

\begin{definition}
    \label{definition_17_2}
    Let $A$ be a ring and $I\subseteq A$ be an ideal. The $I$-adic topology on $A$ is the topology for which $(A,+)$ is a topological group and $(I^n)_{n=0}^\infty$ is a neighbourhood base of $0$.
\end{definition}

\begin{proposition}
    \label{proposition_17_2}
    For a nat ring $A$, the following conditions are equivalent.
    \begin{enumerate}[a).]
        \item $A$ is $\id$-adic satisfies (tn;1) and has an open bounded subring $B$.
        \item $A$ is $\id$-adic, satisfies (tn;0) and has an open bounded subring $B$.
        \item $A$ has an open bounded subring $B$ for which there is a finitely generated ideal $I\subseteq B$ such that the topology on $B$ is the $I$-adic one.
    \end{enumerate}
\end{proposition}

\begin{proof}
    For b).$\Rightarrow$ a), as $A$ is $\id$-adic and $A^\dcirc$ generates an open ideal ((tn;0)), there are $(a_i)_{i=1}^m\in {A^\dcirc}^m$ such that $A^m\stackrel{\vec{a}}{\to} A$ is an open map. The image of $B^m$ under this map is the open and topologically nilpotent since $B$ is powerbounded and all $a_i$ are topologically nilpotent.
    \par a).$\Rightarrow$ b), is trivial. 
    \par b),$\Rightarrow$ c), If the $(a_i)$ are as above, then for each $n\in\N$, the ideal generated by the $a_i^n$ is also open, as $A$ has (oisq). Then without loss of generality, all $a_i\in B$ (replace $a_i$ by $a_i^n$ which is in $B$ when $n$ is large). Then $I=\langle a_1,\cdots,a_m\rangle_B$ can be taken. Indeed, $B^m\stackrel{\vec{a}}{\to} B$ is open and $I^{\cdot(j+1)}$ is the image of $(I^{\cdot j})^{\oplus m}$ (using $\cdot$, $\oplus$ to distinguish ideal powers and Cartesian powers) under this open map, hence open by induction on $j$. If $U$ is a neighbourhood of $0$ in $A$, then since $B$ is bounded, there is a neighbourhood $U'$ of $0$ such that $U'B\subseteq U$ where $U$ is a subgroup of $(A,+)$ without loss of generality. As the $a_i$ are topologically nilpotent, there is $N$ such that $\vec{a}^\alpha = a_i^{\alpha_i}\cdots,a_m^{\alpha_m}\in U'$ where $\vert\alpha\vert\geq N$. Then $I^{\cdot N}\subseteq U$, showing $I^n$ are indeed a neighbourhood base of $0$.
    \par c)$\Rightarrow$ a). Since $I\subseteq A^\dcirc$ and (tn;1) holds. Then, only $\id$-adicity need to be verified. If $J$ is an open ideal in $A$ there is $N$ such thta $I^{\cdot N}\subseteq J$ (as the $I^n$ form a neighbourhood base of $0$). Then if $I^{\cdot N}\langle b_1,\cdots,b_l\rangle_B$ then $A^l\stackrel{\vec{b}}{\to}A$ is open. (it maps $(I^{\cdot j})^{\oplus l}$ onto $I^{\cdot j +N}$ which is open and hence $I^{\cdot j}$ form a neighbourhood base of $0$ in $A$). Since all $b_i\in J$, $A$ is $\id$-adic as claimed.
\end{proof}

\begin{definition}
    \label{definition_17_3}
    A nat-ring with these properties in Proposition \ref{proposition_17_2} is called a Huber ring (In Huber's note, it is refered as $f$-adic).
\end{definition}

\begin{definition}
    \label{definition_17_3_b}
    A ring of Def of a Huber ring is an open bounded subring of $A$. A pair of def is a pair $(B,I)$ where $B$ is an open bounded subring and $I\subseteq B$ a finitely generated ideal such that the topology on $B$ is $I$-adic.
\end{definition}

\begin{corollary}
    Every ring of def of a Huber ring can be extended to a pair of def.
\end{corollary}
\begin{proof}
    Proved in b)$\Rightarrow$c) in Proposition \ref{proposition_17_2}.
\end{proof}

\begin{proposition}
    \label{fact_17_7}
    Every Huebr ring is (oisq), (oifin) (from Fact 3.1.19), and (tn;1) (from Proposition \ref{proposition_17_2} b)). Moreover $A7\circ$ is clopen and integrally closed subring of $A$. If $A$ is (sequentially) complete (note that the $I^{\cdot j}$ form a countable neighbourhood base of $0$) then $A^\times$ is an open subset of $A$ and $A$ has no proper dense ideal (every maximal ideal is then closed).
\end{proposition}

\begin{proposition}
    \label{fact_17_8}
    Let $B$ be an open subring of the nat ring $A$. Then $B$ is Huber if and only if $A$ is. In this case, $B$ has a pair of def which is a pair of def for a).
\end{proposition}

\begin{proposition}
    \label{fact_17_9}
    Let $(B,I)$ be a pair of def for the Huber ring $A$. Then a neighbourhood $X$ of $A$ is bounded if andonly if there is $l\in\N$ suchthat $I^{\cdot l}X\subseteq B$ (where $I^{\cdot l}$ is an ideal power of $I$ in $B$ as before).
\end{proposition}

\begin{proposition}
    Every Huber ring has a countable neighbourhood base of $0$, the completion of a Huber ring is a Huber ring (If $(B,I)$ is a pair of def, then $\hat{B}$ is an open subring of $\hat{A}$ and the topology of $\hat{B}$ is $\hat{I}$-adic with $\hat{I}=\langle f_1,\cdots,f_n\rangle_{\hat{B}}$ where $I= \langle f_1,\cdots,f_m\rangle_B$).
\end{proposition}

\begin{definition}
    \label{definition_17_4}
    A Huber pair (affnoid pair in Huber) is a pair $A=(A^{\triangleleft},A^+)$ such that $A^{\triangleleft}$ is a Huber ring and $A^+\subseteq A^{\triangleleft\circ}$ is an open and integrally closed subring of $A^{\triangleleft}$ (hence $A^{\triangleleft\dcirc}\subseteq A^+$). 
    \par A Huber pair is complete if $A^{\triangleleft}$ (or, equivalently $A^+$) is (sequential) complete. A morphism $A\to B$ of Huber pairs is a continuous ring homomorphism $A^\triangleleft\to B^\triangleleft$ mappign $A^+$ to $B^+$. it is adic if $A^\triangleleft\to B^\triangleleft$ is adic
\end{definition}

\begin{theorem}
    \label{theorem_17_2}
    Let $A=(A^\triangleleft,A^+)$ be a Huber pair and $X=\Spa A$.
    \begin{enumerate}[a).]
        \item $A^+=\{a\in A^\triangleleft\sep \ab{a}_x\leq 1\text{ for all $x\in X$}\}$.
        \item $X$ is a spectral space.
    \end{enumerate}
    If $A$ is complete, we also have,
    \begin{enumerate}[a).]
        \setcounter{enumi}{2}
        \item The ideal $\langle a_1,\cdots,a_m\rangle_{A^\triangleleft}$ is open if and only if $\max_{1\leq i\leq m}\ab{a_i}_x>0$ for all $x\in X_a$ (the analytic points of $X$). This is the case if and only if $A^{\triangleleft m}\stackrel{\vec{a}}{\to}A$ is an open map.
        \item $A^{\triangleleft}=\langle a_1,\cdots,a_m\rangle_{A^\triangleleft}$ if and only if $\max_{1\leq i\leq m}\ab{x_i}_x>0$ for all $x\in X$.
    \end{enumerate}
    A morphism of Huber pairs defines a continuous map $\Spa B\to \Spa A$ which is the case of an adic morphisms is a spectral map.
\end{theorem}

\begin{proof}
    a) follows from Propositin 3.2.2. b) is Corollary 3.2.2. c) follows from Proposition 3.2.5 + $\id$-adicity. d) follows from Corollary 3.2.3. That $Y = \Spa B\to X=\Spa A$ is continuous is easily observed for the strong topologies. If $B^{\triangleleft}/A^{\triangleleft}$ is adic and $\Omega\in\Rat_X$, $\Omega=\cR(a_0|a_1,\cdots,a_m)$ with $\langle a_0,\cdots,a_m\rangle_A$ open. As $B/A$ is adic, $B^{m+1}\stackrel{\vec{a}}{\to}B$ is open hence $\langle a_0,\cdots,a_m\rangle_B$ is open hence $\cR_Y(a_0|a_1,\cdots,a_m)$ the preimage of $\Omega$ in $Y$ is rational open. It follows from Fact 1.2.6 that such maps are spectral.
\end{proof}

\begin{remark}
    Note that $D(T) = \cR(T|)$ is quasicompact in $\Spv\T_{1,\Q_p}$ but not in $\Spa(\T_1,\T_1^\circ)$ as 
    \begin{equation*}
        D(T) = \bigcup_{l=0}^\infty \cR(T|p^l).
    \end{equation*}
    Thererfore $\Spa(\T_1,\T_1^\circ)\to \Spv(\T_1)$ is not spectral.
\end{remark}
%6/26

\begin{proposition}
    \label{proposition_18_3}
    Let $A$ be a complete Huber ring (or sequentially complete, nat, $\id$-adic, (tn;1)). Let $m\in \N$ and $(f_i)_{i=0}^m\in A^{m+1}$ such that 
    \begin{equation*}
        \langle f_0,\cdots,f_m\rangle_A
    \end{equation*}
    is open in $A$. Le t$X=\Cont(A)$. Put $f_i=0$ when $i>m$. Then there is a neighbourhood $U$ of $0$ in $A$ such that for $n\geq m$ and $(g_i)_{i=0}^m\in A^{m+1}$ such that $g_i-f_i\in U$ for $0\leq i\leq n$, we have,
    \begin{enumerate}[1).]
        \item $I=\langle g_0,\cdots,g_n\rangle_A$, 
        \item $\cR_X(f_0|f_1,\cdots,f_m)=\cR_X(g_0|g_1,\cdots,g_n)$.
    \end{enumerate}
\end{proposition}
\begin{proof}
    Let $\vec{f}=(f_i)_{i=0}^m$ and $\vec{g}=(g_i)_{i=0}^m$. By the assumption, $A^{m+1}\stackrel{\vec{f}\cdot}{\to}A$ is open, hence the image $U$ of $(A^\dcirc)^{m+1}$ under $\vec{f}\cdot$ is a neighbourhood of $0$. By our assumption about $\vec{g}$, all $g_i-f_i\in U\subseteq I$, hence all $g_i\in I$. We have $\vec{g} = (I_{m+1}+N)\vec{f}$ where 
    $I_{m+1}$ is an identity $(m+1)$-square matrix and $N$ is an $(m+1)$-matrix with coefficients from $A^\dcirc$. Then $det(I_{m+1}+N)\in 1+A^\dcirc\subseteq A^\times$ as the completeness of $A$ is assumed. Hence $1_{m+1}+N$ is invertible and,
    \begin{equation*}
        \langle f_0,\cdots,f_m\rangle = \langle g_0,\cdots,g_m\rangle.
    \end{equation*}
    For 2)., we distinguish the case $x\not\in X_a$, ($\p=\supp x$ open) and $x\in X_a$, $(\p\in\supp x$ not open).
    In the first case, $A^\dcirc\subseteq\p$ ($a\in A^\dcirc\Rightarrow a^n\in\p$ when $n$ is large $\Rightarrow a\in\p$) hence $U\subseteq \p$ hence $g_i-f_i\in\p$ and $\ab{f_i}_x=\ab{g_i}_x$ and $x$ is in the right hand side of 2). 
    \par In the second case, $M= \max_{0\leq i\leq m}\ab{f_i}_x$ is positive (otherwise $I\subseteq \p$ hence $\p$ is open ad $I$ is). For $a\in A^\dcirc$, $\ab{a}_x<1$ as,
    \begin{equation*}
        \lim_{n\to\infty}\ab{a}_x^n=\lim_{n\to\infty}\ab{a_i^n}_x = \left\vert \lim_{n\to\infty}a^n\right\vert_x = 0
    \end{equation*}
    by the continuity. Let $\rho = \max_{\substack{0\leq i\leq n\\0\leq j\leq m}}\ab{n_{ij}}_x$ where 
    \begin{equation*}
        g_i - f_i = \sum_{j=1}^mn_{ij}f_j.
    \end{equation*}
    Then $\ab{f_i-g_i}_x\leq\rho M$ hence,
    \begin{equation*}
        \max(\ab{f_i}_{x},\rho M) = \max(\ab{g_i}_x,\rho M).
    \end{equation*}
    Thus, $\ab{g_i}_x=M$ when $M=\ab{f_i}_x$. Hence $M$
\end{proof}

\begin{remark}
    Of course from 2), the ? about $\Spa(A,A^+)\subseteq X$ follows. 
\end{remark}

\begin{definition}
    \label{definition_18_5}
    A Tate ring is the sense of Hueber is a nat ring $A$ such that $A^\dcirc\subseteq A^\times\not=\emptyset$ (ie. 3.1 (T)) and such that $A$ contains a bounded open subring. A pair of definition of a Tate ring $A$ is a pair $(B,s)$ where $B$ is a bounded open subring of $A$ and $s\in B^{\dcirc}\cap A^\times$. A ring of definition is (as in the general case of Huber rings) as open bounded subring.
\end{definition}

\begin{proposition}$\:$
    \label{fact_18_11}
    \begin{enumerate}[a).]
        \item If $s\in A^\dcirc\cap A^\times$ and $B\subseteq A$ open bounded then $s^j\in B$ for large enough $j\in\N$. Then $(B,s^{j})$ is a pair of definition. Hence every ring of definition can be completed to a pair of definition.
        \item If $(B,s)$ is a Tate pair of definition for $A$ then $(B,sB)$ is a pair of definition in the sense of Huber rings. Hence any Tate ring is Huber.
        \item If $(B,s)$ is a Tate pair of definition for $A$, $(s^jB)_{j\in\Z}$ or $(s^jB)^{\infty}_{j=n}$ for any $n\in\N$ are neighbourhood bases of $0$, and moreover $X\subseteq A$ is bounded if and only if $X\subseteq s^jB$ for some $j\in\Z$.
    \end{enumerate}
\end{proposition}

\begin{remark}
    \label{remark_18_3}
    Even in the Tate case, $A^\circ$ can fail to be bounded. For example, 
    \begin{equation*}
        A=\Q_p[T]/(T^2), A^\circ = \Z_p+T\Q_p.
    \end{equation*}
\end{remark}

\begin{remark}
    \label{remark_18_4}
    If $A$ is a Tate ring then the image of a bounded or power bounded subset of $A$ under the morphism $A\to B$ of topological rings is bounded or power bounded. Then the universal property of $A\langle X_1,\cdots,X_0\rangle$ (Lemma 3.4.1 or Proposition 3.4.3) will take a 
\end{remark}

\subsection{The structure presheaf on $\Spa(A^{\triangleleft},A^+)$}

We start with the special case of a topological localization of $A^\triangleleft$ with respect to $f\in A^\triangleleft$ such that $A^{\triangleleft}\stackrel{f}{\to} A^\triangleleft$ is open, giving $\Ouv_X(\cR_X(f|))$. 
\begin{proposition}
    \label{proposition_18_1}
    Let $A$ be a complete nat ring with a countable neighbourhood base of $0$ (e.g. a complete Huber ring). Let $f\in A$ such that $A\stackrel{f}{\to}A$ is an open map and $I=\bigcup_{j=0}^\infty \Ker(A\stackrel{f^j}{\to}A)$ and $\overline{I}$ the closure of $I$ in $A$.
    Then $B=A/\overline{I}$ is a complete nat ring, $B\stackrel{f\tmod \overline{I}}{\to}B$ open and $A\to B$ has the universal property for morphisms $A\stackrel{\tau}{\to}T$ such that $T\stackrel{\tau(f)}{\to}T$ is injective and $T$ complete nat ring.
\end{proposition}
\begin{proof}
    It is clear that any such $\tau$ ?. $\ker(\tau)$ is an ideal and $\ker(\tau)f=\ker(\tau)$. Thus, the ? easily follows if $\overline{I}\cdot f=\overline{I}$ (or equivalently, $B\stackrel{f}{\to}B$ injective). For this, we have to consider $a\in A$ such that $fa\in\overline{I}$ and have to show $a\in\overline{I}$. 
    \par For this, let $(U_k)_{k=0}^\infty$ be a neighbourhood base of $0$ in $A$, such that the $U_k$ are open subgroups of $(A,+)$ and 
    \begin{equation*}
        A= U_0\supseteq U_1\supseteq \cdots
    \end{equation*}
    By the assumption $f\cdot U_k$ is open, hence ? $fa\in\overline{I}$. There are $i_k\in I$ such that $fai_k\in fU_K$, where $i_0=0$. Let $(d_k)_{k=1}^\infty$ be such that 
    \begin{equation*}
        d_k\in U_k, fd_k = i_{k-1}-i_k.
    \end{equation*}
    Then $fd_k\in I$ hence $d_k\in I$ as $I\cdot f = I$. Then 
    \begin{equation}
        j = \sum_{k=0}^\infty d_k\in\overline{I}
    \end{equation}
    and $f(a - j) = 0$. Hence $a-j\in I\subseteq\overline{I}$ hence $a\in\overline{I}$. 
\end{proof}
\begin{proposition}
    \label{proposition_18_1_a}
    If $A$ is Huber (Tate, Analytic), then so is $B$ (as these are identified by $A/I$).
\end{proposition}
\begin{lemma}
    If $C\stackrel{\gamma}{\to}C$ is open then $(C/I)\stackrel{r\tmod I}{\to}(C/I)$ is always open.
\end{lemma}
\begin{proof}
    Let $U$ be a neighbourhood of $0$ in $C/I$ then $V = \pi^{-1}U$ is also a neighbourhood of $0$ in $C$ where $C\stackrel{\pi}{\to}C/I$ then $\pi^{-1}(rU)\supseteq rV$ and $rV$ is also a neighbourhood of $0$ in $C$.
\end{proof}
Thus in Proposition \ref{proposition_18_1}, $B\stackrel{f}{\to}B$ is indeed open.
\begin{corollary}
    \label{corollary_18_1}
    Let $C$ be a complete nat ring with a countable neighbourhood base of $0$, $J_0\subseteq C$ a closed ideal and $f\in C$ such that $C\stackrel{f}{\to}C$ is an open map. Let $J = \bigcup_{j=0}^\infty J_0f^j$ and $\overline{J}$ be the colosure of $J$ in $C$, then $B=C/\overline{J}$ is a complete nat ring and $B\stackrel{f\tmod\overline{J}}{\to}B$ an injectiev open map. Moreover $G\to B$ has the universal property for morphismsms $C\stackrel{\tau}{\to}T$ to complete topological rings such that 
    $\tau(J_0)=0$ and $T\stackrel{\tau(f)}{\to}T$ is injective. In addition, if $C$ is Huber (Tate, analytic), then so is $B$.
\end{corollary}
\begin{proof}
    We would like to apply Proposition \ref{proposition_18_1} to $A = C/J_0$. Let 
    \begin{equation*}
        I=\bigcup_{j=0}^\infty(0(f\tmod J_0)^j)_A = \bigcup_{j=0}^\infty J_j/I = J/I
    \end{equation*}
    where $J_j = J\cdot f^j$. ? the ideal contained in Proposition \ref{proposition_18_1} ? ? follow if $\pi(\overline{J}) = \overline{I}$ whereLet $B\stackrel{\pi}{\to}A$ be a projection. We have $\pi(J) = I$. We have $\pi(\overline{J})\subseteq \overline{I}$. On the other hand, 
    \begin{equation*}
        \overline{I}/J_0 = \pi(\overline{J})
    \end{equation*}
    is also closed as $\pi^{-1}(\overline{I}/J_0) = \overline{I}$. Indeed in ? $\pi(\overline{I})$ is closed but $I\subseteq \pi(J)\subseteq\pi(\overline{J})$ hence also $\overline{I}\subseteq\pi(\overline{J})$. Then $\overline{I}=\pi(\overline{J})$ as needed.
\end{proof}
\begin{proposition}
    \label{propositino_18_2}
    Let $A$ be a nat ring and $f\in A$ such that $A\stackrel{f}{\to}A$ is an injective and open map. Consider $A$ as a subring of the localization $A_f = A[T]/\langle 1-T\cdot f\rangle_{A[T]} = A[f^{-1}]$. Equip $(A_f,+)$ with the topology fro which $U$ is a neighbourhood of $0$ in $A_f$ if and only if $U\cap A$ is a neighbourhood of $0$ in $A$. Then we have the following.
    \begin{enumerate}[a).]
        \item $A_f$ is a nat ring and has the universal property for morphisms $A\stackrel{\tau}{\to}T$ of topological rings such that $\tau(f)\in T^\times$.
        \item $A$ is an open subring of $A_f$.
        \item $A$ is Huber (respectively comp)
    \end{enumerate}
\end{proposition}
%6/29
\begin{proof}
The only hard point is the continuity of $A_f\times A_f\to A_f$. Certainley $A\times A\to A\subseteq A_f$ is continuous.  Note that 
$(A_f,+)\stackrel{f\cdot}{\to}(A_f,+)$ is an isomorphism of topological groups. If $U$ is a neighbourhood of $0$, $(f\cdot U)\cap A\supseteq f(U\cap A)$ which is open by out assumption. Thus $fU$ is a neighbourhood of $0$, showing the continuity of the inverse of $A_f\stackrel{f\cdot}{\to}A_f$. Now take $m\in\Z,(f^{-m}A,+)$ is an open subgroup of $(A,+)$ and the continuity of 
\begin{equation*}
    (f^{-m}A)\times(f^{-m}A)\to A_f
\end{equation*}
follows from the diagram belowin which the vertical arrows are homeomorphisms and the bottom horizontal arrow is continuous by our initial observation. As $(A_f,+)$ is the ascending union of its open subgroups $f^{-m}A$, this shows the continuity of $A_f\times A_f\to A_f$.
\begin{center}
    \begin{tikzcd}
(f^{-m}A)\times(f^{-m}A) \arrow[r, "\cdot"] \arrow[d, "(f^m\cdot)\times(f^m\cdot)"'] & A_f \arrow[d, "f^{2m}"] \\
A\times A \arrow[r]                                                                  & A_f                    
\end{tikzcd}
\end{center}
If $A\stackrel{\tau}{\to}T$ is a morphism of topological rings and $\tau(f)\in T^\times$, then we have a ring homomorphism $A_f\stackrel{\tilde{\tau}}{\to}T$. This is continuous. If $U$ is a neighbourhood of $0$ in $T$.
Then $\tilde{\tau}^{-1}U\cap A = \tau^{-1}U$ is a neighbourhood of $0$ in $A$ by continuity of $\tau$, hence $\tilde{\tau}^{-1}U$ is a neighbourhood of $0$ in $A_f$. The rest is also easy.
\end{proof}
\begin{corollary}
    \label{corollary_19_2}
    In the situation of Corollary \ref{corollary_18_1}, $B_{f\tmod\overline{J}}$ is a complete nat-ring and $C\to B_{f\tmod{\overline{J}}}$ has the universal property for morphisms $C\stackrel{\tau}{\to}T$ such that $\tau(J_0)$ and $\tau(f)\in T^\times$. If $C$ is Huber, (Tate, analytic), then so is $B_{f\tmod{\overline{J}}}$.
\end{corollary}

\begin{definition}
    \label{definition_19_1}
    Let $A$ be a nat-ring. We equip $A[X_1,\cdots,X_n]$ with the topology such that $(A[X_1,\cdots,X_n])$ is a topological group and a neighbourhood base of $0$ is given by 
    \begin{equation*}
        \{U[X_1,\cdots,X_n]\sep \text{$U$ a neighbourhood of $0$ in $A$}\},
    \end{equation*}
    where 
    \begin{equation*}
        U[X_1,\cdots,X_n] = \left\{f = \sum_{\alpha\in\N^n}f_\alpha X^\alpha\in A[X_1,\cdots,X_n]\sep \forall \alpha\in\N^n, f_\alpha\in U\right\}.
    \end{equation*}
    We define $A\langle X_1,\cdots,X_n\rangle $ as the completion of $A[X_1,\cdots,X_n]$. If $M$ is an $A$-module, then $M[X_1,\cdots,X_n]$ a module over $A[X_1,\cdots,X_n]$, is the set of formal sum $\sum_{\alpha\in\N^n}m_\alpha X^\alpha$ where $m_\alpha\in M$ , all but finitely many $m_\alpha=0$. 
    The neighbourhood base of $0$ is givenby $U[X_1,\cdots,X_n]$, where $U$ is a neighbourhood of $0$ in $M$, as above. $M\langle X_1,\cdots,X_n\rangle$ is defined as the compltion of $M[X_1,\cdots,X_n]$. 
\end{definition}
\begin{lemma}
    \label{lemma_19_1}
    Let $A$ be a nat-ring, $B$ be a nat-$A$-algebra.
    \begin{enumerate}[a).]
        \item $A[X_1,\cdots,X_n]$ is a nat ring in which $\{X_1,\cdots,X_n\}$ are powerbounded.
        \item If $C\subseteq A$ is an open subring, $C[X_1,\cdots,X_n]$ is an open subring of $A[X_1,\cdots,X_n]$. (Take $U=C$ in Definition \ref{definition_19_1}).
        \item If $A$ is Huber(Tate, analytic), then so is $A[X_1,\cdots,X_n]$. (If $(C,I)$ is a pair of definition for $A$, so is $(C[X_1,\cdots,X_n],I[X_1,\cdots,X_n])$ for $A[X_1,\cdots,X_n]$. If $(s_1,\cdots,s_n)\in A^{\dcirc n}$ such that $\langle s_1,\cdots,s_n\rangle_A =A$ then they do the same in $A[X_1,\cdots,X_n]$).
        \item If $(b_i)_{i=1}^n$ are power bounded in $B$, then the homomorphism,
        \begin{equation*}
            A[X_1,\cdots,X_n]\ni p\mapsto p(b_1,\cdots,b_n)\in B,
        \end{equation*}
        is continuous. 
        \item If $B$ is an adic $A$-algebra (e.g. as $A$ is Tate or analytic), then we have a bijection,
        \begin{equation*}
            \Hom_{(A-\mathrm{alg})^{\mathrm{top}}}(A[X_1,\cdots,X_n],B)\ni\varphi\stackrel{\cong}{\mapsto}(\varphi(X_i))_{i=1}^n(B^\circ)^n.
        \end{equation*}
        (use d) for the surjectivity and well-definedness follows from Fact 3.3.6).
        \item If $\langle f_0,\cdots,f_n\rangle_A$ is open, $I=\langle f_0X_1-f_1,\cdots,f_0X_n-f_n\rangle_{A[X_1,\cdots,X_n]}$ then $A[X_1,\cdots,X_n]/I\stackrel{f_0}{\to}A[X_1,\cdots,X_n]/I$ is an open map.
        \item If $M$ is an adic $A$-module, then so is $M[X_1,\cdots,X_n]$. in particular $A\to A[X_1,\cdots,X_n]$ is adic if $A$ is $\id$-adic.
    \end{enumerate}
\end{lemma}
\begin{proof}
    That $(A,+)$ is indeed a nat-group follows from the fact that the $U$ in Definition \ref{definition_19_1} can be taken to be open subgroups of $(A,+)$. Then $(U[X_1,\cdots,X_n],+)$ is a subgroup of $A[X_1,\cdots,X_n]$.
    \par If $P,Q\in A[X_1,\cdots,X_n]$ and $U$ an open subgroup of $(A,+)$, we have a neighbourhood of $0$, $V$ such that 
    \begin{enumerate}[$\cdot$]
        \item All $P_\alpha V\subseteq U$,
        \item All $Q_\alpha V\subseteq U$,
        \item $V\cdot V\subseteq U$.
    \end{enumerate}
    (Note that only finitely many $\alpha\in\N^n$, $P_\alpha,Q_\alpha\not=0$). Then 
    \begin{equation*}
        (P+V[X_1,\cdots,X_n])(Q+V[X_1,\cdots,X_n])\subseteq (PQ)+U[X_1,\cdots,X_n].
    \end{equation*}
    Thus, multiplication is continuous.
    \par For the powerboundedness assertion, we need to show that 
    \begin{equation*}
        \Theta = \{X^\alpha\sep \alpha\in\N^n\}
    \end{equation*}
    is bounded. For every neighbourhood $U'$ of $0$, we must find a neighbourhood $V'$ of $0$ such that $\Theta V'\subseteq U'$. By Definition \ref{definition_19_1}, it suffices to consider $U'=U[X_1,\cdots,X_n]$ in which case $\Theta U'\subseteq U'$ such that $V'=U'$ will do.
    \par d). Consider $A[X_1,\cdots,X_n]\stackrel{\varphi}{\rightarrow}B$ such that $\varphi(P) = P(b_1,\cdots,b_n)$. This is a ring homomorphism. To show its continuity, consider a neighbourhood $\Omega$ of $0$ in $B$. We must find a neighbourhood $V'$ of $0$ in $A[X_1,\cdots,X_n]$ such that $\varphi(U')\subseteq\Omega$. Without loss of generality, $\Omega$ is a subgroup of $(B,+)$.
    By the powerboundedness of $\{b_1,\cdots,b_n\}$, there is a neighbourhood $V$ of $0$ in $B$ such that $b^\alpha V\subseteq\Omega$. Lt $U'=U[X_1,\cdots,X_n]\subseteq A$ where $U$ is the preimage of $V$ in $A$. If $P\in U[X_1,\cdots,X_n]$, then all $P_\alpha b^\alpha\in b^\alpha V\subseteq\Omega$ hence $\varphi(U')\subseteq\Omega$ as $\Omega$ is a subgroup of $(B,+)$. 
    \par f). Let $U$ be an open subgroup of $(A,+)$ and $V\subseteq A$ the image of $U^{n+1}$ under $A^{n+1}\stackrel{\vec{f}}{\to}A$. if $B=A[X_1,\cdots,X_n]/I$ then the image of $U[X_1,\cdots,X_n]$ in $B$ under $(f_0\tmod{I})$, contains the image of $V[X_1,\cdots,X_n]$ in $B$, which is an open subgroup of $(B,+)$. This will prove f) as the image of $V[X_1,\cdots,X_n]$, $U$ running over open subgroups of $A$, form a neighbourhood base of $0$ in $B$.
    \par For the proof of our claim, let $P = \sum_{\alpha\in\N^n}P_\alpha X^\alpha\in V[T]$. There are $(q_{\alpha,j})_{j=0}^n\in U^{n+1}$ for each $\alpha\in\N^{n+1}$ such that $p_\alpha = \sum_{j=0}^nf_\alpha q_{\alpha,j}$, where we assume $q_{\alpha,j}=0$ when $p_\alpha=0$. Then,
    \begin{align*}
        Q_0 & = \sum_{\alpha\in\N^n}\left(q_{\alpha,j}+\sum_{j=1}^n q_{\alpha-d_j,j}\right)X^\alpha,\\
        Q_j & = \sum_{\alpha\in\N^n} q_{\alpha,j}X^\alpha
    \end{align*}
    where $d_j = (\delta_{ij})_{i=0}^n$, are elements of $U[X_0,\cdots,X_m]$ and $P=f_0 Q_0+\sum_{j=1}^n(f_j-f_0X_j)Q_j$, hence $(P\tmod{I})=f_0(Q_0\tmod{I})$ as needed.
    %7/3
   \par If $\Omega$ is a neighbourhood of $0$ in $A$, then by the assumption, there are $(\omega_i)_{i=1}^n\in\Omega$ such that $M^m\stackrel{\vec{\omega}}{\to}M$ is an open map. We claim that $M[X_1,\cdots,X_n]^m\stackrel{\vec{\omega}}{\to}M[X_1,\cdots,X_n]$ is also open which shows g). But this is obvious. If $U$ is an open subgroup of $M$ then so is the image $V$ of $U^m$ under $\vec{\omega}$, and the $\vec{\omega}$ maps $U[X_1,\cdots,X_n]^n$ onto $V[X_1,\cdots,X_n]$ which is an open subgroup of $M[X_1,\cdots,X_n]$ by Definition \ref{definition_19_1}.
\end{proof}
\begin{proposition}
    \label{fact_20_2}
    Let $X\stackrel{f}{\to}Y$ be a continous homomorphism between abelian nat-groups which have countable neighbourhood bases of $0$. Let $\hat{X}$ and $\hat{Y}$ be their completions. If $f$ is an open map then so is $\hat{X}\stackrel{\hat{f}}{\to}\hat{Y}$.
\end{proposition}
\begin{proof}
    There are open subgroups $(X_i)_{i=1}^\infty$ of $X$ forming a neigh 
    base of $0$, with 
    \begin{equation*}
        X_1\supseteq X_2\supseteq\cdots
    \end{equation*}
    As $f(X_i)$ is open, there is a countable neighbourhood base,
    \begin{equation*}
    Y_1\supseteq Y_2\supseteq\cdots 
    \end{equation*}
    of $Y$ such that $Y_i\subseteq f(X_i)$ is an open subgroup of $Y$. It sufices to show that $\hat{f}(\hat{X}_i)$ is an open subgroup of $\hat{Y}$ and for this, it suffices to show that $\hat{Y}_i\subseteq \hat{f}(\hat{X}_i)$. An element $y$ of $\hat{Y}_i$ is represented as an equivalence class of Cauchy sequence $(y_j)_{j=1}^\infty$ with elements $y_i\in Y_i$. Without loss of generality (taking subsequences), we assume $y_{j+i}-y_j\in Y_{i+j}$ for all $j\geq 1$. By assumption, there is $x_1\in X_1$ such that $f(x_1) = y_1$ and there are $(\xi_j)_{j=1}^\infty$ such that $\xi_j\in X_{i+j}$ such that 
    \begin{equation*}
        f(\xi_j) = y_j+iy_j.
    \end{equation*} 
    Let 
    \begin{equation*}
    x_j = x_1+\sum_{l=1}^{j-1}\xi_l.
    \end{equation*}
    Then 
    \begin{equation}
        x_{j+1}-x_j = \xi_j\in X_{i+j}.
    \end{equation}
    Therefore $(x_j)_{j=1}^\infty$ is Cauchy and $f(x_j)=y_j$. The equivalence class $x$ of $(x_i)_{i=1}^\infty$ is al element $\hat{X}_i$ such that $f(x) = y$. 
\end{proof}
\begin{proposition}
    \label{proposition_20_3}
    Let $A$ be a nat-ring with a countable neighbourhood base of $0$ and $B$ be a complete nat-$A$-algebra.
    \begin{enumerate}[a).]
        \item $A\langle X_1,\cdots,X_n\rangle$ is a nat ring in which $\{X_1,\cdots,X_n\}$ is power bounded. (use $U\langle \cdots\rangle = \widehat{U[\cdots]}$ are $X_i\cdot$-invariant)
        \item If $C$ is an open subring of $A$ then $C\langle X_1,\cdots,X_n\rangle$ is an open subring of $A\langle X_1,\cdots,X_n\rangle$ in which $C[X_1,\cdots,X_n]$ is dense.
        \item If $A$ is Huber (Tate, analytic) then so is $A\langle X_1,\cdots,X_n\rangle$ (For Huber if $C$ in b) is $\langle c_1,\cdots,c_n\rangle_C$-adic, then $C\langle\cdots\rangle$ is $\langle c_1,\cdots,c_n\rangle$-adic.
        \item If $(b_i)_{i=1}6n\in B^{\circ n}$, then there is a unique continuous $A\langle X_1,\cdots,X_n\rangle\to B$ sending $X_i$ to $b_i$ (Lemma \ref{lemma_19_1} + Completeness).
        \item If $B$ is adic overr $A$ (e.g. $A$ is Tate or analytic or ? 3.1(E)), then we have a bijection,
        \begin{equation*}
        \Hom_{\mathrm{nat} A-\mathrm{alg}}(A\langle X_1,\cdots,X_n\rangle,B)\ni\varphi\stackrel{\simeq}{\mapsto}(\varphi(X_i))_{i=1}^n\in(B^\circ)^n.
        \end{equation*}
        (All such $\varphi$ are adic Fact 3.2.21. use Fact 3.3.6 for powerboundedness).
        \item If $\langle f_1,\cdots,f_n\rangle_A$ is open and $I=A\langle X_1,\cdots,X_n\rangle$ is an arbitrary closed ideal containing all $f_0X_i-f_i$ and $B=A\langle X_1,\cdots,X_n\rangle/I$. We have $B\stackrel{f_i}{\to}B$ is open. (Proposition \ref{fact_20_2} applied with $I = \overline{\langle f_0X_i-f_i\rangle}_{A\langle X_1,\cdots,X_n\rangle}$ Then any larger $I$ can be derived from this (Proposition \ref{fact_20_2})).
        \item if $M$ is an adic $A$-module with a countable neighbourhood base of $0$ (e.g. $M=A$ if $A$ is $\id$-adic). Then $M\langle X_1,\cdots,X_n\rangle$ is an adic $A$-module (Proposition \ref{fact_20_2}).
    \end{enumerate}
\end{proposition}
\begin{theorem}
    \label{theorem_20_3}
    Let $A=(A^\triangleleft,A^+)$, $X=\Spa A,\Omega = \cR_X(f_0|f_1,\cdots,f_m)$ a rational open subset, $\langle f_0,\cdots,f_n\rangle_A$ being open in $A$. Let $\mathcal{P}$ be the category of complete Huber pair $(B^\triangleleft,B^+)$.
    \begin{enumerate}[a).]
        \item The functor $F_\Omega$ such that 
        \begin{equation}
            C\mapsto F_\Omega(C) = \{\varphi\in\Hom(A,C)\sep (\Spa(\varphi)(\Spa(C)))\subseteq\Omega\}
        \end{equation}
        is represented by an object $B=(B^\triangleleft,B^+) = (\Ouv_X(\Omega),\Ouv_X^+(\Omega))$. We denote the universal element of $F_\Omega(B)$ by $b\in \Hom(A,B)\subseteq \Hom(A^\triangleleft,\Ouv_X(\Omega))$ and $y = \Spa B$ and $Y\stackrel{\beta}{\to}X$ where $\beta = \Spa(b)$. 
        \item If $\langle f_0,\cdots,f_m\rangle_A=A$ (which is automatic if $A$ is analytic or Tate), then one can take
        \begin{equation*}
            \Ouv_X(\Omega)=B^\triangleleft = A\langle X_1,\cdots,X_m\rangle/I
        \end{equation*}
        when $I\subseteq A\langle X_1,\cdots,X_n\rangle$ is the closure of the ideal generated by $f_0X_i-f_i$ for $1\leq i\leq m$. For $\Ouv_X^+(\Omega)$, one must then take the integral closure of the image of $A^+\langle X_1,\cdots,X_m\rangle$.
        \item In general, one can take,
        \begin{equation}
            \Ouv_X(\Omega) = (A\langle X_1,\cdots,X_m\rangle/J)_{f_0}\label{equation_omega}
        \end{equation}
        where $J\subseteq A\langle X_1,\cdots,X_n\rangle $ is the closure of $\bigcup_{l=0}^\infty I\cdot f_0^l$. For $\Ouv_X^+(\Omega)$, one can take the integral closure of the image of $A^+\langle X_1,\cdots,X_n\rangle$ in Equation \eqref{equation_omega}.
        \item The image of $f_0$ in $\Ouv_X(\Omega)$ is a unit and the image of the ? morphism $A_{f_0}^\triangleleft\to\Ouv_X(\Omega)$ is a dense subring of $\Ouv_X(\Omega)$.
        \item The subring $\Ouv_X^+(\Omega)$ is integral over the closure in $\Ouv_X(\Omega)$ of the subring generated by $b(A)$ and the ${\frac {b(f_i)} {b(f_0)}}$.
        \item The map $Y\stackrel{\beta}{\to}\Omega$ is a homeomorphism sends following $\Rat_Y$ with $(\Rat_Y)_\Omega \defeq\{\Theta\in\Rat_X\sep \Theta\subseteq\Omega\}$.
        \item Let $\Theta\in\Rat_Y$ and $C = (\Ouv_Y(\Theta),\Ouv_Y^+(\Theta)),\Theta_X=\beta(\Theta)$ and $c\in\mathcal{F}_\Theta(C)$ the universal element. The composition $A\stackrel{b}{\to}B\stackrel{c}{\to}C$ is an element of $F_{\Theta_X}$ and has the universal property for that functor. The morphism,
        \begin{equation*}
            (\Ouv_X(\Theta_X),\Ouv_{X^-}^+(\Theta_X))\to(\Ouv_Y(\Theta),\Ouv_Y^+(\Theta))=C
        \end{equation*}
        is an isomorphism.
    \end{enumerate}
\end{theorem}
The proof of the theorem occupies the ret? of this situation Applying c) in the situation of b), $(f_0\tmod\:I)\in (A\langle X_1,\cdots,X_n\rangle/I)^\times$. Since $f_0X_i = f_i$< 
\begin{equation*}
    f_0A\langle X_1,\cdots,X_m\rangle /I = \langle f_0,\cdots,f_m\rangle_{A\langle X_1,\cdots,X_m\rangle/I}=A\langle X_1,\cdots,X_m\rangle/I.
\end{equation*}
Then b) follows from c). Also d) follows from c) and the remainder results in Proposition 1 - 3. c) $\Rightarrow$ a) is clear. c) $\Rightarrow$ e) Follows from the density of $A[X_1,\cdots,X_m]$ in $A\langle X_1,\cdots,X_m\rangle$ (Proposition \ref{proposition_20_3}). c) $\Rightarrow $ f) follows from the density of $A^+[X_1,\cdots,X_m]$ in $A^+\langle X_1,\cdots,X_m\rangle$ again by Proposition \ref{proposition_20_3}.
\par We first show c) (which takes a while). When this is done g) and h) ($\Omega\simeq\Spa B$) remains. Let 
\begin{equation*}
    B^\triangleleft = (A\langle X_1,\cdots,X_m\rangle/J)[f_0^{-1}], 
\end{equation*}
and $B^+$ be the integral closure of the image of $A^+\langle X_1,\cdots,X_m\rangle$ in $B^\triangleleft$ as in Equation \ref{equation_omega}. It follows from the results ? Proposition \ref{proposition_20_3} that $B^\triangleleft$ is complete Huber, and that the following morphisms are open,
\begin{equation}
    A^\triangleleft\langle X_1,\cdots,X_m\rangle\to A^\triangleleft\langle X_1,\cdots,X_m\rangle/J\to (A^\triangleleft\langle X_1,\cdots,X_m\rangle/J)_{f_0}=B^\triangleleft.\label{equation_20_3_a}
\end{equation}
Since $A^+\langle X_1,\cdots,X_m\rangle$ is open in $A^\triangleleft\langle X_1,\cdots,X_m\rangle$, its follows that $B^+$ is an open subring of $B$. Since $A^+\subseteq {A^{\triangleleft}}^\circ$, it can be easily shown that 
\begin{equation*}
    A^+\langle X_1,\cdots,X_n\rangle\subseteq A^\triangleleft\langle X_1,\cdots,X_m\rangle^\circ.
\end{equation*}
By Proposition \ref{fact_15_16}, the images of elements of $A^+\langle X_1,\cdots,X_m\rangle$ is ? ? in Equation \ref{equation_20_3_a} is an open embedding. It easily follows from the definition of "bounded that the image of $A^+\langle X_1,\cdots,X_m\rangle$ in $B^\triangleleft$ is contained in $B^{\triangleleft\circ}$. By Proposition \ref{fact_16_5}, this also holds for the integral closures of ? image. It follows that $B=(B^\triangleleft,B^+)$ is a Huber pair and the composition of morphisms in Equation \ref{equation_20_3_a} with $A^\triangleleft\to A^\triangleleft\langle X_1,\cdots,X_m\rangle$ defines a morphis $A\stackrel{b}{\to}B$.
We now claim that $b\in\mathcal{F}_\Omega(B)$. Indeed, the image $q_i$ of $X_i\in A^\triangleleft\langle X_1,\cdots,X_m\rangle$ in $B$ is in $B^+$ by the definition.
Thus $\ab{q_i}_Y\leq 1$ for all $y\in Y\defeq\Spa B$. Since $b(f_i)= q_ib(f_0)$ (as $f_0X_i-f_i\in I\subseteq J$), we have 
\begin{equation*}
    \ab{b(f_i)}_y\leq \ab{b(f_0)}_y.
\end{equation*}
Set $\beta = \Spa(b)$ then for $x=\beta(y)$,
\begin{equation*}
    \ab{f_i}_x=\ab{b(f_i)}_y\leq \ab{b(f_0)}_y=\ab{f_0}_x.
\end{equation*}
Also 
\begin{equation*}
    \ab{f_0}_x=\ab{b(f_0)}_y\not=0
\end{equation*}
as $b(f_0)\in B^{\triangleleft\times}$. Then $x\in\Omega$ and $b\in \mathcal{F}_\Omega(B)$ as claimed.
\par We want to check that $b$ has the universal property having turning $B$ a representing object for $\mathcal{F}_\Omega$
\begin{equation*}
    \Hom_{\mathrm{op}}(B,C)\stackrel{\varphi\mapsto\varphi b}{\to}\mathcal{F}_\Omega(C)
\end{equation*}
is bijective. For this, let $C$ be a complete Huber pair and $A\stackrel{c}{\to}C$ a morphism of Huber pairs such that $\gamma(\Spa C)\subseteq\Omega$ ? $\Spa(c)$ (i.e. $c\in \mathcal{F}_\Omega(C)$). Let $Z=\Spa C$. For $z\in Z,\vert c(f_0)\vert_Z=\vert f_0\vert_{\gamma(z)}\not=0$ as $\gamma(z)\in \Omega$. Thus, $c(f_0)\in C^{\triangleleft\times}$ by Theorem \ref{theorem_17_2}. Then $q_i = {\frac {c(f_i)} {c(f_0)}}$ exists in $C^\triangle$. For $z\in Z$, $\ab{q_i}_z\leq 1$ as 
\begin{equation*}
    \ab{c(f_i)}_z=\ab{f_i}_{\gamma(z)}\leq \ab{f_0}_{\gamma(z)}=\ab{c(f_0)}_z
\end{equation*}
as $\gamma(z)\in\Omega$. By Theorem \ref{theorem_17_2}, $q_i\in C^+\subseteq C^{\triangleleft\circ}$. By Proposition \ref{proposition_20_3}, there is a unique morphism $A^\triangleleft\langle X_1,\cdots,X_m\rangle\stackrel{c_1}{\to}C^\triangleleft$ of nat-rings where composition with $A^\triangleleft\to A^\triangleleft\langle X_1,\cdots,X_n\rangle$ conincides with $c$ and such that $c_1(X_i) = q_i$.
Since $c(f_i)=q_ic(f_0)$ in $C^\triangleleft$, $c_1$ annhilates $f_0X_i-f_i$. By its continuity, $c_1(I)=0$. As $c_1(f_0) = c(f_0)\in C^{\triangle\times}$, $c_1$ annhilates $I:f_0 ?$ By continuity $c_1(J)=0$. Thus $c_1$ factors over $A^\triangleleft\to A^\triangleleft\langle X_1,\cdots,X_m\rangle/J$ and the first morphism in Equation \eqref{equation_20_3_a}. By the universal property of $(A\langle X_1,\cdots,X_m\rangle/J)_{f_0}$ from Proposition \ref{fact_20_2}, $c_1$ factors uniquely as the second morphism in Equation \eqref{equation_20_3_a} followed by $B^\triangleleft=(A\langle X_1,\cdots,X_m\rangle/J)_{f_0}\to C$.
\begin{equation*}
    c_1(A^+) = c(A^+)\subseteq C^+
\end{equation*}
as $c$ is a morphism of Huber pairs and $q_i\in C^+$ (see the above construction). As $A^+[X_1,\cdots,X_n]$ is dense in $A^+\langle X_1,\cdots,X_m\rangle$, and $C^+$ open, $c_1(A^+\langle X_1,\cdots,X_m\rangle)\subseteq C^+$. 
Thus $c_3(\text{image of $A^+\langle X_1,\cdots,X_m\rangle$ in $B^\triangleleft$})\subseteq C^+$. Since $B^+$ is integral over that image and $C^+$ integrally closed in $C^\triangleleft$, $c_3(B^+)\subseteq C^+$ and we have a morphism $B\stackrel{c_3}{\to}C$ of Huber pairs.
\par By construction $c= c_3b$. That $c$ is uniquely determined by this property follows from,
\begin{equation*}
    c_3\left({\frac {b(a)}{b(f_0)^l}}\right) = {\frac {c(a)} {c(f_0^l)}}
\end{equation*}
for $a\in A^\triangleleft$ and $l\in\N$, together with the density of the image of $A^\triangleleft_f$ in $B^\triangleleft$ observed before. We have these checked the derivd uderived universal property for $A\stackrel{b}{\to}B$.
\par It remains to prove the last two assertions of Theorem \ref{theorem_20_3}. Bijectivity of $Y\stackrel{\beta}{\to}\Omega$, plus identification of topologies of $\Rat$ and $(\Ouv(\cdots),\Ouv^+(\cdots))$. That $\beta(Y)\subseteq\Omega$ follows from $b\in\mathcal{F}_\Omega(B)$ which was checkd before. For the bijectivity of $\beta$, we need the following proposition.
\begin{proposition}
    \label{proposition_21_4}
    Let $A$ be the (completion or) sequential completion of a topological ring $\mathcal{A}$. Then $\Cont(A)\to \Cont(\mathcal{A})$ is bijective.
\end{proposition}
\begin{proof}
    Let $u$ and $v$ be ontinuous aluations on $A$. Let $u_0,v_0$ their compositions with $\mathcal{A}\to A$. Then $\Gamma_{u} = \Gamma_{u_0}$. Indeed, the image of $\mathcal{A}$ is dense in $A$. If $\gamma=u(a)\not=0$, $W=K_{v<r}$ is open in $A$. Hence there is $\alpha\in\mathcal{A}$ whose image in $A$ is in $a+W$. It follows that $v_0(\alpha)=v(a) = \gamma\in\Gamma_{u_0}$. Since these $\gamma$ generates $\Gamma_{u}$, $\Gamma_u\subseteq\Gamma_{u_0}$. The other inclusion is obvious. Thus $\Gamma_u = \Gamma_{u_0}$ and $\Gamma_v = \Gamma_{v_0}$. 
    For injectivity, we assume that $u_0,v_0$ are equivalent. Then without loss of generality, 
    \begin{equation*}
        \Gamma_u = \Gamma_{u_0} = \Gamma_{v_0}=\Gamma_{v}
    \end{equation*}
    and hence $u_0=v_0$ and have to show $u=v$. But this is clear by the continuity of $u$ and $v$.Indeed, if $a\in A$ then by the density of $\mathcal{A}$ in $A$, there is a sequence $(\alpha_i)_{i=1}^\infty$ such that 
    \begin{equation*}
        \lim_{i\to\infty}a_i = a
    \end{equation*}
    where $a_i$ is the image of $\alpha_i\in\mathcal{A}$ in $A$. Since $u$ and $v$ are continuous,
    \begin{equation*}
        u(a) = \lim_{i\to\infty}u(a_i) = \lim_{i\to\infty}u_0(\alpha_i) = \lim_{i\to\infty}v_0(\alpha_i) = \lim_{i\to\infty}v(a_i) = v(a).
    \end{equation*}
    Limits in $\Gamma\cup\{0\}$ exist as it is Hausdorff. 
    \par Thus $\Cont(A)\to\Cont(\mathcal{A})$ is injective. For the surjectiity, let $\mathcal{A}\stackrel{u_0}{\to}\Gamma = \Gamma_{u_0}$ be a continuos valuation which we must extend to $A$. Elements of $A$ are equivalence classes of Cauchy sequences $(\alpha_i)_{i=1}^\infty$. Let $\gamma_i = v_0(\alpha_i)$. If $\lim_{i\to\infty}\gamma_i=0$, then we put $v(a) = 0$. If $(\widetilde{\alpha}_i)$ is equivalent to the given one and $\delta\in\Gamma$, then $K = K_{v_0<\delta}$ is open in $\mathcal{A}$, hence there is $i$ such that $\alpha_i-\widetilde{\alpha}_i\in K$. Then $v_0(\widetilde{\alpha}_i)<\delta$ provided that $v_0(\alpha_i)<\delta$ which happens for sufficiently large $i$. 
    Thus this case is independent of the choise of representatives of teh equivalence class of Cauchy sequences.
    \par if $\lim_{i\to\infty}\gamma_i=0$ does not hold, there is $\delta\in\Gamma$ such that $\gamma_i\geq\delta$ for sufficiently large $i$. Since $K=K_{v_0<\delta}$ is open and $(\alpha_i)$ is a Cauchy sequence, there is $i'$ such that $\alpha_j-\alpha_k\in K$ when $(j,k)\geq i'$. By what was said before, there is $i>i'$ such that $\gamma_i\geq \delta$. If $j\geq i$, then $\alpha_i-\alpha_j\in K$ hence $\gamma_j=u_0(\alpha_j)=u_0(\alpha_i) = \gamma-i$ as $\alpha_i-\alpha_j\in K$. Thus $(\gamma_i)$ stabilizes to some fixed $\gamma\in \Gamma$ as $j\to\infty$, and we put $v(a) = \gamma$. If $\widetilde{\alpha}\sim\alpha$ then $\widetilde{\alpha}_j - \alpha_j\in K$ for $j>i''$ hence $u_0(\widetilde{\alpha}_j)=\gamma$ when $j>\max\{i,i''\}$. hence this case too is independent of choice of $\alpha$. One easily checks the rest.
\end{proof}
% 7/10
We now prove the bijectivity of $Y\stackrel{\beta}{\to}\Omega$. First this makes sense as $\beta(Y)\subseteq\Omega$ which follows from $b\in\mathcal{F}_\Omega(B)$. For the surjectivity, let $r\in\Omega$ be the equivalence class of the valuation $u$ of $A\langle X_1,\cdots,X_n\rangle\stackrel{v}{\to}\Gamma\cup\{\infty\}$ where we assume $\Gamma = \Gamma_u$. We have to extend $u$ to $B^\triangleleft$ which is done in several steps.
We first show the existence and the uniquness of a continuous map,
\begin{equation*}
    A\langle X_1,\cdots,X_n\rangle \stackrel{v}{\to}\Gamma\cup\{0\}
\end{equation*}
such that $v(I)=0$ and the composition,
\begin{equation*}
    A^\triangleleft\to A^\triangleleft\langle X_1,\cdots,X_m\rangle\stackrel{v}{\to}\Gamma\cup\{0\}
\end{equation*}
equals $u$. The uniqueness still holds for $v:\langle X_1,\cdots,X_m\rangle\stackrel{v}{\to}\Gamma'\cup\{0\}$ are allowed where $\Gamma'$ is an ordered group containing $\Gamma$ as an ordered subgroup.
Since $x\in\Omega,\ab{f_0}_x = v(x)\not=0$, hence $x$ extends in a unique way to the valuation below. 
\begin{equation*}
    A^\triangleleft_{f_0}\stackrel{u_1}{\to}\Gamma\cup\{0\}.
\end{equation*}
The uniuqness holds in the same general sense as before, (i.e. values $\Gamma'\supseteq\Gamma$ are also true).
\par We define,
\begin{equation*}
    A[X_1,\cdots,X_n]\stackrel{v_0}{\to}\Gamma\cup\{0\}
\end{equation*}
by 
\begin{equation*}
    (P) = u_1\left(P\left({\frac {f_1} {f_0}},\cdots,{\frac {f_m} {f_0}}\right)\right)
\end{equation*}
where 
\begin{equation*}
    P\mapsto P\left({\frac {f_1} {f_0}},\cdots,{\frac {f_m} {f_0}}\right)
\end{equation*}
is the $A^\triangleleft$-algebra morphism $A^\triangleleft[X_1,\cdots,X_m]\to A^\triangleleft_{f_0}$. For 
\begin{equation*}
    P = \sum_{\alpha\in\N^m}p_\alpha X^\alpha\in A[X_1,\cdots,X_m]
\end{equation*}
we have, 
\begin{align*}
    v_0(P) &= u_11\left(\sum_{\alpha\in\N^m}p_\alpha\prod_{i=1}^m\left({\frac {f_i}{f_0}}\right)^{\alpha_i}\right),\\
    &\leq \max_{\alpha\in\N^m}u_1\left(p_\alpha\prod_{i=1}^m\left({\frac {f_i} {f_0}}\right)^{\alpha_i}\right),\\
    & \leq \max_{\alpha\in\N^m}u(p_\alpha)\prod_{i=1}^m\left({\frac {v(f_i)} {v(f_0)}}\right)^{\alpha_i},\\
    &\leq \max_{\alpha\in\N^m}v(p_\alpha) \tag{b}
\end{align*}
It follows that $v_0$ is continuous for the topology on $A^\triangleleft[X_1,\cdots,X_n]$ considered in Lemma \ref{lemma_19_1} (as $v_(K_{v< r})[X_1,\cdots,X_m]\subseteq [0,r)_{\Gamma_\cup\{0\}}$ by (b) 
and $K_{v<r}[X_1,\cdots,X_m]$ is open in $A^\triangleleft[X_1,\cdots,X_m]$).
\par By Proposition \ref{proposition_21_4}, there is a unique continuous extension,
\begin{equation*}
    A^\triangleleft\langle X_1,\cdots,X_m\rangle\stackrel{r}{\to}\Gamma,
\end{equation*}
satisfying our conditions.
We have 
\begin{equation*}
    v(f_0X_i-f_i)=v_0(f_0X_i-f_i) = v_0 = v_1\left(f_0{\frac {f_i} {f_0}}-f_i\right) =1.
\end{equation*}
Hence $r$ is annhilates the ideal generated by the $f_0X_i-f_i$. 
By the continuity, it also annhilates its closure $I$. Conversely, let $A^\triangleleft\langle X_1,\cdots,X_m\rangle\stackrel{v'}{\to}\Gamma'$ with $\Gamma\subseteq\Gamma'$ (ordered preserving inclusion of a subgroup) such that 
\begin{equation*}
    A^\triangleleft\to A^\triangleleft\langle X_1,\cdots,X_m\rangle\stackrel{v'}{\to}\Gamma'\cup\{0\}
\end{equation*}
equals to $v$. Let $P\in A\polring{m}\backslash\{0\}$ be of degree at most $d$, 
\begin{equation*}
    Q(X_0,\cdots,X_m)=X_0^dP\left({\frac {X_1} {X_0}},\cdots{\frac {X_m} {X_0}}\right)
\end{equation*}
be the homogeneousation of $P$, and let $q = Q(f_0,\cdots,f_m)\in A$. One easily checks that 
\begin{equation}
    f_0^dP-q\in\langle f_0X_1-X_1,\cdots, f_0X_m-f_m\rangle_{A^\triangleleft[X_1,\cdots,X_m]}\tag{c}
\end{equation}
Since $v$ annhilates $I$, $v'(f_0X_i-f_i)=0$ and it follows that 
\begin{equation*}
    u(f_0)^dv'(P) = v'(f_0^dP) = v'(q) = u(q).
\end{equation*}
hence,
\begin{equation*}
    v'(P) = {\frac {u(q)} {v(f_0^d)}} = u_1\left({\frac {q} {f_0^d}}\right) = u_1\left(P\left({\frac {f_1} {f_0}},\cdots,{\frac {f_m} {f_0}}\right)\right) = v(P).
\end{equation*}
By the uniquness part of Proposition \ref{proposition_21_4}, $v'=v$. 
\par Since $v(f_0) = u(f_0) \not=0$, $J$ is also annhilates by $v$. This is then a unique continuous valuation
\begin{equation*}
    A^\triangleleft\compring{m}/J\stackrel{v_1}{\to}\Gamma\cup\{0\}
\end{equation*}
when compsosed with $A^\triangle\compring{m}\to A^\triangleleft\compring{m}/J$ equals $v$. As $v_1(f_0) = v(f_0)\not=0$, this is a unique continuous,
\begin{equation*}
    B^\triangleleft\stackrel{w}{\to}\Gamma\cup\{0\}
\end{equation*}
where composition with $A\compring{m}/J\to B^\triangleleft$ the second morphism in $(a)$ equals $v_1$. Since this morphism is an open inclusion and $v_1$ is continuous, one easily check the continuity of $w$. 
\par We claim that the equivalence class of $w$ is an element of $Y=\Spa(B)$. By (b), 
\begin{equation*}
    v(f)= v_0(f)\leq 1,
\end{equation*}
when $f\in A^+\polring{m}$. Since this ring is dense in $A\compring{m}$, by the continuity of $v$,
\begin{equation*}
    v(f)\leq 1
\end{equation*}
when $f$  belongs to the latter ring.
\par Thus, the image in $B^\triangleleft$ of $A^+\compring{m}$ is contained in $V_{w\leq 1}$. Since $K_{w\leq 1}$ is the integrally closed subring of $B^\triangleleft$ of the image of $A^+\compring{m}$. We obtain,
\begin{equation*}
    w(B^+)\subseteq [0,1]_{\Gamma\cup\{0\}}
\end{equation*}
as claimed.
\par Thus, the equivalence class $y$ of $w$ is an element of $Y$ such that 
\begin{equation*}
    \beta(y). =x.
\end{equation*}
If $y'$ has the same property, then as $A^\triangleleft\to B^{\triangleleft,\ab{\cdot}_{y'}}\to \Gamma' = \Gamma_{y'}$ equivalent to $v$ and $\Gamma_x = \Gamma_u = \Gamma$, we may assume without loss of generality that $\Gamma\subseteq\Gamma'$ (with order preserving inclusion of a subgroup) and that the composition,
\begin{equation*}
    A^\triangleleft\to B^\triangleleft\stackrel{\ab{\cdot}_{y'}}{\to} \Gamma'
\end{equation*}
actually equals to $v$. By the uniqueness assertion about $v$, the composition $A^\triangleleft\compring{m}\to B^\triangleleft\stackrel{\ab{\cdot}_{y'}}\to \Gamma'$ equals to $v$. 
Since $A^\triangleleft\compring{m}\to A^\triangleleft\compring{m}/J$ is injective, the composition $A^\triangleleft\compring{m}/J\to B^\triangleleft\to\stackrel{\ab{\cdot}_{y'}}\Gamma'$ and $A^\triangleleft\compring{m}/J\stackrel{v_1}{\to}\Gamma\subseteq\Gamma'$ also coincide.
Finally $w=\ab{\cdot}_{y'}$ as it is uniquely determined by $v_1$. Thus $y=y'$ and the bijectivity of $\beta$ is shown.
\par For the comparison of topological and of rational open subsets,it suffices to consider, 
\begin{equation*}
    \Theta = \cR_Y(g_0|g_1,\cdots,g_j),
\end{equation*}
where $g_i\in B^\triangleleft$ and $\langle g_1,\cdots,g_j\rangle_{B^\triangleleft}$ is open in $B^\triangleleft$, 
and  show that $\beta(w)\in\Rat_X$. This will identify the topology base as claimed when ($\Theta_X = \Rat_X(f_0|f_1,\cdots,f_m)$), one can take $g_i = (\text{image of $f_i$})$ in $B^\triangleleft$ which generate open ideal in $B^\triangleleft$ as $A^\triangleleft\to B^\triangleleft$ is adic.
\par By Proposition 3.3.3, this is an open subgroup $U\subseteq B^\triangleleft$ such that 
\begin{equation*}
    \cR_Y(g_0|g_1,\cdots,g_j) = \cR_Y(r_0|r_1,\cdots,r_j,\cdots,r_m)
\end{equation*}
and 
\begin{equation*}
    \langle g_0,\cdots,g_j\rangle_{B^\triangleleft} = \langle r_0,\cdots,r_m\rangle_{B^\triangleleft}
\end{equation*}
when $g_i-r_i\in U$ for $0\leq i\leq j$ and $r_iin U$ when $j<i\leq m$.
As $A^\triangleleft_{f_0}$ has a dense image in $B^\triangleleft$, we can take $r_i,0\leq i\leq j$, for this subring,
\begin{equation*}
    r_i = \left(\text{image in $B^\triangleleft$ of ${\frac {\varphi_i} {f_0^{a_i}}}\in A_{f_0}$}\right)
\end{equation*}
Since $f_0$ is a unit in $B^\triangleleft$, we can multiply all $r_i$ and $g_i$ by $f_0^a$ with out changin $\Theta$ or $\langle g_0,\cdots,g_j\rangle$ or $\langle r_0,\cdots,r_j\rangle$. 
THus it can be assumed that $r_i$ ins the image of $\varphi_i\in A^\triangleleft$ in $B^\triangleleft$. Then $\beta$ maps $\Theta$ onto $\Omega\cap\cR_X(\varphi_0,\cdots,\varphi_j)$ but $\langle \varphi_0,\cdots,\varphi_j\rangle_{A^\triangleleft}$. Let $V\subseteq A^\triangleleft$ be the preimage under $A^\triangleleft\to B^\triangleleft$ of the open subgropu $f_0^aU$. Then we still have $\beta(\Theta) = \Omega\cap\cR_X(\varphi_0,\cdots,\varphi_0)$ if $\varphi_i\in V$ for $j<i\leq m$. As $A$ is Huber it has (oifin) and (tn:0) and there a finite subset $S\subset A^{\triangleleft\dcirc}$ generating an open ideal.
By Fact 3.1.5, $S$ is tn. Assume $S^d\subseteq V$, when $V$ is sufficinetly large. As $A^\triangleleft$ has (0;fq), $S^d$ still generates an open ideal. Selecting $\varphi_0,\cdots,\varphi_m$ such that they form a ? of $S^d$, we may assume that 
\begin{equation*}
    \langle \varphi_0,\cdots,\varphi_m\rangle
\end{equation*}
is open in $A^\triangleleft$. Then $\cR_X(\varphi_0|\varphi_1,\cdots,\varphi_m)\in\Rat_X$.
\par Finally,
\begin{equation*}
    \beta(\Theta) = \Omega\cap\cR_X(\varphi_0|\varphi_1,\cdots,\varphi_m) = \cR_X(f_0\varphi_0|\text{ all other }f_i\varphi_j)
\end{equation*}
The topological claim then follows. If $\Theta_X = \beta(\Theta)$, one checks that 
\begin{align*}
    \mathcal{F}_\Theta(C)&\rightarrow \mathcal{F}_{\Theta_X}(C),\\
    B\stackrel{v}{\to}C\Rightarrow A\stackrel{\beta}{\to}B\stackrel{r}{\to}C
\end{align*}
is bijective usig the definition and the bijectivity of $\beta$. Thus the composition of structure presheaves also holds.
\begin{corollary}
    \label{corollary_22_1}
    \begin{equation*}
        \Spa(A)\cong \Spa(\hat{A})
    \end{equation*}
\end{corollary}
% \par Recall $I\subseteq A^\triangleleft\langle X_1,\cdots,X_m\rangle$ is an ideal generated by the closure of 
% \begin{equation*}
%     \{f_0X_i-f_i\}_{i=1}^n.
% \end{equation*}
% And we set,
% \begin{equation*}
%     J = \bigcup_{i\in\Z_{j\geq0}} If_0^j.
% \end{equation*}

%7/13

\section{Acyclicity proofs}

The Acyclicity proofs usually consider a special ? of Čech cohomology to show that $\Ouv_X$ extends from $\Rat_X$ to $\Ouv_X$. Acyclicity will typically follow using ? the follows.
\begin{proposition}[Grothendieck Tohoku Paper Lemma 3.3.1]
    Let $\mathscr{A}\stackrel{F}{\to}\mathscr{B}$ be a left exact functor between abelian categories where $\mathscr{A},\mathscr{B}$ contain sufficineltly many injective objects. Let $\mathscr{X}$ be a class of objects of $\mathscr{A}$
    such that 
    \begin{enumerate}[a).]
        \item Every direct summand of (in particular, every object isomorphic to) an element of $\mathscr{X}$ is in $\mathscr{X}$.
        \item Every objet $A$ of $\mathscr{A}$ has a monomorphism $A\to X$ with $X\in\mathscr{X}$.
        \item If 
        \begin{equation*}
            0\to X'\to X\to X''\to 0,
        \end{equation*}
         is a short exact sequence in $\mathscr{A}$ such that $X',X\in\mathscr{X}$, then $X'\in\mathscr{X}$ and 
        \begin{equation*}
            0\to FX'\to FX\to FX''\to 0.
        \end{equation*}
    \end{enumerate}
\end{proposition}

\begin{corollary}
    All injective objects of $\sA$ belongs to $\sX$ and furthermore $X\in\sX$ and $p>0$, we have $R^pFX = 0$.
\end{corollary}

\subsection{Acyclicity Criteria for General Sheaves on $\Spa A$}

In this subsection, we denote $X=\Spa A$ where $A$ is a Huber pair.

\begin{proposition}
    \label{propositin_23_1}
    Let $x\in X$ be such that $\cGamma_x = \Gamma_x$ and $U$ be an open subset of $X$ such that $x\in U$. Then there is $m\in\N$ and $(f_i)_{i=0}^m\in(A^\triangleleft)^{m+1}$ such that 
    \begin{equation*}
        \langle f_0,\cdots,f_m\rangle_{A^\triangleleft} = A^\triangleleft
    \end{equation*}
    and 
    \begin{equation*}
        x\in\cR_X(f_0|f_1,\cdots,f_m)\subseteq U.
    \end{equation*}
    In particular, Fact 3.2.11, this holds when $x$ is closed.
\end{proposition}

\begin{proof}
    Since rational open subsets form a topology base, there are $(\varphi_i)_{i=0}^m$ such that 
    \begin{equation*}
        x\in\cR_X(\varphi_0|\varphi_1,\cdots,\varphi_m)\subseteq U.
    \end{equation*}
    If $a\in A$ is such that $\ab{a}_x\not=0$, then this implies
    \begin{equation*}
        x\in\cR(a\varphi_0|\varphi_1,\cdots,\varphi_m)\subseteq\cR(\varphi_0|\varphi_1,\cdots,\varphi_m)\subseteq U.
    \end{equation*} 
    As $\ab{\varphi_0}_x\in\Gamma_x = \cGamma_x$, such $a$ can be chosen such that 
    \begin{equation*}
        \ab{\varphi_0a}_x\geq 1.
    \end{equation*}
    Then 
    \begin{equation*}
        x\in\cR(a\varphi_0|a\varphi_1,\cdots,a\varphi_m,1)\subseteq\cR(a\varphi_0|a\varphi_1,\cdots,a\varphi_m)\subseteq U.
    \end{equation*}
    Note that 
    \begin{equation*}
        \langle a\varphi_0,\cdots,a\varphi_m,1\rangle_{A^\triangleleft} = A^\triangleleft.
    \end{equation*}
\end{proof}

\begin{remark}$\:$
    \begin{enumerate}
        \item It also holds in the Tate case and the analytic case for arbitrary $x$.
        \item Thus $\mathcal{O}_X(\cR(f_0|f_1,\cdots,f_m))$ can be bounded by the ? Then 3b) (where ? $f_0$ is not necessary).
    \end{enumerate}
\end{remark}

\begin{definition}$\:$
    \begin{enumerate}[a).]
        \item A rational covering of $X$ is a covering $X = \bigcup_{i=1}^m\cR_X(f_i|f_1,\cdots,f_m)$ such that 
        \begin{equation*}
            \langle f_1,\cdots,f_m\rangle_{A^\triangleleft} = A^\triangleleft.
        \end{equation*}
        \item The rational sieve defined by such $(f_i)_{i=1}^m$ is 
        \begin{equation*}
            \{U\in\Ouv_X\sep U\subseteq\cR(f_i|f_1,\cdots,f_m)\text{ for some $i\in\{1,\cdots,m\}$}\}.
        \end{equation*}
        i.e. the sieve generated by this covering.
    \end{enumerate}
\end{definition}

\begin{proposition}
    \label{proposition_23_2}
    $\:$
    \begin{enumerate}
        \item Every open covering $\U$ of $X$ has a refinement which is a rational covering.
        \item Eery coering sieve (i.e. $X=\bigcup_{U\in\sS}U$) of $X$ contains a rational sieve.
    \end{enumerate}
\end{proposition}

\begin{proof}
    Let $\sS$ be a covering sieve of $X$. For $x\in X$, we have $U_x\in\S$ such that $x\in U_x$ and $(f_{x,i})_{i=1,\cdots,m_x}$ with 
    \begin{equation*}
        \langle f_{x,1},\cdots,f_{x,m_x}\rangle_{A^\triangleleft} = A^\triangleleft
    \end{equation*}
    such that $x\in\cR(f_{x,0}|f_{x,1},\cdots,f_{x,m_x})\subseteq U_x$. Then 
    \begin{equation*}
        X = \bigcup_{x\in X_{\mathrm{cl}}}\cR_X(f_{x,0}|f_{x,1},\cdots,f_{x,m_x}).
    \end{equation*} 
    Indeed, the right hand side $V$ contains all closed points. If $V\subsetneqq X$, then $X\backslash V$ would be a non-empty closed subset without closed points, a contradiction by Fact 1.2.7.
    \par As $X$ is spectral, it is quasi-compact, hence there are finitely many closed points $(x_i)_{i=1}^n$ such that 
    \begin{equation*}
        X = \bigcup_{j=1}^n\cR(f_{x_j,0}|f_{x_j,1},\cdots,f_{x_j,m_j}).
    \end{equation*}
    Let $f_{j_i} = f_{x_j,i},n_j = m_{x_j},U_j = \cR(f_{j,0}|f_{j,1},\cdots,f_{j,m_j})$. Let $J$ be the set of fnctions $\iota$ defined on $\{1,\cdots,n\}$ such that 
    \begin{equation*}
        0\leq \iota(j)\leq m_j,
    \end{equation*}
    and let $J_0\subseteq J$ be the set of all $\iota\in J$ such that there is $j\in\{1,\cdots,n\}$ such that $\iota(j)=0$. For all $\iota\in J$, let 
    \begin{equation*}
        f_\iota = \prod_{j=1}^nf_{j,\iota(1)}
    \end{equation*}
    then 
    \begin{equation*}
        \langle f_\iota\sep \iota\in J\rangle_{A^\triangleleft} = \prod_{j=1}^n\langle f_{j,0},\cdots,f_{j,m}\rangle_{A^\triangleleft}=A^\triangleleft
    \end{equation*} 
    by the assumption. Note that 
    \begin{enumerate}[a).]
        \item $X = \bigcup_{i=1}6n\cR(f_{j,0}|f_{j,1},\cdots,f_{j,m_j}) = \bigcup_{j=1}^m U_j \in\sS$.
        \item $\langle f_{j,0},\cdots,f_{j,m_j}\rangle_{A^\triangleleft}=A^\triangleleft$.
        \item $\max_{\iota\in J}\ab{f_\iota}_x = \max_{\iota\in J_0}\ab{f_\iota(x)}$ when $x\in X$. Note that when $x\in U_k$ and $\iota(j) = \begin{cases}
            \iota(j),j\not=k,\\
            0,j=k.
        \end{cases}$ then $\iota'\in J_0$ and $\ab{f_\iota}_x\leq \ab{f_{\iota'}}_x$.
        \item $\langle f_\iota\sep \iota\in J_0\rangle_{A^\triangleleft}=A^\triangleleft$. (Follows from b), c) and Theorem \ref{theorem_17_2}).
    \end{enumerate}
    Then the $(f_\iota)_{\iota\in J_0}$ defines a rational sieve/covering. We claim that this rational sieve is contained in $\sS$. This will finish the proof. 
    \par To prove this claim, it suffices to consider,
    \begin{equation*}
        \Omega_\lambda = \cR(f_\lambda|\text{ all $f_\iota$ with $\iota\in J_0$}).
    \end{equation*}
    with $\lambda\in J_0$ and to show that $\Omega_\lambda\in\sS$, it suffices to find $j\in\{1,\cdots,n\}$ such that 
    \begin{equation*}
        \Omega_\lambda\subseteq U_j.
    \end{equation*}
    We claim that this holds when $\lambda(j) = 0$ (by Definition such $j$ exists for $\lambda\in J_0$). Indeed, let $x\in\Omega_\lambda$ then 
    \begin{equation*}
        \ab{f_\lambda}_x = \prod_{k=1}^n\ab{f_{k,\lambda(k)}}\not=0.
    \end{equation*}
    Hence 
    \begin{equation*}
        0\not=\ab{f_{j,\lambda(j)}}_x = \ab{f_{j,0}}_x.
    \end{equation*}
    If $x\in U_j$, there would be $i\in \{1,\cdots,m_j\}$ such that 
    \begin{equation*}
        \ab{f_{j,i}}_x>\ab{f_{j,0}}.
    \end{equation*}
    Let 
    \begin{equation*}
        \lambda'(k) = \begin{cases}
            \lambda(k),\quad k\not=j,\\
            i,\quad k=j.
        \end{cases}
    \end{equation*}
    Then 
    \begin{equation*}
        \ab{f_{\lambda'}(x)} = \ab{f_{j,i}(x)}\prod_{k\not=j}\ab{f_{k,\lambda(k)}}_x>\ab{f_{j,0}}_x\prod_{k\not=j}\ab{f_{k,\lambda(k)}}_x = \ab{f_\lambda(x)}
    \end{equation*}
    For the inequality in the middle use that this is not $0$ as they are factors of $\ab{f_\lambda}_x\not=0$.
    By c), there is $\theta\in J_0$ such that $\ab{f_\theta(x)}\geq \ab{f_{\lambda'}}_x>\ab{f_\lambda(x)}$, contradicting $x\in\cR(f_\lambda|\text{ all $f_\iota$ for $\iota\in J_0$})\subseteq\cR(f_\lambda|f_\theta)$.
\end{proof}

%7/17

\begin{proposition}
    \label{proposition_24_3}
    Let $A$ be a Huber pair and $X=\Spa A$ and $\mathcal{G}$ be a presheaf of abelian groups on $\Rat_X$ such that for every $\Omega\in\Rat_X$ and every rational covering $\mathcal{U}$ of $\Omega$, the map,
    \begin{equation}
        \mathcal{G}(\Omega)\to \hat{H}^0(\mathcal{U},\mathcal{G})
    \end{equation}
    is bijective. Then 
    \begin{equation}
        \mathcal{G}\to(\underline{\Sheaf}(\mathcal{G}))_{\Rat_X}
    \end{equation}
    is an isomorphism in the category of presheaves on $\Omega$.
\end{proposition}

\begin{proposition}
    Let $X=\Spa A$ where $A$ is a Huber pair. The class $\mathscr{X}$ of sheaves $\mathcal{F}$ of abelian groups on $X$ such that 
    \begin{equation}
        \hat{H}^p(\mathcal{U},\mathcal{F}) = 0
    \end{equation}
    where $p>0$ and $\mathcal{U}$ is a rational covering of $\Omega\in\Rat_X$, together with the functor
    \begin{equation}
        \mathcal{F}\mapsto\mathcal{F}(\Omega)
    \end{equation}
    where $\Omega\in\Rat_X$ satisfy the assumption of Proposition 2.4.3. Then all injective sheaves are in $\mathscr{X}$ and $H^p(\Omega,\mathcal{F})=0$ when $\Omega\in\Rat_X,\mathcal{F}\in\mathscr{X}$ and $p>0$.
\end{proposition}

We have somewhat more convenient acyclicity criterion in the Tate case.

\begin{definition}
    \label{definition_24_2}
    Let $X=\Spa A$ where $(A^\triangleleft,A^+)$ is a Huber pair. We inductively define the property of a sieve $\sS$ in $\Omega\in\Rat_X$ to have a Laurent series $\leq o$ as follows.
    \begin{enumerate}
        \item A sieve has Laurent order $0$ if it is the all sieve.
        \item If $o>0$ and Laurent order $\leq o-1$ has already been defined then $\sS$ has Laurent order $\leq o$ if there is $g\in \Ouv_X(\Omega)$ such that $\sS|_{\cR_\Omega(g|1)}$ and $\sS|_{\cR_\Omega(1|g)}$ both have Laurent order $\leq o-1$. 
    \end{enumerate}
    The Laurent order of $\sS$ is the smallest $o\in N$ such that $\sS$ has Laurent order $\leq. o$ or $\infty$ if no such $o$ exists. We denote $o_L(\sS)$ to be the Laurent order of $\sS$.
\end{definition}

\begin{proposition}
    \label{fact_24_1}$\:$
    \begin{enumerate}[a).]
        \item $\sS\subseteq\mathcal{T}\Rightarrow o_L(\mathcal{T})\leq o_L(\sS)$. 
        \item $\Theta\subseteq\Omega\Rightarrow o_{L,\Theta}(\sS|_\Theta)\leq o_{L,\Omega}(\sS)$ 
        \item $o_L(\sS)<\infty\Rightarrow \sS$ is a covering sieve.
    \end{enumerate}
\end{proposition}

The following lemma does not yet need the Tate assumption.

\begin{lemma}
    \label{lemma_24_1}
    Let $G=\{g_i\sep 1\leq i\leq n\}$ be a non-empty finite subset of $A^{\triangleleft\times}$. Then the rational sieve defined by $G$ has finite Laurent order.
\end{lemma}
\begin{proof}
    We write $g_i\preceq g_j$ if $\Omega=\cR(g_j|g_i)$. Obviously, this is a partial order with allowed equivalence. We call $(i,j)$ in order if at least one of $g_i\preceq g_j$ or $g_j\preceq g_i$ holds. 
    We prove the claim by induction on the number of pairs not in order. 
    If all pairs are in order then without loss of generality by renumbering $G$, we can assume $g_1\preceq \cdots\preceq g_m$. Thus, all $g_i\leq g_m$. and the rational sieve contains $\Omega=\cR_\Omega(g_m|g_0,\cdots,g_{m-1})$ hence is the all sieve thus has Laurent order $0$.
    \par Let $(i,j)$ be not in order and the assertion shown for a smaller number of pairs not in order. Let $q ={\frac {g_i} {g_j}}$ and $\Theta_1 = \cR_\Omega(q|1)$ and $\Theta_2=\cR_\Omega(1|q)$. Then the pairs in order stay in order after restricting to $\Theta_i$ and $g_j|_{\Theta_1}\preceq g_i|_{\Theta_1}$ and $g_i|_{\Theta_2}\preceq g_j|_{\Theta2}$.
    Thus $(i,j)$ becomes in order for $G_k = \{g_i|_{\Theta_k}\sep 1\leq i\leq n\}$ where $k=1,2$.By the induction assumption and Definition \ref{definition_24_2}, the rational sieve defined by $G$ has finite Laurent order.
\end{proof}

\begin{lemma}
    \label{lemma_24_2}
    Let $A$ be a Huber pair such that $A^\triangleleft$ is Tate and there is $(f_i)_{i=0}^m\in A^{m+1}$ such that $\langle f_0,\cdots,f_m\rangle_{A^\triangleleft}=A^\triangleleft$. Then there is $\varepsilon\in A^{\triangleleft\times}\cap A^{\triangle\dcirc}$ such that 
    \begin{equation}
        \ab{\varepsilon}_x \leq\max_{0\leq i\leq n}\ab{f_i}_x.\label{eq:lemma_24_2_3}
    \end{equation}
\end{lemma}

\begin{proof}
    Let $s\in A^{\triangleleft\dcirc}\cap A^{\triangleleft\times}$. By assumptino, there are $(\lambda_i)_{i=0,\cdots,m}\in (A^\triangleleft)^{m+1}$ such that 
    \begin{equation*}
        \sum_{i=0}^m \lambda_i f_i = 1.
    \end{equation*}
    Since $A^+$ is an open subring and $\lim_{j\to\infty} s^j\lambda_i = 0$, there is $1\leq j\leq m-1$ such that all $\lambda_is^j\in A^+$. Then one can define $\varepsilon = s^j$. 
    \begin{equation*}
        \ab{\varepsilon}_x =\left\vert\sum_{i=0}^m (s^j\lambda_i)f_i\right\vert_x\leq\max_{0\leq i\leq m}\ab{\lambda_is^j}_x\ab{f_i}_x\leq \max_{0\leq i\leq m}\ab{f_i}_x.
    \end{equation*}
    since $\ab{\lambda_is^j}\leq 1$ since $x\in\Spa A$ and $\lambda_i s^j\in A^+$.
\end{proof}

\begin{proposition}
    \label{proposition_24_5}
    In the situation of Definition \ref{definition_24_2}, assume $A$ (i.e $A^\triangleleft$) to be Tate. Then a sieve $\sS$ in $\Omega\in\Rat_X$ covers $\Omega$ if and only if $\omega_L(\sS)<\infty$. 
\end{proposition}
\begin{proof}
    The proof will work by showing that every rational sieve has finite $o_L$ which proves the claim by Proposition \ref{proposition_23_2} and \ref{fact_24_1} a).
    \par As was said before it is sufficent to show that every rational sieve has finite Laurent order. Let $\sS$ be the rational sieve defined by the finite subset of $G$ of $\Ouv_X(\Omega)$. We use induction on the number of elements of $G$ which are not in $\mathcal{O}_X^\times(\Omega)$.
    If this number is $0$, Lemma \ref{lemma_24_1} can be applied. Let $g_1\not\in\mathcal{O}_X^\times(\Omega)$ and the assertion shown for a smaller number of non-elements.
    Since $A$ is Tate, there is $\delta\in\mathcal{O}_X^\times(\Omega)$ such that $\ab{\delta}_\omega<\max_{1\leq i\leq m}\ab{g_i}_\omega$ holds for all $\omega\in\Omega$. Consider 
    \begin{equation*}
        \Theta_1 = \cR_\Omega(g_1|\delta)=\cR\left({\frac {g_1} \delta}\bigg{\vert}1\right),\Theta_2 = \cR_\Omega(\delta|g_1) = \cR\left(1\bigg{\vert}{\frac {g_1} \delta}\right).
    \end{equation*}
    On $\Theta_1,g_1$ becomes a unit as $\ab{g_1}_\theta\geq\ab{\delta}_\theta\not=0$ when $\theta\in\Theta_1$ hence the induction assumption applies to $G|_{\Theta_1}$. But the induction assumption also applies to $G|_\Theta$ since
    \begin{equation*}
        \cR_{\Theta_2}(g_1|g_2,\cdots,g_m)=\emptyset 
    \end{equation*}
    as for $\vartheta\in\Theta_2$, $\ab{g_1}_\vartheta\leq\ab{\delta}_\vartheta<\max_{1\leq i\leq m}\ab{g_i}_\vartheta$. And,
    \begin{equation*}
        \cR_{\Theta_2}(g_i|g_1,g_2,\cdots,g_m) = \cR(g_i|g_2,\cdots,g_m)
    \end{equation*}
    as for $\vartheta\in\Theta_2$, 
    \begin{equation}
        \ab{g_1}_\vartheta\leq\ab{\delta}_\vartheta<\max_{2\leq i\leq m}\ab{g_i}_\vartheta.\label{eq:proposition_24_5_star}
    \end{equation}
    \par Hence $g_1|_{\Theta_2}$ can simply be removed from $G_{\Theta_2}$ without affection the rational sieve by \eqref{eq:proposition_24_5_star} it follows that $g_2|_{\Theta_2},\cdots,g_m|_{\Theta_2}$ have no common zero in $\Theta_2$, hence define a rational sieve.
\end{proof}

\begin{proposition}
    \label{proposition_24_6}
    Let $X=\Spa A$ where $A$ is Tate, $\mathcal{G}$ be a presheaf on $\Rat_X$ (of abelian groups) such that for all $\Omega\in\Rat_X$ and all $g\in\mathcal{O}_X(\Omega)$, the sequence,
    \begin{equation}
        \begin{tikzcd}
0 \arrow[r] & \mathcal{G}(\Omega) \arrow[r, "{r\mapsto(r|_{\Omega_1},r|_{\Omega_2})}"] & \mathcal{G}(\Omega_1)\oplus\mathcal{G}(\Omega_2) \arrow[r, "{(r_1,r_2)\mapsto(r_1|_{\Omega_{12}}-r_2|_{\Omega_{12}})}"] & \mathcal{G}(\Omega_{12}) \arrow[r] & 0
\end{tikzcd}
    \end{equation}
    is exact where $\Omega_1 = \cR_\Omega(g|1),\Omega_2 = \cR_\Omega(1|g)$ and $\Omega_{12} = \Omega_1\cap \Omega_2$. Then $\mathcal{G}\to(\underline{\Sheaf}(\mathcal{G}))|_{\Rat_X}$ is an isomorphism in the category of pressheaves of abelian groups on $X$.
\end{proposition}

\begin{proposition}
    \label{proposition_24_7}
    Let $X=\Spa A$ where the Huber pair $A$ is Tate. Then the assumption in Proposition 2.4.3 hold for $\mathscr{A}=(\text{sheaves of abeliangroups on $X$})$ and of the functor $\mathcal{F}\to \mathcal{F}(\Omega)$ and $\mathscr{X}$ the class of sheaves $\mathcal{F}$ on $X$ such that 
    \begin{equation}
        \begin{tikzcd}
0 \arrow[r] & \mathcal{F}(\Omega) \arrow[r, "{f\mapsto(f|_{\Omega_1},f|_{\Omega_2})}"] & \mathcal{F}(\Omega_1)\oplus\mathcal{F}(\Omega_2) \arrow[r, "{(f_1,f_2)\mapsto(f_1|_{\Omega_12}-f_2|_{\Omega_{12}})}"] & \mathcal{F}(\Omega_{12}) \arrow[r] & 0
\end{tikzcd}
    \end{equation}
    is exact for all $\Omega\in\Rat_X$ and all $g\in\mathcal{O}_X(\Omega)$ where $\Omega_1,\Omega_2$ and $\Omega_{12}$ are as in Proposition \ref{proposition_24_6}.
\end{proposition}

\begin{thebibliography}{9}
\bibitem{Markus}
Huber, R \emph{Continuous valuations}, Mathematische Zeitschrift
\end{thebibliography}

\end{document}